\documentclass[11pt,a4paper]{article}
\usepackage[utf8]{inputenc}
\usepackage[T1]{fontenc}
\usepackage{geometry}
\geometry{margin=1in}
\usepackage{fancyvrb}
\usepackage{parskip}
\usepackage{hyperref}
\hypersetup{colorlinks=true,linkcolor=black,urlcolor=blue}
\usepackage{framed}

\title{Maya2C Technical Reference}
\author{Generated from the workspace's module documentation}
\date{\today}

\begin{document}
\maketitle

\begin{framed}
\noindent This document is generated. It is a \emph{technical reference}
assembled from the module documentation in the Maya2C source tree, not a
whitepaper: it describes what the code says about itself, and it does not argue
a design.

\medskip

\noindent Sections describing modules that are not reachable by consensus carry
a banner saying so. Those banners are taken from the modules' own status
documentation, not from an editorial judgement made here.
\end{framed}

\tableofcontents
\newpage
\section{The node}
\noindent\texttt{crates/node\textbackslash{}src}

\subsection{config.rs}
\texttt{config.toml}: the operational settings a node reads from disk.

\subsubsection{What is deliberately absent, and why it is absent from the *type*}

There is no consensus field here. Not "there is one and you should not set
it" — there is no field, so no configuration file can express one.

[\texttt{crate::consensus::ChainConfig}] carries three values and every one of them
is consensus:

\begin{verbatim}
| Field | What a per-host override does |
|---|---|
| `verify_pow` | `false` makes a node that accepts **any** block |
| `pow_limit` | a different difficulty floor is a different network |
| `dag.activation_height` | decides which proof-of-work rule applies when |
\end{verbatim}

A node whose file disagrees with its peers forks — silently, and usually
only under load. Chain rules come from \texttt{genesis.json}, which is one shared
artefact every operator diffs against the ceremony's \texttt{COMMITMENT.txt}. That
is the whole reason it is a distributed file rather than a local one, and
putting any of it here would turn a typo into a chain split.

The same reasoning \texttt{crates/governance/src/limits.rs} uses: the safest way to stop a
value being set is for there to be no way to name it.

\subsubsection{Precedence}

\begin{verbatim}
command-line flag  >  config.toml  >  compiled default
\end{verbatim}

A flag beats the file so an operator can override one setting for one run
without editing state that outlives the run — which is what people actually
do at three in the morning.

\subsubsection{Unknown keys are an error}

\texttt{deny\_unknown\_fields} on every struct. A typo like \texttt{p2p\_prot = 30333} would
otherwise be silently ignored and the node would listen on the default,
which is the failure mode where the config *looks* applied and is not.

\subsection{consensus/chain.rs}
Block index, fork choice, and chain reorganization.

\subsubsection{Fork choice}

The active chain is the one with the greatest \textbf{cumulative work}, never the
greatest height. A long branch of easy blocks represents less expended
hashpower than a short branch of hard ones, and height cannot distinguish
them.

Ties are broken by keeping the incumbent tip. Switching on equal work would
let a peer flip the active chain back and forth at no cost.

\subsubsection{Reorganization}

Switching branches is: walk both tips back to their common ancestor, revert
the abandoned blocks newest-first through the undo journal, then apply the
new branch oldest-first.

If applying the new branch fails partway — an invalid transaction in a block
a peer sent — the chain must not be left straddling two branches. Everything
applied so far is rolled back and the original branch is restored before the
error is returned, so a failed reorg is a no-op rather than corruption.

\subsection{consensus/chain/store.rs}
The persistent half of [\texttt{Chain}]: opening a stored chain, importing headers
ahead of their bodies, adopting a snapshot, and pruning.

\subsection{consensus/difficulty.rs}
Difficulty retargeting and cumulative chain work.

\subsubsection{Retargeting}

Every [\texttt{RETARGET\_INTERVAL}] blocks the target is rescaled by how far the
window's actual elapsed time diverged from the intended
\texttt{RETARGET\_INTERVAL × TARGET\_BLOCK\_TIME}:

\begin{verbatim}
new_target = old_target × actual_timespan ÷ expected_timespan
\end{verbatim}

A larger target is easier, so a window that ran slow (large actual timespan)
raises the target and speeds the chain back up.

The ratio is clamped to \texttt{[1/4, 4]} per retarget. That bound is what stops a
single window — whether from an honest hashrate cliff or a manipulated
timestamp — from moving difficulty arbitrarily far in one step.

\subsubsection{Work}

Comparing branches by height is wrong: a long chain of easy blocks must lose
to a shorter chain of hard ones. The comparable quantity is *expected hashes
spent*, which for a target \texttt{t} is \texttt{2²⁵⁶ ÷ (t + 1)}.

\subsection{consensus/miner.rs}
Multi-threaded proof-of-work search, under whichever rule the target height
calls for.

Two rules exist and both are permanent: ArgonBlake below
[\texttt{DAG\_ACTIVATION\_HEIGHT}] and the memory-hard DAG at or above it. The caller
picks with [\texttt{PowMode}], because the caller is the one that knows what height
the candidate is for.

[\texttt{DAG\_ACTIVATION\_HEIGHT}]: crate::crypto::dag::DAG\_ACTIVATION\_HEIGHT

\subsubsection{Nonce partitioning}

Threads stride through the nonce space rather than taking contiguous blocks:
thread \texttt{i} of \texttt{n} tries \texttt{i, i+n, i+2n, …}. Contiguous ranges would make the
first thread do all the work whenever a solution is found early, since low
nonces are tried first.

\subsubsection{Memory}

Under [\texttt{PowMode::Argon}] each concurrent hash allocates a 32 MiB Argon2id
buffer, so thread count is a memory decision as much as a CPU one: 8 threads
is a quarter-gigabyte of working set. [\texttt{suggested\_threads}] accounts for
that.

The DAG modes allocate nothing per thread. Every worker reads the one shared
cache or dataset, so the footprint is 64 MiB or 4 GiB total rather than per
worker — which is the whole point of a shared dataset, and why a GPU can run
thousands of lanes against it.

\subsection{consensus/mod.rs}
Consensus: difficulty retargeting, cumulative work, fork choice, and mining.

\subsection{consensus/uint.rs}
Minimal unsigned 256-bit integer.

Difficulty targets are 256-bit big-endian thresholds, and both retargeting
(\texttt{target × actual ÷ expected}) and cumulative chain work (\texttt{2²⁵⁶ ÷ (target+1)})
need arithmetic wider than \texttt{u128}. This is the smallest type that does the
job; it is not a general-purpose bignum.

Limbs are stored least-significant first. Every operation is explicit about
overflow — nothing here silently wraps.

\subsection{core/batch.rs}
\begin{framed}\noindent\textbf{Status.} This module is a research branch. It is present in the tree and no consensus path reaches it. Nothing described here is part of the running chain.\end{framed}

Batch identifiers, and the shard access set of a transaction.

\subsubsection{Status}

\textbf{Research branch. No consensus path reaches this module.} It is the
node-side half of [\texttt{maya\_blockgraph}], which carries the bounds and the
scheduler and the reasons they live in a crate of their own. The network
worker that would disseminate batches and the header field that would
reference them are not written. See \texttt{docs/blockgraph.md}.

\subsubsection{Why the derivation is here and not there}

\texttt{maya-blockgraph} has no dependencies, so Kani can compile it. That rules
out BLAKE3 and therefore rules out computing a digest. The same boundary
\texttt{maya-dex} and \texttt{maya-lattice-pow} live behind.

\subsection{core/codec.rs}
Bounds-checked reader for decoding untrusted wire bytes.

Everything here operates on data received from the network, so no path may
panic: a malformed frame from a hostile peer must surface as an error, never
as an index-out-of-bounds or a huge allocation.

\subsection{core/dex\_payload.rs}
Wire forms of the trading transactions.

Split out of [\texttt{crate::core::payload}] because that module was already at the
size where finding anything meant scrolling. The tags themselves stay there:
they are one numbering, and one numbering belongs in one place.

Every structure here is fixed-width apart from a route's leg list. Fixed
width means the decoder cannot be told to allocate, and it means a
transaction's size is a function of what it does rather than of who sent it.

\subsubsection{Two paths, priced differently}

[\texttt{SwapRequest}] does \textbf{not} execute where it appears in the block. It joins
that pair's batch and settles at the batch's single clearing price once
every transaction has been staged — see [\texttt{maya\_dex::batch}] for why.

[\texttt{SwapRoute}] does execute in place, leg after leg, because its legs depend
on each other and there is nothing to net a multi-hop path against. The
asymmetry is deliberate. A route is the arbitrage primitive: whoever sends
one is competing on speed and accepts that ordering matters. A plain swap is
the ordinary user's, and they are the ones being protected.

It does not reopen the sandwich, because routes run *before* the block's
batches. A miner can put a trade in front of a victim's batch, but there is
no place left in the same block to unwind it, and a front-run that cannot be
unwound is just a position.

\subsection{core/governance\_payload.rs}
Wire forms of the governance transactions.

All six are fixed-width apart from a proposal's change list, so a decoder
cannot be told to allocate, and a transaction's size is a function of what
it does rather than of who sent it.

\subsubsection{Why a work claim carries a beneficiary rather than a signature}

Nothing in [\texttt{WorkClaim}] proves the sender mined anything, and it does not
need to. The claim is only valid inside a block, at most one per block, and
putting it there required the proof of work the block already carries. The
miner is the party that chose the block's contents, so the miner is the
party entitled to say who is credited.

\subsubsection{Why a ballot carries no weight}

The voter's weight is read from chain state at the height the vote lands —
their locked stake plus their decayed work credit. A ballot that carried its
own weight would be a ballot whose weight the voter chose.

\subsection{core/htlc\_payload.rs}
The wire forms of the lattice HTLC transitions.

Three actions, and the split between them is who bears a failure. A lock is
the sender's own transaction, so a bad one — no balance, an expiry already
past — is an error. A claim and a refund are not: they race each other at
the expiry boundary as a matter of course, and the loser of that race must
not void the block the winner is in. See \texttt{crate::state::htlc\_exec}.

Sizes: a lock carries a 4,448-byte commitment and a claim a 1,408-byte
opening. Both are smaller than the 11,165-byte hybrid signature on the same
transaction, so HTLC-L does not change what bounds a block.

\subsection{core/identity\_payload.rs}
The wire forms of the identity transitions.

\subsubsection{Rotation is authorised by the key it replaces}

That is the whole security argument, and it needs no new cryptography. A
transaction is already signed by a hybrid pair and the sender's address is
BLAKE3 over both halves — so a rotation transaction *sent from the subject's
own address* is, by construction, signed by the key the document currently
names. The new key is data inside that transaction.

ML-DSA-65 is a lattice signature. The brief asked for "post-quantum lattice
proofs" for rotation, and this is one: the proof that the old key authorised
the new is the old key's signature over a payload naming it. A second proof
system alongside would be another thing to get wrong for nothing.

\subsubsection{Why the new key rides in the payload rather than being derived}

The subject signs with the old key and names the new one. The alternative —
deriving the next key from the current one — would make the whole chain of
keys computable from any single one of them, which is the opposite of what
rotation is for.

\subsubsection{What none of these carry}

A claim. \texttt{AnchorAttestation} carries a Merkle root; \texttt{SetRevocationBit}
carries an index. There is no field in any of them a personal fact could go
in, and that is checked by \texttt{tests/identity\_tests.rs} scanning committed
state rather than by reading the types.

\subsection{core/iot\_payload.rs}
The wire forms of the IoT anchor transactions.

Every layout and bound belongs to \texttt{maya-iot-anchor}; each message has one
fixed size, so this file only slices that many bytes off the reader.

\subsection{core/mod.rs}
Core chain data structures.

\subsection{core/oracle\_payload.rs}
Wire forms of the oracle transactions.

\subsubsection{Why a feed update is large, and why that is not fixable here}

A quorum's worth of post-quantum signatures is what it costs to know that a
quorum signed. ML-DSA and SLH-DSA do not aggregate — there is no post-
quantum analogue of a BLS multisignature that would let \texttt{n} signers produce
one constant-size proof — so a submission carries one public key and one
signature per authority, and grows linearly in the quorum.

Measured, at 13,157 bytes per signer — a 1,984-byte hybrid public key beside
an 11,165-byte hybrid signature:

\begin{verbatim}
| Quorum | Submission | Per 8 MiB block |
|---|---|---|
| 3 | 38.6 KiB | 212 |
| 5 | 64.3 KiB | 127 |
| 11 | 141.4 KiB | 57 |
| 21 | 269.9 KiB | 30 |
\end{verbatim}

That is a fact about post-quantum cryptography, not a tuning knob, and it is
the reason [\texttt{MAX\_OBSERVATIONS}] exists. Reproduce it with
\texttt{cargo bench --bench oracle}.

The keys have to be carried, incidentally, for the same reason a transaction
carries its sender's: an address is a *hash* of a hybrid key, so a registry
entry cannot yield the key needed to check a signature.

\subsubsection{The beacon proof is small}

Eighty bytes plus a height. It is a VRF proof, not a signature quorum, so it
costs nothing to carry and one verification to check.

\subsection{core/payload.rs}
Typed transaction payloads.

The base chain has no scripting, so state-channel operations reach the state
transition as typed payloads carried by an ordinary [\texttt{Transaction}].

\subsubsection{Wire compatibility}

A [\texttt{TxKind::Transfer}] carries \textbf{no payload section at all} — not an empty
one. That keeps a plain transfer's signing bytes byte-identical to the
pre-payload format, so existing signatures, transaction ids, and stored
blocks all remain valid. Adding even a zero-length marker would have
invalidated every signature ever produced.

Non-transfer kinds append \texttt{tag || fields} after the nonce. Because the
preceding fields are fixed-width and a payload'd transaction is strictly
longer, no payload encoding can collide with a transfer's signing bytes.

\subsubsection{Batch authorization}

A [\texttt{ChannelClosure}] carries \textbf{both} parties' signatures over the channel
state. The enclosing transaction's own signature pays for inclusion and
nothing more: whoever submits a batch cannot move a single unit of value
that both channel participants did not already authorize off-chain. That
separation is what makes cross-channel batching safe — a batch settling a
hundred channels has two hundred authorizers, and one outer signature could
never speak for them.

[\texttt{Transaction}]: crate::core::Transaction

\subsection{core/rwa\_payload.rs}
The wire forms of the RWA transitions.

Five actions, and the split between them is about who bears the cost of a
failure. Issuance, distribution and attestation are the issuer's own
transactions, so a bad one is an error and fails only that transaction's
block-worth of intent. A \textbf{DvP is not}: it names a counterparty, and an
error there would let anyone void a block by submitting a swap they know
cannot settle. See \texttt{crate::state::rwa\_exec}.

\subsection{core/sealed\_payload.rs}
Wire forms of the two sealed-mempool transactions.

\subsubsection{Why a sealed transaction is a payload and not a transaction}

The obvious design encrypts a whole [\texttt{crate::core::Transaction}] and has the
chain decrypt and execute it. That needs the inner transaction to carry its
own signature, its own nonce, and its own fee — a second, nested copy of
every rule the outer one already enforces, with a second chance to get each
of them wrong.

So a [\texttt{SealedEnvelope}] is an ordinary transaction whose *payload* is
encrypted. The sender signs it, pays for it, and burns a nonce for it, all
through the machinery that already exists. What the committee reveals is a
[\texttt{crate::core::TxKind}] — the action — which then executes under the
envelope's authority. There is exactly one signature, one nonce, and one
sender, and the encryption covers only the part that needs to be secret.

\subsubsection{Why the share is a separate transaction}

A committee member publishing a decryption share is doing something on the
public record: it is attributable, it is checkable, and a member who does it
wrong should be identifiable afterwards. Gossiping shares off-chain would be
cheaper and would leave nothing behind.

\subsection{core/suite\_tx.rs}
Wire version 7: suite-tagged transactions — ADR-007.

\begin{verbatim}
7 ‖ inputs ‖ outputs ‖ suite:u8 ‖ pk_len:u32le ‖ pk ‖ nonce:u64le
  ‖ has_sig:u8 [‖ sig_len:u32le ‖ sig] ‖ has_payload:u8 [‖ payload]
\end{verbatim}

Both lengths must equal the registry's for the suite — checked before the
bytes are read, so a hostile length costs nothing — and the signing bytes
commit to the suite id as well as the key, so a signature cannot be
re-presented under another suite that happens to accept the same bytes.

Dark until \texttt{SUITE\_ENVELOPE\_ACTIVATION\_HEIGHT}: see
[\texttt{Transaction::verify\_at}].

\subsection{core/threat\_payload.rs}
The wire form of an attack attestation.

The layout and every bound belong to \texttt{maya-threat-intel}; this file only
slices the transaction reader so that crate sees exactly one attestation.
The header is read first and the declared length checked against the
evidence cap before a byte of evidence is read.

\subsection{core/transaction.rs}
Transactions and their hybrid ML-DSA-65 + SLH-DSA-SHA2-128s authorization.

\subsubsection{The address is no longer the key}

Under ed25519 a 32-byte public key *was* the address: outputs named a key,
state was keyed by that same value, and a verifier recovered the key it
needed straight from the address it was checking. A hybrid public key is
1984 bytes and cannot play that role.

So an address is \texttt{blake3} over *both* keys (see
[\texttt{crate::crypto::hybrid::address\_of}]) and stays 32 bytes, while the keys
themselves travel in the transaction. Two consequences follow, and both are
load-bearing:

- [\texttt{Transaction::public\_key}] is not redundant with the sender's address. It
  is the only place the keys appear, and [\texttt{Transaction::signing\_bytes}]
  commits to both so neither signature can be re-presented under another
  key — including beside a different partner key from the other scheme.
- The sender's address is *derived*, never taken from the wire. A transaction
  that could name its own sender independently of the keys that signed it
  would let anyone spend from any account.

\subsubsection{Both signatures, always}

[\texttt{Transaction::verify}] passes only if the lattice proof *and* the
hash-based proof check out. There is no single-signature frame, no optional
second half, and no wire version that accepts one: a transaction with one
valid signature is exactly as rejected as a transaction with none.
[\texttt{crate::crypto::hybrid}] explains why that is worth 11 kilobytes.

\subsubsection{Suite-tagged transactions (wire version 7)}

ADR-007's envelope: one signature under the suite its id names, carried in
[\texttt{Transaction::suite\_auth}], encoded by [\texttt{crate::core::suite\_tx}]. Dark:
[\texttt{Transaction::verify}] refuses every one, and only
[\texttt{Transaction::verify\_at}] at or past \texttt{SUITE\_ENVELOPE\_ACTIVATION\_HEIGHT}
(\texttt{u64::MAX}) can accept one. The hybrid fields are unused on such a frame.

\subsection{crypto/argon\_blake.rs}
ArgonBlake: a memory-hard hybrid proof-of-work hash.

Three stages:

1. \textbf{BLAKE3 pre-hash} of the raw header bytes, compressing arbitrary-length
   input to a fixed 32-byte block that becomes the Argon2 "password".
2. \textbf{Argon2id} over 32 MiB, which supplies the memory hardness — the reason
   this construction resists ASIC and GPU acceleration.
3. \textbf{BLAKE3 squeeze} (XOF) over the Argon2 output, producing the final
   digest.

\subsubsection{Performance, and why there is no SIMD tuning here}

Measured on this codebase with \texttt{cargo bench --bench argon\_blake}:

\begin{verbatim}
| Stage | Time |
|---|---|
| Full ArgonBlake hash | ~25.4 ms |
| BLAKE3 pre-hash alone | ~120 ns |
\end{verbatim}

The BLAKE3 stages are \textbf{0.0005\%} of the total. Argon2id over 32 MiB is
essentially the entire cost, by design — that memory-hardness is the point
of the construction.

Three facts follow, and together they close off SIMD tuning as an option:

- \texttt{blake3} is already SIMD-accelerated with runtime CPU dispatch. The crate
  exposes only \texttt{no\_avx2} / \texttt{no\_avx512} / \texttt{no\_sse41} *opt-outs*; there is
  nothing to enable.
- \texttt{argon2 0.5} ships no SIMD feature at all. It is a portable scalar
  implementation.
- Building with \texttt{-C target-feature=+avx2} was measured at \textasciitilde{}27.0 ms against a
  \textasciitilde{}25.4 ms baseline: no improvement, well inside run-to-run noise.

So no target-feature flags are set. Beyond being useless here, \texttt{target-cpu}
tuning in the release image would risk \texttt{SIGILL} on any deployment host older
than the build host.

The only remaining lever would be replacing the Argon2 implementation, which
would change the digest. That is a hard fork: see
\texttt{argon\_blake\_matches\_the\_frozen\_known\_answer} in \texttt{tests/crypto\_tests.rs}.

\subsubsection{Why the salt is derived, not random}

Proof-of-work must be *verifiable*: any node holding only the header bytes
has to recompute an identical digest. A random salt would make that
impossible, so the salt is derived deterministically from the header under
its own domain-separation context. Distinct headers still get distinct
salts, so Argon2's memory cost cannot be amortized across candidates.

\subsection{crypto/dag/cache.rs}
The verification cache: what every node holds, and what the dataset is
derived from.

\subsubsection{Why a cache exists at all}

A miner needs the full 4 GiB dataset resident, because it reads pages of it
at random and recomputing one on demand would cost far more than the memory
it saves. A *validator* reads only [\texttt{ACCESSES}](super::ACCESSES) pages per block, so it can
afford to recompute each one — and recomputing needs nothing but this
\textasciitilde{}64 MiB cache.

That asymmetry is the whole reason the design is usable. Without it, running
a node would mean holding 4 GiB resident to check other people's work, and
the set of machines that can validate the chain would collapse to the set
that can mine it.

\subsubsection{Construction}

Sequentially memory-hard, following Ethash's use of Lerner's RandMemoHash:

\begin{verbatim}
cache[0] = H(seed)
cache[i] = H(cache[i-1])                       for i in 1..n

repeat CACHE_ROUNDS times:
    for i in 0..n:
        v        = cache[i][0] mod n
        cache[i] = H(cache[(i-1+n) mod n] ⊕ cache[v])
\end{verbatim}

The first pass is a plain chain — cheap, and enough to fill the array with
material that depends on the seed. The rounds are what make it hard: each
rewrite reads a *data-dependent* index, so the whole array has to be present
for any of it to be reproduced. An implementation that tried to recompute
cache items on demand would find each one depends on an unpredictable
earlier one, transitively on most of the array.

Three rounds, as Ethash uses. The cost is 4 × n hashes — around 4.2 million
for a 64 MiB cache. Measured on this codebase: **0.9 s in release, 4.4 s in
a dev build**. Paid once per [\texttt{EPOCH\_LENGTH}] blocks, which is 5.21 days, so
even the debug figure is four orders of magnitude inside its budget.

[\texttt{EPOCH\_LENGTH}]: crate::crypto::dag::EPOCH\_LENGTH

\subsection{crypto/dag/dataset.rs}
The mining dataset: 4 GiB derived from the 64 MiB cache.

\subsubsection{The asymmetry this file exists to create}

Every dataset item is a pure function of the cache and its own index
([\texttt{dataset\_item}]). A miner materialises all 67 million of them once per
epoch and then reads pages at random; a validator materialises none of them
and recomputes the handful a block actually touched. Same function, two
footprints — 4 GiB against 64 MiB.

That is what makes the work bandwidth-bound rather than arithmetic-bound.
Recomputing an item costs 256 cache reads and two BLAKE3 compressions, so an
implementation that tried to avoid holding the dataset pays roughly a
hundredfold in compute for what it saves in memory. Holding it is the only
sane strategy, and holding it means buying DRAM on the same open market
everyone else buys it on.

\subsubsection{Cost}

\texttt{2 × dataset\_items()} BLAKE3 compressions plus \texttt{256 × dataset\_items()} FNV
mixes — about 134 million compressions at [\texttt{Params::MAINNET}]. Generation is
spread across threads because the items are mutually independent given the
cache: item *i* never reads item *j*. Measured single-threaded this is
minutes; across eight threads it is tens of seconds, paid once per 5.21-day
epoch and paid ahead of the boundary rather than on it.

\subsubsection{Why the buffer is zeroed first}

[\texttt{Vec::resize}] writes 4 GiB of zeros before the real fill overwrites them.
That costs a fraction of a second against a generation measured in tens, and
it buys a \texttt{\&mut [Item]} that can be split across threads with no \texttt{unsafe} at
all. Handing out uninitialised memory to eight threads to save 0.3 s in a
consensus-critical allocation is not a trade worth making.

\subsection{crypto/dag/hashimoto.rs}
Hashimoto: the hash function the chain's proof of work actually is.

One digest is [\texttt{ACCESSES}] pseudo-random 128-byte reads through the epoch's
dataset, folded into a 128-byte mix and squeezed to 32 bytes. Each lookup
index depends on the mix so far, so the reads are strictly sequential in
their dependency and cannot be issued ahead of time. Arithmetic between
reads is a handful of multiply-xors — deliberately trivial, so that the DRAM
round trip is the entire cost and the hash rate a machine achieves is a
direct reading of its memory bandwidth.

\subsubsection{Two implementations of one function}

[\texttt{hashimoto\_full}] reads a materialised 4 GiB [\texttt{Dataset}]. [\texttt{hashimoto\_light}]
recomputes each page it touches from the 64 MiB [\texttt{Cache}]. They must return
identical bytes — a miner uses the first, the node that validates the miner's
block uses the second — and \texttt{light\_and\_full\_agree\_on\_every\_nonce} asserts it
rather than trusting it.

Light costs about a hundred times more per hash: 128 items recomputed at 256
cache reads each. That ratio is the design's load-bearing number. Too small
and mining from the cache becomes competitive, which deletes the memory
requirement the whole scheme rests on; too large and validation stops being
cheap. Around 100× is where Ethash sits and where this sits.

\subsubsection{Cost, measured against what it replaces}

A light hash is roughly 33,000 BLAKE3 compressions, on the order of 1–3 ms.
ArgonBlake, the pre-fork rule, is \textasciitilde{}25.4 ms. Validation gets an order of
magnitude cheaper at the same time as mining gets bound to bandwidth.

\subsection{crypto/dag/mod.rs}
A memory-hard, DAG-based proof of work.

Structurally this is Ethash: a small \textbf{cache} that every node holds, and a
large \textbf{dataset} derived from it that only miners hold. A hash reads
pseudo-random pages of the dataset, so the work is bound by memory bandwidth
rather than by arithmetic. Verification recomputes the handful of pages it
needs from the cache, which is what keeps a validator's footprint at tens of
megabytes instead of gigabytes.

\subsubsection{What this is for, stated honestly}

The goal is \textbf{not} to make general-purpose hardware beat purpose-built
silicon. Nothing has ever achieved that, and the design this one is modelled
on is the proof: Ethash was the canonical ASIC-resistant DAG, and Ethash
ASICs shipped in April 2018 at roughly 2–2.5× the efficiency of contemporary
GPUs. The mechanism is structural. When an algorithm is bound by commodity
DRAM bandwidth, an ASIC buys the *same* commodity DRAM and then deletes the
display engine, the shader array, the PCIe complex and the driver stack.

What a DAG does buy is a \textbf{ceiling on that advantage}. A SHA-256 ASIC beats
a CPU by roughly four orders of magnitude because the work is pure
arithmetic and arithmetic is what custom silicon is best at. Here the work is
DRAM traffic, which custom silicon has to buy on the same open market as
everyone else. The realistic outcome is a 2–5× gap rather than a 10,000×
one — which is the difference between GPU mining being worthless and GPU
mining being viable. That is the claim this module makes, and it is the only
one it can support.

\subsubsection{Design target: GDDR6/6X/7, not HBM}

Deliberately tuned for 400–1000 GB/s commodity graphics memory. Tuning for
HBM3 would favour datacenter accelerators — the most centralized and least
commodity hardware there is — which is the opposite of the point.

\subsubsection{Fixed size, rotating contents}

Ethash grows its dataset every epoch, which is what eventually made 4 GB
cards unable to mine Ethereum at all. Here [\texttt{DATASET\_BYTES}] is constant and
only the *contents* rotate, once per [\texttt{EPOCH\_LENGTH}] blocks. A card that
can mine today can mine indefinitely, and the hardware floor is a published
constant rather than a moving deadline.

\subsubsection{Hashing primitive}

BLAKE3 throughout, where Ethash uses Keccak. It is already the chain's hash
(\texttt{src/core/block.rs:76}, \texttt{src/crypto/argon\_blake.rs:84}), so this adds no
dependency and no second primitive to audit.

Domain separation uses derived keys rather than \texttt{Hasher::new\_derive\_key} in
the hot loop: \texttt{new\_derive\_key} re-hashes the context *string* on every call,
and dataset generation makes on the order of 10⁸ of them. The context is
hashed once into a 32-byte key (\texttt{cache\_key}, \texttt{item\_key}, \texttt{mix\_key})
and keyed hashing is used thereafter, which is one compression per call.

\subsection{crypto/dag/params.rs}
Sizing parameters for the DAG, and the prime search that fixes them.

\subsubsection{Why the sizes are a value and not only a constant}

Consensus runs at [\texttt{Params::MAINNET}]: a 64 MiB cache and a 4 GiB dataset.
Nothing else may run there, and no code path here can select a different
size for a block on the real chain — the chain's parameters are pinned in
[\texttt{crate::consensus::chain::ChainConfig}] and default to \texttt{MAINNET}.

Tests are the reason the sizes are a value at all. A determinism test wants
to build *whole* datasets and compare them, twice, from two epochs, and a CI
runner cannot hold 8 GiB to do it. [\texttt{Params::TESTING}] is the same
construction at 1/64th the cache and 1/512th the dataset, which builds in
milliseconds and exercises every branch the mainnet sizes do. The alternative
— testing the 4 GiB path only under \texttt{\#[ignore]} — means the determinism
property is checked by hand, on someone's workstation, if at all.

\subsubsection{Why the counts are prime}

Both counts are rounded *down* to a prime. A composite count shares factors
with the strides the mixing function produces, which lets the walk fall into
cycles far shorter than the array: the cache would be nominally 64 MiB and
effectively a fraction of it, the memory-hardness claim would be false, and
every test would still pass. Ethash rounds to a prime for the same reason.

\subsection{crypto/dag/registry.rs}
Which epoch's cache a block is checked against, and how many of them a node
keeps.

\subsubsection{Why more than one}

A block is validated against the cache for *its own* epoch, and a reorg can
cross an epoch boundary: a chain that reorganises from height 30,001 back to
29,998 and forward down another branch needs epoch 0 and epoch 1 within a
few milliseconds of each other. Holding one cache would mean regenerating —
seconds of CPU — inside fork choice, every time a boundary was recrossed,
which is a denial of service a peer could trigger deliberately.

Two is the whole requirement and the whole policy: the current epoch and the
one before it. At [\texttt{Params::MAINNET}] that is 128 MiB, held by every
validating node.

\subsubsection{Why generation happens under the lock}

[\texttt{CacheRegistry::cache\_for\_epoch}] holds its mutex across generation, so a
second thread asking for the same epoch waits rather than building a second
copy of the same 64 MiB array. The wait is bounded by one generation — about
a second — and only ever happens on the first block of an epoch, because a
node that follows the chain has already pulled the next epoch in through
[\texttt{CacheRegistry::prepare}]. The alternative, dropping the lock and letting
both threads build, spends 64 MiB and a second of CPU to avoid a wait of the
same length.

\subsection{crypto/hybrid.rs}
Hybrid ML-DSA-65 + SLH-DSA-SHA2-128s transaction authorization.

Every transaction on this chain carries two signatures over the same bytes:
a lattice proof under FIPS 204 and a hash-based proof under FIPS 205. Both
must verify. Neither is a fallback for the other.

\subsubsection{Why two}

ML-DSA-65 is already post-quantum, so a second post-quantum scheme looks
redundant. It is not, because the two are post-quantum for unrelated reasons.

ML-DSA rests on Module-LWE — a structured lattice assumption about fifteen
years old, believed hard, not proven hard. The ring structure that makes it
fast is the same structure a future attack would attack. SLH-DSA rests on
nothing but the preimage and collision resistance of SHA-2: no algebra, no
structure, and a security argument that predates the chain by decades.

A forgery must therefore break lattices *and* hashes. A cryptanalytic
advance against either one alone leaves the ledger intact, which is the only
reason the cost below is worth paying.

\subsubsection{What it costs}

Measured by \texttt{cargo bench --bench hybrid\_signing}, over a real transaction's
signing bytes:

\begin{verbatim}
| | ML-DSA-65 | SLH-DSA-SHA2-128s | hybrid |
|---|---|---|---|
| public key | 1952 B | 32 B | 1984 B |
| signature | 3309 B | 7856 B | 11165 B |
| keygen | 0.14 ms | 14 ms | 14 ms |
| sign | 0.22 ms | 105 ms | 105 ms |
| verify | 0.099 ms | 0.107 ms | 0.18 ms |
\end{verbatim}

Signing is roughly five hundred times more expensive than it was.
Verification is under twice. That asymmetry is the whole design: signing
happens once, in a wallet, per transaction; verification happens on every
node, for every transaction, in every block, forever.

The verify column also explains the choice of the \texttt{s} parameter set over
\texttt{f}. Measured side by side, \texttt{f} signs twenty times faster but verifies at
0.56 ms against 0.22 ms and carries a 17088-byte signature — it optimizes
the operation that happens once at the expense of the one that happens
forever. See \texttt{maya\_crypto\_pq} for that comparison.

\subsubsection{The address must commit to both keys}

This is the part that is easy to get wrong, and getting it wrong makes the
entire feature decorative.

Suppose the address kept committing only to the ML-DSA key, as it did in v2.
An adversary who broke Module-LWE could forge the lattice half for any
account — and then simply *generate their own SLH-DSA keypair*, sign with
it, and attach it to the transaction. Both checks pass. The hash-based half
would have added 7856 bytes and zero security, because nothing pinned it to
the victim.

So an address is \texttt{blake3(v3 ‖ ml\_dsa\_pk ‖ slh\_dsa\_pk)}, and
[\texttt{HybridPublicKey}] travels as a unit. Breaking one scheme now yields a
forged signature under a key that hashes to a *different* address, which
spends nothing.

[\texttt{crate::core::Transaction::signing\_bytes}] additionally commits to both
keys, so each signature covers the other's key: neither half can be lifted
out of one transaction and replayed beside a different partner key.

\subsubsection{Consequence: every address changed}

v2 addresses are hashes of an ML-DSA key alone and have no v3 preimage. This
is a hard fork with a new genesis, not a migration. There is deliberately no
compatibility path: accepting a v2 address would mean accepting a
single-signed transaction, which is the exact thing this module exists to
prevent.

\subsection{crypto/keys.rs}
ML-DSA-65 key handling.

Thin wrappers over \texttt{fips204} that keep the rest of the crate free of direct
dependency on its types where practical.

\subsubsection{This is one half of a signature}

Nothing on the chain is authorized by ML-DSA alone. Every transaction
carries a second, hash-based proof under FIPS 205, and an address commits to
both keys — see [\texttt{crate::crypto::hybrid}], which is the module callers
should reach for. The types here are its lattice half, exposed because the
keystore and the benchmarks need to name the two halves separately.

Consequently there is no \texttt{address()} on these types. An ML-DSA key does not
name an account by itself, and a function that pretended otherwise would be
the easiest way to reintroduce the attack \texttt{hybrid} exists to close.

\subsubsection{Why ML-DSA-65 and not ed25519}

An adversary with a cryptographically relevant quantum computer recovers an
ed25519 private key from its public key in polynomial time. That matters more
for a chain than for a session protocol: a transaction signature is published
permanently, and every address that has ever sent a transaction has therefore
published the public key that a future adversary would attack. "Harvest now,
decrypt later" against a ledger is "harvest now, *spend* later".

ML-DSA-65 is the FIPS 204 parameter set at NIST security category 3. The cost
is bulk and speed, both of which are real:

\begin{verbatim}
| | ed25519 | ML-DSA-65 |
|---|---|---|
| public key | 32 B | 1952 B |
| signature | 64 B | 3309 B |
| secret key | 32 B | 4032 B |
\end{verbatim}

\subsubsection{Scope of the guarantee}

This module makes *transparent transaction authorization* quantum-resistant.
The shielded pool's half is \texttt{maya\_zk\_stark::pool}: a hash-based STARK with no
setup (ADR-008), replacing the Groth16 pool whose pairing-based soundness a
quantum adversary could break. What that half still lacks is an independent
audit of its circuit — \texttt{CIRCUIT\_IS\_AUDITED} — not a post-quantum argument.

\subsubsection{Deterministic signing}

Signatures are produced with the FIPS 204 deterministic variant — the \texttt{rnd}
input fixed to 32 zero bytes — rather than the hedged default.

This is forced by consensus, not preference. [\texttt{crate::core::Transaction::txid}]
hashes the signature, and the txid is what the block Merkle root commits to.
Under hedged signing, signing one transaction twice yields two different
txids for the same authorized payload; a wallet that retried would produce a
transaction the network treats as unrelated to the first. Ed25519 is
deterministic, so this preserves the property the chain already relied on
rather than introducing a new one.

The trade is that deterministic signing is the weaker choice against fault
injection: an attacker with physical access who can glitch the signer while
it signs the same message twice can recover the key. That attack needs the
hardware in hand, and it is the same exposure ed25519 has always had here.
The alternative — hedged signing plus removing the signature from the txid —
would make two distinct signatures over one payload share a txid, which is
precisely the malleability the wire format is built to prevent.

\subsection{crypto/lattice/mod.rs}
\begin{framed}\noindent\textbf{Status.} This module is a research branch. It is present in the tree and no consensus path reaches it. Nothing described here is part of the running chain.\end{framed}

Turning a block header into the lattice that block must be mined against.

\subsubsection{Status}

\textbf{Research branch. Nothing in consensus calls this.} It is the node-side
half of [\texttt{maya\_lattice\_pow}], which carries the verification rules and the
reasons they are in a crate of their own. The consensus questions — chiefly
that lattice reduction is not progress-free — are open, and are recorded in
that crate's documentation and in \texttt{docs/lattice-pow.md}.

\subsubsection{Why the derivation lives here and not there}

\texttt{maya-lattice-pow} has no dependencies, so that Kani can compile it. That
rules out BLAKE3, and therefore rules out deriving anything from a header.
The same boundary \texttt{maya-dex} lives behind: the arithmetic that decides
validity is dependency-free and provable, and the hashing that feeds it
happens in \texttt{custom-l1-node}, which already links BLAKE3.

\subsubsection{Why the miner must not choose the lattice}

This is the property the whole scheme rests on. Finding a short vector is
only hard in a lattice you were handed; a miner who chose the basis would
choose one they had already reduced, and mine every block for free.

So the coefficients come from the header and nowhere else, through a keyed
BLAKE3 XOF. The domain string is consensus in exactly the way the layer
domains in [\texttt{crate::state::proof}] are: one definition, and both the prover
and the verifier reach it.

\subsubsection{Rejection sampling, not modular reduction}

Each coefficient is drawn as a \texttt{u32} and kept only if it is below the
largest multiple of \texttt{q} that fits in \texttt{u32}. Taking \texttt{draw \% q} instead would
be biased towards the low residues whenever \texttt{q} does not divide \texttt{2\textasciicircum{}32}, and
a lattice family with a skewed coefficient distribution is a lattice family
whose hardness estimates do not apply to it.

The loop is bounded: [\texttt{MAX\_DRAWS\_PER\_COEFFICIENT}] attempts, after which
derivation fails rather than spinning. With the smallest permitted modulus
the rejection probability per draw is below \texttt{2\textasciicircum{}-31}, so the bound is
unreachable in practice and exists so that a malformed parameter set cannot
hang a validating node.

\subsection{crypto/mod.rs}
Cryptographic primitives: the ArgonBlake hash, proof-of-work evaluation,
and hybrid ML-DSA-65 + SLH-DSA-SHA2-128s key handling.

[\texttt{hybrid}] is the module callers want. [\texttt{keys}] is its lattice half and
\texttt{maya\_crypto\_pq} is its hash-based half; both are exposed because the
keystore and the benchmarks need to name the two schemes separately, but
neither authorizes anything on its own.

\subsection{crypto/pow.rs}
Proof-of-work target evaluation.

A target is a 256-bit big-endian threshold. A digest satisfies the target
when, interpreted as a big-endian integer, it is less than or equal to that
threshold. Because both values are 32 bytes, lexicographic byte comparison
is exactly integer comparison — no bignum arithmetic required.

\subsection{crypto/suites.rs}
The node's side of ADR-007: suite ids, addresses, and the activation gate.

The registry, the envelope and the migration engine live in
\texttt{maya\_crypto\_pq}; this module is the thin layer that decides *when* the
chain accepts them. Until [\texttt{SUITE\_ENVELOPE\_ACTIVATION\_HEIGHT}] every
suite-tagged transaction is refused, so the v5/v6 hybrid format stays the
only one any block can contain.

[\texttt{SUITE\_ENVELOPE\_ACTIVATION\_HEIGHT}]: crate::state::context::SUITE\_ENVELOPE\_ACTIVATION\_HEIGHT

\subsection{error.rs}
Error types for the node.

Every fallible operation returns [\texttt{Result}]; no production path panics.

\subsection{genesis.rs}
Genesis configuration.

A testnet's genesis must be reproducible: every node has to derive an
identical genesis block and identical starting balances from the same file,
or they will disagree about the chain from block one and never converge.

Reproducibility here rests on three things:

- \textbf{Allocations are sorted by address} before hashing, so the JSON file's
  ordering — which a human editing it will not preserve — cannot change the
  state root.
- \textbf{The state root is computed from the config alone}, with the same Merkle
  construction the live chain uses. A node can therefore verify that what it
  wrote to disk matches what the config declared.
- **\texttt{chain\_id} is folded into the genesis \texttt{prev\_hash}**, so two networks with
  different ids can never share a genesis block, and a block mined for one
  cannot be replayed onto the other.

\subsection{governance/mod.rs}
On-chain governance: what the chain stores about proposals, stake, and the
rules currently in force.

The decisions live in [\texttt{maya\_governance}]. This module is the part that has
to know about storage — how a record is laid out, what key it files under,
and how it enters the Merkle state root.

\subsubsection{Voting power}

Two sources, added:

- \textbf{Locked stake.} Native coin debited into a lock record with an unlock
  height. The lock must outlive the proposal's *execution*, not merely its
  vote — otherwise the cheapest way to decide something is to acquire
  weight, vote, and be gone before the decision binds anyone who stayed.
- \textbf{Recent work.} See [\texttt{work}]. This chain has no block reward and no miner
  field in the header, so there was no on-chain record of who mined
  anything; [\texttt{crate::core::governance\_payload::WorkClaim}] creates one
  without touching the proof-of-work preimage.

Quorum is measured against locked stake alone, for the reason
[\texttt{maya\_governance::tally}] gives: work credit decays per address at a
different moment for each address, so a chain-wide total of it would drift
from the sum of its parts.

\subsubsection{One keyspace, one prefix}

Everything here lives under \texttt{g:}, joining the trading subsystem's \texttt{d:} and
the oracle's \texttt{o:}. No pre-existing prefix begins with that byte, so one scan
collects the whole subsystem — which is what lets it fold into the state
root as a single layer and journal into the undo record with no new section
at all.

\subsection{governance/record.rs}
What the chain stores about proposals, stake, and the rules in force.

\subsection{governance/work.rs}
Recent hash power, as a number the chain can hold.

\subsubsection{The problem this solves}

This chain has no block reward. \texttt{bins/pool-service/src/config.rs} states it in
the code: \texttt{apply\_block\_checked} credits no subsidy and fees burn to the fee
sink. There is no coinbase, and [\texttt{crate::core::BlockHeader}] carries
\texttt{prev\_hash}, \texttt{state\_root}, \texttt{timestamp}, \texttt{nonce}, and \texttt{difficulty\_target} —
nothing that says who mined it.

So before this module there was \textbf{no on-chain record of who did any work},
and "voting power based on hash power" had no data source at all.

\subsubsection{Why a claim transaction rather than a header field}

Adding a miner field to the header would change the proof-of-work preimage,
and therefore the miner, the pool protocol, the Stratum implementation, and
the CUDA kernel — \texttt{core::block::NONCE\_RANGE} exists precisely because those
four have to agree byte for byte about the header's layout.

A claim transaction changes none of them. At most one is admitted per block
and it credits that block's work to an address the *miner* names, because
the miner is the party that chose what went in the block. That is the
correct answer to "who did this work" and it needs no new consensus field.

The honest consequence: hash-power voting is \textbf{miner-directed}, not
pool-participant-directed. A pool's miners contribute the hashes; the pool
operator decides who is credited. That is a real property of this design and
not one the mechanism can fix — attributing work to the individual hashers
would require the share data the pool holds off chain.

\subsubsection{Why credit decays}

Voting power should reflect hash power *now*. A miner who ran a farm for a
month two years ago is not securing the chain today, and an all-time
accumulator would let them govern it forever.

Decay is by halving, not by a multiplier: a value loses half its weight
every [\texttt{WORK\_HALF\_LIFE\_BLOCKS}]. Halving is exact in integer arithmetic —
it is a shift — where a fractional decay would need a rounding rule, and a
rounding rule in voting weight is a bias in whichever direction it falls.

\subsection{iso20022\_bridge.rs}
The chain side of the ISO 20022 bridge: bank accounts to chain addresses,
and committed state back to a statement.

\texttt{maya\_iso20022} deliberately knows nothing about blocks. This module is the
other half of that seam — it is where a [\texttt{PaymentIntent}] becomes something
the chain can hold, and where committed state becomes a \texttt{camt.053} a
counterparty can reconcile against.

\subsubsection{Virtual accounts}

A bank account has no Maya2C keypair. So each one is mapped to an address
\textbf{derived} from its list identifier rather than generated: the same IBAN
always maps to the same address, on every node, without anybody storing a
table. That matters twice over — a table is state that can disagree between
nodes, and a table is a thing an operator can edit to redirect a payment.

Nobody holds the private key to a derived address, which is the point:
value at one moves only through the bridge, because there is no signature
that could move it any other way. It also means value there is
\textbf{unrecoverable} if the bridge is switched off, which is stated here rather
than discovered.

\subsubsection{What a statement is drawn from}

[\texttt{StatementSource::for\_block}] reads a committed block only. It never reads
an overlay, a mempool, or a block that has not applied — a
statement is what a bank reconciles against, and reporting a payment that a
reorg then removes is worse than reporting it late. Every entry it renders
is \texttt{BOOK}, because there is no other status a committed block can produce.

\subsection{metrics/mod.rs}
Prometheus instrumentation.

\subsubsection{Why metrics live on their own port}

Peer topology and mempool contents are operational data, not public data.
Serving them alongside JSON-RPC would publish the node's peer set and its
pending transactions to anyone who can reach the RPC endpoint, which for a
load-balanced node is the whole internet. The exporter therefore binds
separately and is expected to stay inside the pod network.

\subsubsection{Measuring propagation honestly}

"Block propagation latency" is usually computed as the gap between a block's
header timestamp and the moment a node saw it. That number is real, but it
measures clock skew between the miner and the observer just as much as it
measures network delay — an operator whose NTP has drifted will see a
latency spike caused entirely by their own clock.

So both are exported, under names that say which is which:

- [\texttt{Metrics::observe\_block\_age}] records \texttt{maya\_block\_observed\_age\_seconds},
  the skew-sensitive end-to-end figure.
- [\texttt{Metrics::observe\_import}] records \texttt{maya\_block\_import\_duration\_seconds},
  measured entirely against the local clock and therefore skew-free.

Alert on the second; watch the first for trends.

\subsection{metrics/server.rs}
The \texttt{/metrics} exporter.

A deliberately minimal HTTP server. It answers \texttt{GET /metrics} and
\texttt{GET /healthz} and nothing else — every other path is a 404. Prometheus
needs no more than that, and an exporter with a wider surface is a wider
surface on an endpoint that carries operational data.

Hyper is used directly rather than a framework because the node already
pulls it in through jsonrpsee, so this costs no new dependency.

\subsection{network/behaviour.rs}
Composed network behaviour: gossipsub, Kademlia, and identify.

The three compose deliberately:

- \textbf{gossipsub} broadcasts blocks and transactions.
- \textbf{kad} provides dynamic peer discovery so the mesh is not statically
  configured.
- \textbf{identify} is what makes Kademlia usable in practice: a peer learns its
  dialable addresses from what other peers observe, and those addresses are
  what get inserted into the routing table. Without it, Kademlia knows peer
  IDs it has no way to reach.

\subsection{network/evidence\_tap.rs}
Carries each inbound gossip message's author signature to the driver, so a
rejected message can become evidence.

Gossipsub checks the signature in its codec — Strict mode, before the
behaviour sees the message — and then drops it: the [\texttt{Message}] handed to
the application has no signature field. This [\texttt{DataTransform}], the one hook
that sees the [\texttt{RawMessage}], puts the signature in front of the message
data, and the driver takes it off again with [\texttt{split}].

In the message rather than in a table beside it, on purpose. A table keyed
by author and sequence number is shared state that a flood of validly signed
messages can evict before the driver claims an entry, which silently turns
evidence capture off. A signature riding inside its own message cannot be
evicted, is gone once the message is judged, and exists only for messages
gossipsub actually delivers.

Nothing else changes. Gossipsub's default message id is the author plus the
sequence number, never the data, and forwarding uses the raw message, not
this one.

\subsection{network/gossip\_score.rs}
Gossipsub's own peer scoring, configured to punish what is provable and
nothing else.

[\texttt{crate::network::peer\_health}] decides quarantine. This is the second,
cheaper layer: gossipsub keeps a score per peer for mesh management, and a
peer below its thresholds stops receiving gossip from us, then stops being
heard (graylisting). Left at its defaults it would also score *quietness* —
mesh delivery deficits, time in mesh, slow delivery — and a quiet or distant
honest peer would sink. So those components are switched off and only P4,
invalid message deliveries, carries weight: a peer loses score here only for
messages this node rejected.

\subsection{network/identity.rs}
The persisted libp2p identity.

\subsubsection{Why this is not private to the node binary}

A seed's \texttt{PeerId} is derived from this key, and cross-region bootnode
multiaddrs are written against that \texttt{PeerId} — they are baked into manifests
long before the pod that owns the key first starts. Something other than the
node has to be able to read the key and report the identity, or the fleet's
bootnode list can only be discovered by reading logs after the fact.

So the loader lives in the library and both \texttt{node} and \texttt{peerid} call it.
Two copies of this logic that disagreed by a byte would produce two different
\texttt{PeerId}s from the same file, and the symptom would be a seed nobody can dial.

\subsubsection{Format}

Exactly 32 bytes: the raw ed25519 secret scalar. Not PKCS\#8, not hex — a
fixed-size opaque blob is what \texttt{kubectl create secret generic --from-file}
round-trips without transformation, which is how these keys reach a cluster.

\subsection{network/mempool.rs}
Thread-safe transaction pool.

\subsubsection{Locking}

The map is a [\texttt{std::sync::RwLock}], not a \texttt{tokio::sync::RwLock}. Every
critical section here is a pure in-memory map operation with no \texttt{.await}
inside it, so the async-aware lock would buy nothing and cost more. The
RocksDB reads that back validation happen *before* the write lock is taken —
holding a std lock across blocking I/O is exactly the mistake this ordering
avoids.

\subsubsection{Validation}

Inbound gossip is untrusted. A transaction is admitted only if it is
well-formed, correctly signed, not already pooled, and consistent with
committed state: nonce not stale, balance sufficient.

\subsection{network/mod.rs}
P2P networking: gossipsub broadcast, Kademlia discovery, and the mempool
that gates inbound transactions.

\subsection{network/node.rs}
Node swarm construction and the event loop driving it.

A node owns a libp2p \texttt{Swarm} and a [\texttt{Mempool}]. Because a \texttt{Swarm} is not
\texttt{Sync} and must be polled from a single place, the node runs as one spawned
task and is controlled through an mpsc command channel; observers subscribe
to a broadcast channel of [\texttt{NodeEvent}]. That keeps all swarm access on one
thread while still letting many callers publish and watch.

\subsection{network/node/build.rs}
Building a [\texttt{Node}]: the transport, the upgrade chain, and the swarm.

Split from \texttt{node.rs} to keep that file within the project's size limit, as
\texttt{guard.rs} and \texttt{relay.rs} were. A child module, so it can fill in the node's
private fields.

\subsection{network/node/guard.rs}
The driver's half of the peer guard: gossip verdicts, quarantine, and block
sync.

What is scored, and why lateness and double proposals are not, is in
[\texttt{crate::network::peer\_health}]. The fetch protocol and its bounds are in
[\texttt{crate::network::sync}]. Split from \texttt{node.rs} to keep that file within the
project's size limit; it is a child module so it can reach the driver's
private state.

\subsection{network/node/relay.rs}
The driver's half of the block relay.

The relay itself — sealed chunks, the XDP program, AF\_XDP — is
\texttt{maya\_ebpf\_net}; the key exchange is [\texttt{crate::network::relay\_key}]. This
module joins them to the node:

- agrees keys on each new connection, and forgets a peer's keys when its
  last connection closes;
- seals a published block into datagrams for every keyed peer, on the
  blocking pool rather than the swarm task;
- takes reassembled bodies from the receive path and hands them to
  [\texttt{NodeDriver::gossiped\_block}], the function a gossiped block reaches;
- mirrors peer-guard quarantines into the XDP blocklist.

\subsubsection{A relay decides nothing}

A relayed body is validated by the same code as a gossiped one and emits the
same [\texttt{NodeEvent::BlockReceived}]. The only rule here that gossip lacks is
framing: a peer that names one block id and delivers another is scored,
because an honest sender never does and a dishonest one could otherwise
re-deliver one body under endless fresh ids.

\subsubsection{One hop}

A node relays the blocks it publishes and never re-relays one it received.
Gossip carries a block onward after validating it; a relay that forwarded
on receipt would turn one block into thousands of datagrams per hop ahead
of any validation.

\subsection{network/node/threat.rs}
The driver's half of the threat-intel registry: turning a rejected gossip
message into evidence, and remembering where peers connected from.

Split from \texttt{node.rs} for size, as \texttt{guard.rs} and \texttt{relay.rs} were.

The node emits evidence; it never submits it. Submitting needs a funded
chain key, and the node holds none — so an operator's tool, or a test, turns
[\texttt{NodeEvent::AttackEvidence}] into a signed \texttt{TxKind::AttestAttack}.

\subsection{network/peer\_health.rs}
Evidence that a peer is Byzantine, and quarantine.

\subsubsection{What counts}

Only what a peer is provably responsible for. With gossip validation on
(\texttt{behaviour.rs}), an honest node forwards a message only after it checked
it, so a message that fails the same stateless check came from a peer that
either wrote it or relayed it unchecked. Either is the peer's fault.

\begin{verbatim}
| Offence | Evidence | Weight |
|---|---|---|
| [`Offence::MalformedFrame`] | gossip bytes that do not decode | 25 |
| [`Offence::InvalidSignature`] | a transaction whose hybrid signature fails | 50 |
| [`Offence::TxRootMismatch`] | a block body that disagrees with its header — invariant 24's relay substitution | 50 |
| [`Offence::InvalidBlock`] | a block the chain refuses for a reason inside the block | 50 |
| [`Offence::BadSyncResponse`] | a sync answer carrying blocks nobody asked for, or undecodable ones | 50 |
\end{verbatim}

\subsubsection{What does not, and why}

- \textbf{Late block propagation.} A distant honest peer is late, and scoring
  lateness would let an attacker get honest peers quarantined by being
  closer — the first step of an eclipse. It is recorded for monitoring
  ([\texttt{PeerHealth::observe\_latency}]) and never scored.
- \textbf{Double proposals.} A proof-of-work block has no proposer. Two blocks at
  one height are a fork, forks are normal, and relaying both is honest.
- \textbf{Transactions refused for state reasons} — a stale nonce, an
  insufficient balance, a full pool. Honest relays race each other, so these
  are \texttt{Ignore}, not \texttt{Reject}.
- \textbf{Unknown parents.} A block whose parent has not arrived is out of order,
  not wrong.

\subsubsection{Quarantine}

A score decays with a half-life. Crossing [\texttt{GuardConfig::quarantine\_score}]
quarantines the peer: the driver blacklists it in gossipsub (every message
it relays or authored is dropped), blocks its connections, and disconnects
it. Quarantine expires, because a node that bans forever eventually bans an
honest peer that was briefly misconfigured; each repeat doubles the length
up to a cap. This is all node-local and changes no consensus rule.

\subsection{network/pq/dual.rs}
\begin{framed}\noindent\textbf{Status.} This module ships disabled. It is reachable only when an operator opts in, and a stock node behaves as though it were absent.\end{framed}

The dual-KEM combiner: one session key from two unrelated hard problems.

\subsubsection{Status}

\textbf{Off by default.} \texttt{/maya/dualkem/1.0.0} is offered alongside
\texttt{/maya/mlkem/1.0.0} and negotiated, and a node only offers it when its
configuration says so. See [\texttt{DualKemPolicy}] for why it ships disabled and
what has to happen before that changes.

\subsubsection{What the combiner buys}

Confidentiality that survives a break in *either* primitive. ML-KEM rests on
Module-LWE, a structured lattice assumption; HQC rests on decoding random
quasi-cyclic codes. An advance against one is not an advance against the
other, so a channel keyed from both stays secret unless both fall — the same
reasoning \texttt{docs/hybrid-signatures.md} gives for ML-DSA plus SLH-DSA on every
transaction, applied to the transport.

\subsubsection{What it costs, and this is the part to read}

\textbf{Confidentiality becomes an OR; availability becomes an AND.}

The channel is secure if either KEM holds. But the connection *works* only
if both implementations are correct. A panic, a mis-derived secret, or an
encoding change in either one breaks every connection that negotiated this
protocol. \texttt{ml-kem} is pinned at a released 0.3.2 against a final FIPS 203;
\texttt{hqc-kem} is a release candidate tracking a \textbf{draft} FIPS 207.

Pairing them means the transport's liveness is gated by the less-settled of
the two, in order to defend the transport's secrecy against a break in the
more-settled one. That can be the right trade. It is not a free one, and it
is why this is negotiated rather than required.

\subsubsection{Why a KDF combiner and not double encryption}

"Requires both secrets" could be built by encrypting twice, once under each
KEM's secret. That is worse in every respect: it doubles the AEAD cost on
every frame, it doubles the nonce management surface, and it has no clean
security argument — nesting two AEADs is not a standard construction.

A KDF over both secrets gives the property outright. The session key depends
on both inputs, so an adversary missing either one faces a PRF with an
unknown input, and one AEAD pass per frame is unchanged from today.

\subsubsection{Why the transcript, and not just the secrets}

The combiner hashes \textbf{both shared secrets and both full transcripts} —
encapsulation keys and ciphertexts alike. Hashing the secrets alone is not
sound in general: a combiner over shared secrets is IND-CCA robust only if
the underlying KEMs are ciphertext-collision-resistant, which neither
specification promises. Binding the ciphertexts removes the assumption
rather than relying on it.

The single-KEM handshake already did this, so the extension inherits the
property rather than inventing it.

\subsubsection{Why the context string changes}

\texttt{KDF\_CONTEXT} is not the single-KEM one. If both protocols derived under
the same context, a dual session and an ML-KEM-only session that happened to
share an ML-KEM transcript would derive related keys, and the second KEM
would be contributing nothing at exactly the moment it was supposed to
matter. A new protocol gets a new context; that is what the string is for.

\subsection{network/pq/handshake.rs}
The ML-KEM-768 key exchange, run inside an established Noise session.

\subsubsection{The exchange}

Two messages, both fixed-width, because ML-KEM has no variable-length
encodings and a length prefix would only add a field an attacker can lie
about:

\begin{verbatim}
responder → initiator   [1 byte version][1184 byte encapsulation key]
initiator → responder   [1 byte version][1088 byte ciphertext]
\end{verbatim}

The responder generates a keypair, sends the public half, and decapsulates
what comes back. The initiator encapsulates to it. Both then hold the same
32-byte secret — or, if anything was tampered with, two different secrets
and a connection that fails on its first frame.

\subsubsection{Why this is safe to keep so small}

It is not an authenticated key exchange, and it does not try to be. It runs
*inside* a completed Noise XX session, which has already authenticated the
peer against its \texttt{PeerId} and encrypted the channel. An active attacker
cannot reach these bytes without first breaking X25519 in real time.

So the job here is narrow: add a secret that a *recording* adversary cannot
recover later, no matter how long they wait. That is the harvest-now-
decrypt-later threat, and it is the one that actually applies to a ledger's
gossip traffic, because that traffic is worth reading years after capture.

Writing a full post-quantum AKE here instead would mean hand-rolling an
authenticated protocol, which is the thing not to do casually.

\subsubsection{Transcript binding}

The derived keys commit to both handshake messages, not just the shared
secret. Without that, an attacker who could reflect or splice messages
between two concurrent connections might get two sessions to derive the same
key from the same encapsulation. Hashing the transcript makes each session's
key a function of the exact bytes that session exchanged.

\subsubsection{Implicit rejection, and why nothing branches on it}

ML-KEM decapsulation never fails. On a forged or corrupted ciphertext FIPS
203 returns a pseudorandom value derived from the ciphertext and a per-key
rejection seed, rather than an error. There is deliberately no check here
that decapsulation "succeeded", because it always does — a mismatch shows up
as the two sides deriving different keys, and the first encrypted frame
failing to authenticate.

\subsection{network/pq/measure.rs}
Handshake cost measurement: bytes on the wire, segments, and wall clock.

\textbf{Test and benchmark infrastructure.} It lives in the library rather than
in \texttt{tests/} for the same reason [\texttt{crate::network::sim}] does: an integration
test cannot reach a \texttt{\#[cfg(test)]} item. Nothing on a production path calls
any of it.

\subsubsection{Why this measures rather than asserts}

The dual-KEM work was scoped against a sub-100 ms connection target. A test
asserting that threshold would be worse than no test:

- On loopback or an in-process duplex, \textbf{both} handshakes finish in well
  under a millisecond. The assertion would pass forever, including after a
  change that tripled the cost, and would report a green tick for a
  property it never examined.
- On a real link the number is dominated by round trips, which are a
  property of the path and not of this code. A CI machine's loopback cannot
  observe it, and a threshold tuned until CI passed would be measuring the
  CI runner.

So this reports. [\texttt{Measurement}] carries the numbers; the caller decides
what to make of them, and \texttt{tests/dualkem\_latency\_tests.rs} prints a table.

\subsubsection{The number that actually matters, and the trap in it}

Segment count — but of the \textbf{largest single message}, not of the handshake
total. Congestion control is per-direction, so the sum of both directions is
the wrong quantity, and dividing it by the MSS gives an answer that is wrong
in the alarming direction.

\begin{verbatim}
| | Wire total | Largest message | Segments | Fits window? |
|---|---|---|---|---|
| `/maya/mlkem/1.0.0` | 2,274 | 1,185 | 1 | yes |
| `/maya/dualkem/1.0.0` | 15,766 | 5,699 | 4 | yes |
\end{verbatim}

The initial congestion window is [\texttt{INITIAL\_CWND\_SEGMENTS}] segments
(RFC 6928). \textbf{Both handshakes fit}, so neither pays an extra round trip.

This is worth stating loudly because the design discussion assumed the
opposite: 15,766 divided by 1,460 is eleven, which exceeds ten, and that
arithmetic is easy to do and wrong. [\texttt{Measurement::exceeds\_initial\_window}]
takes the largest message precisely so the mistake is hard to repeat.

The extra bytes are still real. They cost bandwidth on every connection, and
on a constrained or metered link that is the cost worth counting. They do
not cost latency.

\subsection{network/pq/mod.rs}
Post-quantum confidentiality for the P2P transport.

\begin{verbatim}
TCP / Memory
  └── libp2p-noise   Noise_XX_25519_ChaChaPoly_SHA256   ← unchanged
        └── /maya/mlkem/1.0.0                           ← this module
              ML-KEM-768 (FIPS 203) + ChaCha20-Poly1305
                └── yamux → gossipsub
\end{verbatim}

\subsubsection{Why a layer and not a replacement}

The obvious implementation would be a post-quantum Noise pattern —
\texttt{Noise\_XXhfs\_25519+MLKEM768\_ChaChaPoly\_SHA256}. It is not reachable from
here. \texttt{libp2p-noise} hard-codes its parameters in a private constant behind
a private resolver, and \texttt{Config} exposes no way to change the pattern or the
key exchange. One layer down, \texttt{snow} does expose a \texttt{Kem} trait under its
\texttt{hfs} feature — but \texttt{snow::params::KemChoice} is a closed enum with exactly
one variant, \texttt{Kyber1024}, so a pattern string naming ML-KEM-768 does not
parse. An ML-KEM implementation could be registered under the name
\texttt{Kyber1024}, and that was rejected: a handshake whose advertised protocol
name misstates its primitive is a trap for whoever audits it next.

The remaining options were to fork two security-critical crates, or to write
a full authenticated key exchange from scratch. This layer is the third: it
leaves the audited Noise handshake exactly as it is and adds a small,
self-contained post-quantum stage above it.

\subsubsection{What this protects against}

\textbf{Harvest now, decrypt later.} An adversary recording gossip traffic today
and waiting for a quantum computer must break ML-KEM-768 *and* X25519 to
read any of it. This is the threat that actually applies to a chain: gossip
reveals which peer originated which transaction, and that metadata does not
become less interesting with age.

\subsubsection{What it does not}

\textbf{Authentication stays classical.} The peer is authenticated by the Noise
layer below, against its ed25519 \texttt{PeerId}, and nothing here changes that. An
adversary holding a cryptographically relevant quantum computer *at the
moment a connection is made* could still impersonate a peer. That is an
active attack requiring the machine to exist and be on the wire — a
different and much harder proposition than recording packets — but it is a
real limit and it is not closed here. Closing it means post-quantum peer
identity, which changes what a \texttt{PeerId} is and is separate work.

\subsubsection{No downgrade}

The upgrade is not optional. A peer that does not speak \texttt{/maya/mlkem/1.0.0}
fails multistream-select negotiation and the connection is dropped. There is
deliberately no fallback path: an optional post-quantum layer is one an
attacker strips.

\subsection{network/pq/rotation.rs}
Session rotation: when a connection's post-quantum keys are too old.

\subsubsection{What rotation is and is not for}

The ML-KEM keypair in [\texttt{super::handshake}] is generated fresh for every
connection and dropped the moment the handshake finishes. That is where
forward secrecy comes from, and it is already complete: an adversary who
seizes a node tomorrow cannot decrypt traffic recorded today, because the
decapsulation key that could have is long gone.

What this module adds is a bound on how long any one session key lives. A
node in a quiet mesh can hold a connection open for days; without rotation,
one handshake would protect days of traffic, and a key extracted from a live
process would unlock all of it. Forcing a re-handshake every epoch caps that
window. The property is \textbf{post-compromise security} — recovery after a
leak — not forward secrecy, which was never missing.

\subsubsection{Why the epoch is local, and not negotiated}

An earlier design carried the epoch in the handshake and had both peers
agree on it. That does not survive contact with a real network: peers are at
different heights, a syncing node is thousands of blocks behind, and a node
that has just started has no chain at all. Two honest peers would compute
different epochs and fail to connect — a self-inflicted partition, exactly
when the network is least able to afford one.

So the epoch never goes on the wire. Each node independently notes the epoch
a connection was established in and closes it once its own epoch has moved
on. libp2p redials, the upgrade runs again, and fresh keys are in place.
Whichever side notices first drives it; the other simply sees a reconnect.
Two peers disagreeing about the epoch now means one of them rotates sooner
than the other, which is harmless.

\subsection{network/pq/stream.rs}
The encrypted duplex that sits above the Noise session.

Once [\texttt{super::handshake}] has produced directional keys, every byte the node
sends is framed and sealed with ChaCha20-Poly1305 under a key an adversary
must break ML-KEM-768 to derive. The Noise layer underneath keeps doing its
job; this adds a second, post-quantum one on top.

\subsubsection{Framing}

\begin{verbatim}
[4 bytes big-endian length][ciphertext .. length bytes, tag included]
\end{verbatim}

Fixed 4-byte prefix rather than a varint: a varint saves three bytes on a
frame that is already up to 64 KiB, and costs an incremental decoder that
has to handle a hostile peer stretching a length across many reads.

The length covers the ciphertext *and* its 16-byte tag, so a frame is
self-delimiting and a truncated one fails to authenticate rather than being
silently accepted short.

\subsubsection{Nonces}

A 64-bit counter per direction, little-endian in the low 8 bytes of the
12-byte nonce, starting at zero and never reused. Two properties matter:

- \textbf{Separate counters per direction.} The two directions use different keys
  ([\texttt{super::handshake::SessionKeys}]), so even identical counters are safe;
  separate counters mean the two facts are independent rather than one
  propping up the other.
- \textbf{No wraparound, ever.} At \texttt{u64::MAX} this returns an error and the
  connection dies. Reaching it would take longer than the universe has
  existed at any plausible frame rate, but a counter that silently wrapped
  would repeat a \texttt{(key, nonce)} pair, and repeating one destroys
  ChaCha20-Poly1305's security outright. A hard failure is the only correct
  behaviour and it costs one comparison per frame.

\subsubsection{Hostile input}

Every length is checked against [\texttt{MAX\_FRAME\_LEN}] before a single byte is
allocated, mirroring the discipline in [\texttt{crate::core::codec}]. A peer that
announces a 4 GiB frame gets an error, not an allocation.

\subsection{network/radio\_gateway.rs}
The node side of the radio link: headers out, headers in, and a relay
queue in between.

\subsubsection{Why this is a gateway and not a libp2p transport}

The plan for this subsystem was a \texttt{TransportKind::Radio} beside \texttt{Memory} and
\texttt{Tcp}. Working through the link budget says it should not be, and the reason
is worth writing down rather than discovering later.

A libp2p connection opens with multistream-select, a Noise handshake and a
yamux negotiation. Maya2C's Noise is ML-KEM-768: a 1,184-byte encapsulation
key and a 1,088-byte ciphertext, before any protocol negotiation and before
one byte of payload. At SF12 the link carries 26 bytes of fountain payload
per frame and each frame costs 2.47 seconds of airtime against a 1\% duty
cycle — roughly 246 seconds of enforced silence apiece.

That handshake is therefore \textbf{hours}, repeated per connection, to move a
144-byte header. A transport that cannot complete its own handshake inside
the useful lifetime of what it carries is not a transport.

So the radio path is a gateway: it takes headers from the chain, puts them
on the air as fountain symbols, takes symbols off the air, and hands
recovered headers back. No session, no negotiation, no per-peer state that
has to survive a week of intermittent contact. This is the shape every
working LoRa mesh converges on, and it is what the architecture vision meant
by "low-bandwidth header relay".

\subsubsection{A gateway decides nothing}

\texttt{docs/architecture-vision.md} §3: a transport may not change consensus. So
this module hands a recovered header to exactly the code a TCP peer's header
reaches, and there is no branch anywhere that asks how it arrived. The one
thing it may do is *refuse to carry* — airtime is finite and somebody has to
choose — and that choice is by height, which is public and not a judgement
about validity.

\subsection{network/relay\_key.rs}
Relay key exchange: two contributions over a connection that is already
authenticated and post-quantum sealed.

The block relay (\texttt{maya\_ebpf\_net}) seals every chunk under a per-peer key.
That key needs three things: nobody but the two peers knows it, it binds the
peer the connection authenticated, and it is fresh for each connection. The
libp2p connection already gives the first two — Noise authenticates the
\texttt{PeerId}, and the ML-KEM layer makes what crosses it confidential against a
recording adversary — so the exchange rides on it rather than inventing a
handshake:

\begin{verbatim}
lower PeerId                          higher PeerId
RelayKeyRequest { contribution, port } ──►
                                      ◄── RelayKeyResponse { contribution, port }
both: derive_keys(request, response, requester id, responder id)
\end{verbatim}

Only the peer with the lower id asks, so two nodes that dial each other at
once do not run two exchanges and end up holding different keys.

\subsubsection{Where datagrams go}

A peer names a \textbf{port}, never an address. Relay datagrams are sent to the
IP the libp2p connection came from. Letting a peer name the address would
let it point a node's block bursts — thousands of datagrams per block — at
a victim: reflection, with the node as the amplifier. A peer with no IP
(the in-process memory transport) gets no relay.

\subsubsection{What this does not protect}

The contributions travel through the CBOR codec, which copies them into
buffers this crate cannot zeroize. The derived keys are zeroized on drop;
the contributions' copies in libp2p's buffers are not. They are only useful
to someone who can already read the node's memory.

\subsection{network/sim.rs}
Network simulation support: a transport that injects latency.

\textbf{This is test and benchmark infrastructure.} It lives in the library
rather than in \texttt{tests/} because an integration test cannot reach a
\texttt{\#[cfg(test)]} item, and building a node with a modified transport requires
a constructor the library exports. Nothing on a production path constructs a
\texttt{DelayStream}-wrapped transport, and [\texttt{crate::network::Node::new\_memory}] is unaffected.

\subsubsection{What it is for}

Every default test in this repository runs on libp2p's in-process memory
transport, where a write is a \texttt{Vec} push and latency is effectively zero.
That is the right default — it makes the tests fast and free of OS
flakiness — but it means the whole suite validates the protocol under
conditions no real network provides.

It matters more now than it used to. The post-quantum upgrade adds a
two-message round trip to every connection ([\texttt{crate::network::pq}]), and a
round trip is exactly the thing latency multiplies. On a zero-latency
transport that cost is invisible; at 200 ms RTT it is the dominant term in
connection setup. A simulation that never delays anything cannot tell the
difference between "the handshake works" and "the handshake works when it
costs nothing".

\subsubsection{What it models, and what it does not}

It delays reads. Each \texttt{poll\_read} waits before it resolves, which
approximates one-way propagation delay well enough to expose round-trip
sensitivity and to order events realistically.

It does not model bandwidth, jitter, packet loss, reordering, or queueing —
all of which a real network has and some of which change protocol behaviour.
A passing run here is evidence that latency alone does not break
propagation. It is not evidence that the protocol survives a bad network,
and the tests that use it say so.

\subsection{network/sync.rs}
Block sync: fetching a missing parent from the peer that relayed its child.

The one repair a peer can need that proof of work does not already give.
A node's state cannot be corrupted by a peer — every block is validated and
its state root checked (invariant 24) — but a node that received a child
before its parent holds an orphan until the parent arrives, and gossip does
not resend history. So the importer asks the relay for the missing ids, and
whatever comes back goes through \texttt{Chain::insert\_block} like any other block:
the heavier valid chain wins, and a lie is refused and scored.

There are no "consensus proofs" beyond that. A set of honest nodes vouching
for a state would need a known committee, and on an open proof-of-work
network a vote is free to forge. Work is the proof.

\subsubsection{Bounds}

Requests carry at most [\texttt{MAX\_BLOCKS\_PER\_REQUEST}] ids. A response stops
adding blocks at [\texttt{MAX\_RESPONSE\_PAYLOAD}] bytes and is refused on the wire
past [\texttt{MAX\_RESPONSE\_BYTES}]. Each peer may ask [\texttt{REQUESTS\_PER\_WINDOW}] times
per [\texttt{REQUEST\_WINDOW}]; beyond that it is answered with nothing.

\subsection{network/topics.rs}
Gossip topic identifiers.

\subsection{neural\_gas/features.rs}
Extracting the six features from a block.

Q16 fixed point throughout, each clamped to \texttt{±4.0}. The order is
consensus once the rule activates, and it matches \texttt{fee-market}'s
\texttt{model::Feature}:

\begin{verbatim}
| Index | Feature | Definition |
|---|---|---|
| 0 | fullness | block bytes / target |
| 1 | mean transaction size | (block bytes / transactions) / 16 KiB |
| 2 | access overlap | transactions naming an account another transaction names / transactions |
| 3 | fuel per byte | Σ declared `gas_limit` / block bytes / 1,000 |
| 4 | cross-shard | transactions whose accounts fall in more than one shard / transactions |
| 5 | size trend | (block bytes − previous block bytes) / target, signed |
\end{verbatim}

\subsection{neural\_gas/mod.rs}
Block features for the neural base-fee gain.

The network and the fee rule live in \texttt{fee-market} (\texttt{src/model/},
\texttt{src/rule.rs}), which nothing in this crate links: \texttt{tests/fee\_market\_tests.rs}
fails if any file under \texttt{src/} names that crate, and it stays true here.
This module does the half that needs chain types — reading a block — and
hands back six integers. The fee market's research branch reads them; no
consensus path calls either side. \texttt{tests/neural\_fee\_tests.rs} pins the
constants the two halves share.

\subsubsection{What the brief asked for, and what a validator can compute}

Every feature is a function of the block and nothing else: no clock, no
allocator, no node-local measurement. A fee that depended on something one
node measured differently from another would split the chain, for the reason
invariant 28 gives about the breaker.

\begin{verbatim}
| Brief | Here |
|---|---|
| tx size | mean transaction size and block fullness |
| state access overlap | accounts named by more than one transaction — the static access set, because execution builds one merged overlay per block and tracking per-transaction reads would allocate in the hot loop (execution directive 3) |
| memory allocation depth | **not computable deterministically**; replaced by declared contract fuel per byte |
| inter-shard dependencies | transactions whose accounts span more than one `blockgraph` shard. Hypothetical: no shards execute |
| historical DAG vertex metrics | the DAG here is the proof-of-work dataset, not a transaction graph; the features describe the parent block |
\end{verbatim}

\subsection{oracle/beacon.rs}
The randomness beacon.

\subsubsection{The construction}

\begin{verbatim}
alpha_n   = previous_beacon ‖ height_n
output_n  = VRF(sk_proposer, alpha_n)
beacon_n  = H(previous_beacon ‖ output_n)
\end{verbatim}

Three parties would have to cooperate to steer it, and none of them can act
alone:

- \textbf{The proposer cannot choose its output.} A VRF is unique per
  (key, input), so for a fixed \texttt{alpha} there is exactly one proof it can
  produce. Grinding requires re-keying, which changes its address and takes
  it out of the registry.
- **The proposer cannot choose \texttt{alpha}.** Both halves are fixed before it
  acts: the previous beacon is committed state, and the height is the block
  being built.
- \textbf{The miner cannot compute the output.} It does not hold the key.

\subsubsection{What a miner *can* still do, stated plainly}

Include the proposer's proof, or leave it out. Leaving it out yields
[\texttt{fallback\_beacon}] instead, which is a different value — so a miner holding
a block choice gets \textbf{one bit of influence over each block it mines}, free,
and can retry that choice as often as it mines.

That is bounded and it is not zero. Two things reduce what it is worth:

1. The accumulator chains, so a bit chosen now is composed with every future
   proposer's output. Steering a value k blocks out costs influence over all
   k, and the miner holds influence over only the blocks it wins.
2. The fallback is *deterministic*, so the miner chooses between two known
   values rather than searching a space.

The fix, if that bit ever matters for an application, is to move the proof
into the block header so that choosing between the two values costs a full
re-mine rather than nothing. That is a change to [\texttt{crate::core::BlockHeader}]
— and therefore to the proof-of-work preimage, the miner, the pool protocol,
and the GPU kernel — which is why it is named here rather than done.

**Do not use this beacon for a lottery whose payout exceeds a block reward
without reading the paragraph above.**

\subsubsection{Liveness does not depend on the proposer}

An absent or invalid proof is not an error. The chain folds
[\texttt{fallback\_beacon}] and moves on. The alternative — refusing blocks with no
beacon proof — would hand any single authority the power to halt the chain
by going offline, which is a worse property than the one bit above.

\subsection{oracle/feed.rs}
Price feeds: what a quorum agreed, and when.

\subsubsection{Why the median}

A mean is a weighted vote, and one compromised authority holds a weight
bounded only by the number it is willing to sign. A median is not: until a
liar holds a majority of the quorum, moving the median requires moving the
honest values, which is to say it requires being right.

With an even count there is no single middle, and the two candidates are
averaged in most textbook definitions. Not here — averaging introduces a
division, a division introduces rounding, and a rounding rule in a price is
a systematic bias in whichever direction it falls. The lower of the two
middles is taken instead: deterministic, exact, and conservative in the same
direction every time.

\subsubsection{One scale, chain-wide}

Every price is an integer in units of \texttt{1 / FEED\_SCALE}. There is no
per-feed decimals field, for the reason [\texttt{crate::rpc::market}] already gives
about supply: saying "whole units" in one place and "base units" everywhere
else is how a number ends up wrong by a power of ten. A feed that cannot be
expressed at this scale is a feed this oracle does not carry.

\subsubsection{What a signature commits to}

The feed, the round, the value, and the height the authority observed at —
under a domain of its own. All four matter. Without the feed id a signature
for MAYA/USD is a signature for BTC/USD; without the round it replays
forever; without the height a stale observation is indistinguishable from a
fresh one; and without the domain separator it is a transaction signature.

\subsection{oracle/mod.rs}
The oracle subsystem: a randomness beacon and a multi-signed price feed.

\subsubsection{The chain's first trusted party}

Everything else in this node is trustless in the strict sense — proof of
work decides ordering, signatures decide authority over funds, and a
validator who lies is caught by arithmetic anyone can redo. None of that is
available here.

An external price is a fact about the world, and no amount of cryptography
makes a chain able to observe one. So the oracle introduces a named set of
authorities and believes what a quorum of them agrees on. That is a genuine
reduction in the chain's security model, and every design decision below
exists to bound it rather than to disguise it:

- \textbf{A quorum, not a party.} A feed moves only when
  [\texttt{registry::OracleRegistry::quorum}] distinct authorities sign the same
  round.
- \textbf{The median, not the mean.} One compromised authority moves a mean
  without limit. It moves a median by nothing at all until it holds a
  majority of the quorum.
- \textbf{Rotation from the first block.} A set that cannot be changed is a set
  that is permanent, including after a key leaks.
- \textbf{Freshness is not optional.} A contract reading a feed states how stale
  it will tolerate, and gets nothing if the feed is older. See
  [\texttt{crate::state::oracle\_exec}].

\subsubsection{The beacon is a different kind of thing}

Randomness is *not* a fact about the world, so the beacon needs no trust in
what an authority observes — only that it holds a key. A VRF output is
unique per (key, input), so an authority cannot choose its value; it can
only refuse to produce it, and refusal is handled by a fallback that keeps
the chain moving.

What an authority and a miner *together* can still do is bounded and stated
in [\texttt{beacon}].

\subsection{oracle/registry.rs}
Who the chain believes, and how that set changes.

\subsubsection{What an authority is}

An address and a VRF key. Not a full hybrid public key: an address is 32
bytes and a hybrid key is 1984, and every signature that arrives already
carries the key it was made with — that is forced by the address being a
hash of the key rather than the key itself. Storing the key here as well
would be the same bytes in two places, one of which nobody checks.

\subsubsection{Rotation exists from the first block}

Not because a rotation is expected soon, but because the alternative is a
set that cannot change after a key leaks. A registry without rotation is a
registry whose worst day is permanent.

The authority to rotate is the registry itself: a change needs a quorum of
the *current* set, the same threshold that moves a price. There is no
separate admin key, because a separate admin key is a single point of
failure sitting above a construction built to avoid one.

\subsection{rpc/bootstrap.rs}
Serving and fetching what a pruned node bootstraps from, over JSON-RPC.

Server: \texttt{get\_tip\_height}, \texttt{get\_headers}, \texttt{get\_snapshot\_manifest} and
\texttt{get\_snapshot\_chunk}, next to the existing \texttt{get\_block\_by\_height}.
Client: [\texttt{RpcBootstrapSource}], a
[\texttt{crate::state\_pruner::snapshot::BootstrapSource}] over those methods.

JSON-RPC rather than a libp2p protocol for now, as \texttt{backfill\_loop} already
does for history. Nothing here is trusted: the bootstrap validates every
header, checks every chunk against the manifest, and checks the imported
state against the header's root.

\subsection{rpc/limit.rs}
Per-peer RPC rate limiting.

\subsubsection{Why this exists}

Before it, there was no rate limiting anywhere in \texttt{src/rpc} or
\texttt{src/network}. A \texttt{config.toml} field naming a limit over code that did not
enforce one would be worse than no field at all: an operator reads it, sets
it, and believes they are protected.

\subsubsection{What it protects against, and what it does not}

It bounds how fast \textbf{one address} can make this node work. That is the
cheap attack: a single client issuing \texttt{get\_block\_by\_height} in a loop makes
a node read and serialise blocks instead of validating them, and costs the
attacker almost nothing.

It is \textbf{not} a defence against a distributed flood. A thousand addresses at
the limit each are a thousand times the limit, and no per-peer bucket can
see that. That needs a global ceiling or something in front of the node, and
\texttt{maya-api-gateway} is the layer where that belongs.

Stating the boundary matters more than the feature: a limiter described as
"rate limiting" without it invites someone to treat the node's RPC as
publicly exposable, which it is not.

\subsubsection{A token bucket, not a fixed window}

A fixed window admits twice the rate across a boundary — 50 requests at
0.999 s and 50 more at 1.001 s is 100 in two milliseconds, and every
attacker knows it. A bucket refills continuously, so the sustained rate is
the sustained rate and the burst is stated explicitly.

\subsection{rpc/market.rs}
Market-data endpoints for listing aggregators.

\subsubsection{What aggregators actually take from a chain}

CoinGecko's and CoinMarketCap's \texttt{/ticker} formats are their \textbf{exchange}
integration specs: trading pairs, last price, 24-hour volume, order-book
depth. A node has none of that. There is no price until the asset trades
somewhere, and a node that reported one would be inventing it.

What both aggregators consume from a *project* is a \textbf{supply endpoint} — a
plain number for total and circulating supply, used to verify supply claims
and compute market capitalisation. That is real data a node can produce, and
it is what [\texttt{supply}] computes.

This module is the arithmetic and the wire shapes. The HTTP endpoints that
actually serve them — \texttt{/api/v1/total\_supply}, \texttt{/api/v1/circulating\_supply},
\texttt{/api/v1/ticker} — live in [\texttt{crate::rpc::market\_http}], because aggregators
fetch them with a plain \texttt{GET} rather than a JSON-RPC call.

The exchange-shaped ticker is therefore an adapter, not a data source. With
no market feed configured it answers \texttt{503} rather than fabricating a price.
Populating it requires an exchange integration, which is a listing step
rather than a code one.

\subsubsection{Supply under a shielded pool}

Hiding balances does not hide supply. The shielded pool's holdings are
exactly the sum of every \texttt{public\_in} minus every \texttt{public\_out} and fee ever
applied, all of which are in the clear, so the pool tracks its own balance
and total supply stays a matter of arithmetic. Coins inside the pool are
counted as circulating: they are spendable by whoever holds the notes, and
excluding them would understate supply by the amount most worth hiding.

\subsection{rpc/market\_http.rs}
The aggregator-facing HTTP endpoints.

\subsubsection{Why these are not JSON-RPC methods}

\texttt{get\_supply} already exists on the JSON-RPC server, and it is the wrong shape
for the consumers that need it. CoinGecko and CoinMarketCap fetch a supply
figure with an unauthenticated \texttt{GET} and parse the response body as a number.
They do not POST a JSON-RPC envelope, and they will not be taught to. An
endpoint they cannot call is an endpoint that does not exist as far as a
listing is concerned.

So this is a small HTTP surface built on the same hyper server the metrics
exporter uses, serving the handful of paths the aggregators document.

\subsubsection{Why a separate port from the exporter}

These are the only endpoints the node serves that are *meant* to be public
and unauthenticated. The exporter on 9600 is the opposite — it publishes peer
topology and mempool contents, and the NetworkPolicy confines it to the
monitoring namespace. Putting public supply data behind that policy would
mean either exposing the exporter or not answering aggregators, so the two
get separate listeners and the firewall rules can differ.

\subsubsection{Units}

Every figure is in base units, the same unit \texttt{AccountInfo::balance} and the
chain's transfer amounts use. There is no decimals constant in this codebase,
so no scaling is applied. If one is ever introduced, both aggregators must be
told the exponent — an unannounced change of unit shows up as a supply figure
wrong by orders of magnitude, and it is the kind of error that gets a listing
flagged rather than corrected.

\subsubsection{Prices}

A node has no price. There is none until the asset trades somewhere, and a
node that reported one would be inventing it. The ticker endpoints are
adapters over an operator-supplied [\texttt{MarketFeed}]; with no feed configured
they answer \texttt{503} rather than fabricating a quote.

\subsection{rpc/mod.rs}
JSON-RPC interface.

\begin{verbatim}
| Method | Params | Returns |
|---|---|---|
| `get_balance` | `[address_hex]` | [`types::AccountInfo`] |
| `send_raw_transaction` | `[tx_hex]` | [`types::SubmitTransactionResult`] |
| `get_block_by_height` | `[height]` | [`types::BlockInfo`] |
| `get_mining_candidate` | `[]` | [`types::MiningCandidate`] |
| `submit_block` | `[block_hex]` | [`types::SubmitBlockResult`] |
| `get_supply` | `[]` | [`market::SupplyReport`] |
| `threat_indicators` | `[]` | `Vec<`[`types::ThreatIndicatorInfo`]`>` |
| `threat_peer_addresses` | `[]` | `Vec<`[`types::PeerAddressInfo`]`>` — only with [`RpcContext::with_peers`] |
\end{verbatim}

\texttt{submit\_block} is the counterpart to \texttt{get\_mining\_candidate}: a candidate a
miner can fetch but never return is not a usable interface, and the
end-to-end tests exercise the full fetch-mine-submit round trip.

\subsection{rpc/server.rs}
JSON-RPC server.

Handlers are registered as synchronous methods. Every one of them is a short
RocksDB read or an in-memory index lookup — microseconds — so dispatching
them to a blocking pool would cost more in scheduling than it saves. The one
genuine exception is \texttt{submit\_block}, which verifies proof of work: a 32 MiB
Argon2id pass that would stall the reactor, so it runs on a blocking thread.

The chain sits behind a \texttt{Mutex} because fork choice mutates it and the reorg
path must be serialized — two concurrent reorgs would interleave undo
journals and corrupt state.

\subsection{rpc/types.rs}
JSON-RPC wire types.

These are deliberately *not* \texttt{serde} derives on the core consensus structs.
The consensus encoding is a hash preimage: changing a field name or ordering
there changes block and transaction hashes. Keeping the JSON representation
in its own layer lets the API evolve — renamed fields, added conveniences
like \texttt{height} — without touching anything a signature or proof-of-work
commits to.

Byte arrays cross the wire as lowercase hex strings, since JSON has no byte
type and numeric arrays are both larger and harder to read.

\subsection{sealed/mod.rs}
The sealed mempool's chain state: the committee, the pending envelopes, and
the shares that open them.

The cryptography lives in [\texttt{maya\_mev}]. This module is the part that has to
know about storage — how a record is laid out, what key it files under, and
how it enters the Merkle state root.

\subsubsection{One keyspace, one prefix}

Everything here lives under \texttt{m:}, joining the trading subsystem's \texttt{d:}, the
oracle's \texttt{o:}, and governance's \texttt{g:}. No pre-existing prefix begins with
that byte, so one scan collects the whole subsystem — which is what lets it
fold into the state root as a single layer and journal into the undo record
with no new section at all.

\subsubsection{Why keys are big-endian here and nowhere else}

Payload encodings on this chain are little-endian. Storage keys under this
prefix are not, and the reason is that RocksDB orders keys lexicographically
by byte. A big-endian height sorts numerically; a little-endian one sorts
into nonsense, and the end-of-block pass would visit envelopes in an order
that depends on the bit pattern of the height rather than on the height.

That ordering is not a convenience. It is the execution order of every
transaction the committee opens, so it has to be a property of the data and
not of the encoding.

\subsection{state/asset.rs}
Assets other than the native coin, and the balances that hold them.

\subsubsection{Why the native coin is not one of these}

[\texttt{Account}] already holds a balance, every existing transfer moves it, and
the Merkle state root is computed over those records. Giving the native coin
a second home in the multi-asset keyspace would mean two numbers for one
quantity, kept in step by hand.

So it keeps the one it has. [\texttt{NATIVE\_ASSET}] is all zeros, and every access
here routes it back to [\texttt{Account::balance}]. A chain that never registers an
asset has the same state root it always had, and the account record is
unchanged at sixteen bytes.

Nothing derived can collide with the native identifier: it would take a
BLAKE3 output of thirty-two zero bytes.

\subsubsection{Supply is fixed at registration}

An asset is created once, with its whole supply credited to its creator, and
there is no mint operation. That is a deliberate limitation rather than a
stub: a mint authority is an account that can dilute every holder, and
introducing one is a governance decision, not a data-model one. An asset
that needs elastic supply belongs in a contract, where its rules are code
that holders can read.

\subsubsection{No decimals field}

Every amount on this chain is in base units, and [\texttt{crate::rpc::market}]
already refuses to carry a scale factor for the same reason: saying "whole
coins" in one place and "base units" everywhere else is how a listing ends
up wrong by a power of ten. How a wallet chooses to render an asset is the
wallet's business.

\subsection{state/blocks.rs}
The block store: every block this node has accepted, kept in the same
RocksDB as state.

\subsubsection{Why it exists}

Until 2026-09-12 no block was persisted anywhere. \texttt{Chain} held every block
in memory, and RocksDB held only the current state and the undo journals. A
node that had applied one block could not restart. On boot \texttt{seed\_state}
rewrote the genesis allocations over the evolved state, the genesis root
check then failed, and \texttt{Chain::new} rebuilt from genesis alone.

\subsubsection{Why the same database}

A block's canonical-index entry and the tip pointer are written in the
**same \texttt{WriteBatch}** as the state that block produced (see
\texttt{StateDB::apply\_canonical} and \texttt{StateDB::revert\_canonical}). A crash can
therefore never leave state at one block and the tip at another. A second
database would need a two-phase commit to promise that.

\subsubsection{Layout}

\begin{verbatim}
| Key | Value |
|---|---|
| `blk:h:<id>` | header (144 B) ‖ height (u64 LE) ‖ cumulative work (U256 BE) |
| `blk:b:<id>` | the block's wire encoding |
| `blk:n:<height BE>` | id of the active-chain block at that height |
| `blk:meta` | genesis id ‖ tip id ‖ prune horizon (u64 LE) |
\end{verbatim}

Headers stay for every block, forever: 144 bytes each, carrying the
\texttt{state\_root} and \texttt{tx\_root} a pruned body is verified against. Bodies and
undo journals are what \texttt{state\_pruner} removes.

Every key here is under
[\texttt{BLOCK\_STORE\_PREFIX}](crate::state::commitments::BLOCK\_STORE\_PREFIX), which is local-only: never
in the state root and never carried in a snapshot.

\subsection{state/channel.rs}
On-chain payment channel records.

\subsubsection{Lifecycle}

\begin{verbatim}
           OpenChannel
                |
              Open ----- CooperativeClose ----> Closed
                |
          DisputeClose
                |
            Disputed ---- PenaltyClaim -------> Closed  (all funds to victim)
                |
        FinalizeDispute (after the deadline) --> Closed  (disputed balances)
\end{verbatim}

A unilateral close does \textbf{not} pay out immediately. It records the claimed
state and opens a window in which the counterparty can prove that state was
revoked. Paying out at once would make fraud costless: publish a stale state
showing yourself richer, and take the money before anyone can object.

\subsubsection{Why only a commitment is stored}

Each submitted state carries a hash committing to its revocation secret. The
chain stores that one hash. A victim punishes fraud by supplying the
preimage, which they only possess because the closer handed it over when
advancing past that state. Storing the secrets themselves would grow
on-chain state without bound for no gain.

\subsection{state/commitments.rs}
The two state layers that used to sit outside the state root: contracts
and the nullifier set.

Until 2026-09-11 the root folded accounts, channels, trading, oracle,
governance, the sealed mempool and the shielded pool, and nothing else. Two
things RocksDB persisted were missing:

- \textbf{Contract code and storage} (\texttt{code:}, \texttt{cstate:}). Two nodes could
  disagree about a contract's storage and agree on every state root. That
  hid the divergence until it reached a balance.
- \textbf{The nullifier set} (\texttt{null:}). The pool's Poseidon root commits to the
  notes that exist, not to which of them were spent. A state snapshot that
  left a nullifier out would have verified against the root, and the node
  that loaded it would accept a second spend of that note.

Both fold in the way every other layer does: only when non-empty, so a chain
that never deployed a contract and never spent a note keeps the root it
always had.

\subsection{state/context.rs}
Execution context for a block.

\subsection{state/contracts.rs}
Contract deployment and invocation against chain state.

\subsubsection{The host adapter}

\texttt{Store<T>} in wasmtime requires \texttt{T: 'static}, so the VM's host cannot hold a
borrow of the database. [\texttt{ChainHost}] therefore owns everything the call
might touch:

- the contract's \textbf{entire} storage, read up front
- a snapshot of the balances the caller pre-declared

Reading a contract's whole key space before execution is a real cost, and it
bounds how large a contract's storage can usefully get. It buys two things:
the host is \texttt{'static} without any unsafe lifetime games, and reverting is
simply dropping the map — a trapped call cannot leave a partial write
because its writes never touched the database.

\subsubsection{Determinism}

Every method is a pure function of committed state. No clock, no randomness,
no iteration whose order depends on hashing — \texttt{BTreeMap} throughout, so the
order a contract observes is the same on every node.

\subsection{state/db.rs}
Persistent account state backed by RocksDB.

\subsubsection{Atomicity model}

Block execution never writes to RocksDB incrementally. Every mutation lands
first in an in-memory overlay; only after *all* transactions in the block
have passed validation is that overlay flushed as a single
[\texttt{rocksdb::WriteBatch}]. A failure at any transaction returns \texttt{Err} before
the batch is ever handed to the database, so a rejected block leaves state
byte-for-byte unchanged. There is no partial application and nothing to undo.

The overlay is also what makes intra-block double-spend detection work: each
transaction reads the effects of its predecessors in the same block, not the
stale on-disk values.

\subsection{state/dex.rs}
On-chain records for the trading engine: pools, resting orders, and the
keys they live under.

The arithmetic is in [\texttt{maya\_dex}]. This module is the part that has to know
about storage: how a record is laid out, what key it files under, and how it
enters the Merkle state root. Keeping the split there is what lets the
arithmetic be model-checked — see \texttt{crates/dex/Cargo.toml}.

\subsubsection{One keyspace, one prefix}

Everything here lives under \texttt{d:}. No pre-existing prefix (\texttt{acct:}, \texttt{chan:},
\texttt{code:}, \texttt{cstate:}, \texttt{null:}, \texttt{shld:}, \texttt{undo:}) begins with that byte, so one
scan collects the whole trading subsystem — which is what makes the state
root fold and the undo journal a single generic layer rather than five
parallel ones. Five parallel ones is how the three commit paths in
[\texttt{crate::state::db}] drift apart.

\subsubsection{Order keys are the priority queue}

A resting order files under \texttt{d:ord:<pair><side><price><sequence>}, with the
price and sequence big-endian. RocksDB iterates lexicographically, so a
prefix scan yields one side of one book already in price-time order. The
matcher never sorts, and two nodes cannot disagree about the queue because
neither of them chose it.

A second key, \texttt{d:oidx:<order\_id>}, holds the first. Cancelling by identifier
would otherwise be a scan of the whole book.

\subsubsection{Escrow}

A resting order holds its own assets. Placing a bid debits quote from the
owner there and then; the record carries what is left of it, so a
cancellation or an expiry refunds an exact number rather than a recomputed
one. Recomputing it would mean deriving the same rounding twice and being
right both times.

\subsection{state/dex\_exec.rs}
Trading operations against the block-execution overlay.

The arithmetic lives in [\texttt{maya\_dex}]. This is the part that moves balances,
writes records, and decides what happens when the engine says no.

\subsubsection{Two failure classes, and why they are not the same}

A transaction can be \textbf{wrong} — it names a pool that does not exist, its
route's legs do not join up, its asset symbol has a control character in it.
That is an \texttt{Err}, and because a failing transaction fails its whole block,
it is also a statement that no honest miner would have included it.

A transaction can also merely \textbf{lose}. Its slippage bound was missed
because somebody else traded first; its arbitrage was taken by a faster
arbitrageur. That is not an error and must never be one: two traders racing
the same opportunity is the ordinary case, and if the loser's transaction
took the block down with it, every block carrying a competitive trade would
be invalid.

So a lost trade is a \textbf{no-op}. The nonce advances, nothing moves, and the
record of what happened is that nothing did. Getting this wrong is the
difference between a DEX and a denial-of-service surface, which is why it is
stated here and repeated at each site.

\subsubsection{When a swap actually happens}

Not where it sits in the block. \texttt{StateDB::stage\_swap} takes the trader's
input into escrow and adds them to that pool's batch; every batch settles in
\texttt{StateDB::settle\_trading} once the whole block has been staged, at one
price per pool. See [\texttt{maya\_dex::batch}] for what that buys and what it does
not.

Routes are the exception and execute in place — they are the arbitrage path,
and there is nothing to net a multi-hop trade against.

\subsection{state/governance\_exec.rs}
Governance operations against the block-execution overlay.

\subsubsection{Where each thing happens}

Locks, votes, proposals, and work claims execute in place — they are
ordinary transactions with an immediate effect, and one that fails is one
somebody built wrong.

\textbf{Rule changes do not.} Finalizing a closed vote and applying a passed
proposal both happen in \texttt{StateDB::settle\_governance}, once, at the end of
the block, in proposal-identifier order. Two reasons, and the first is
fatal:

1. A parameter change applied where a transaction sits would mean two
   transactions in one block executing under different rules. A node that
   ordered them differently would reach a different state.
2. Finalizing depends on the whole block's votes. A tally read halfway
   through is a tally that is still being written.

\subsubsection{What governance can and cannot reach}

It moves values in a table, and [\texttt{maya\_governance::params}] fixes a hard
range for each. It cannot move the quorum floor, the approval floor, the
minimum voting period, or the minimum timelock — those are compiled in and
reachable by no transaction, because the first move of a hostile majority is
to remove the checks that would have let anyone react.

It cannot reach native code at all. There is no key whose value is a
program.

\subsection{state/htlc.rs}
Lattice HTLC records in the state, and their place in the root.

Invariant 25: one prefix in \texttt{RECORD\_LAYERS}, one
[\texttt{StateLayer}](crate::state::proof::StateLayer), and the undo journal covers
everything written through \texttt{put\_record}. The layer folds only when
non-empty, so a chain with no locks keeps the root it had before HTLC-L.

\begin{verbatim}
| Prefix | Holds |
|---|---|
| `h:lk:<lock id>` | one [`LockRecord`], including the opening once claimed |
\end{verbatim}

One record per lock rather than a record plus a revelation index: a lock
changes exactly once after it is written, and the watcher always knows the
lock id it is waiting on.

\subsection{state/htlc\_exec.rs}
Executing the lattice HTLC transitions: lock, claim, refund.

\subsubsection{A claim or refund that loses is a no-op, never an error}

Invariant 7. A claim and a refund for the same lock race each other at the
expiry boundary as the *ordinary* case: the recipient's watcher claims late,
the sender's watcher refunds early, and whichever is in the earlier block
wins. If the loser were an \texttt{Err} it would void the winner's block — and a
claim carrying a wrong opening would be a way for anyone to void any block.
So every losing outcome returns, the nonce advances, and nothing moves.

A *lock* is different. It is the sender's own transaction with no
counterparty whose block it could void, so a lock that cannot be made is an
error, as issuance is in \texttt{crate::state::rwa\_exec}.

\subsubsection{Neither claims nor refunds are gated by the circuit breaker}

[\texttt{Module::of}](crate::state::Module::of) maps them to no module. A breaker
that halted claims for \texttt{BREAKER\_BLOCKS} while a lock's expiry passed would
let the refund through and hand the swap to the refunder — the safety
mechanism committing the theft. Halting *new locks* is safe, so locks are
gated and settlement is not, which is invariant 28's rule about transfers
applied to value already promised.

\subsubsection{Permissionless}

Anyone may submit a claim or a refund. A claim pays the recipient the lock
named and a refund the sender, so the submitter chooses only *when*, and
only inside the window the timelock already allows. That is what lets a
watcher settle for its owner, and why copying an opening out of the mempool
gets a front-runner nothing but the chance to pay the right party early.

\subsection{state/identity.rs}
Identity records in the state: their prefixes, and their place in the root.

\subsubsection{Everything here is under the state root}

Invariant 25: every persisted consensus record is under the state root, or
it is on an explicit local-only list. These are consensus records — a
verifier decides whether to accept a credential by reading them — so they
get a prefix in [\texttt{RECORD\_LAYERS}](crate::state::commitments) and a
[\texttt{StateLayer}](crate::state::proof::StateLayer) of their own, and the undo
journal covers them for free
because they go through \texttt{put\_record} into \texttt{overlay.records}.

The layer folds only when it is non-empty, like every other. So a chain
where nobody has registered a DID produces exactly the state root it would
have produced before this subsystem existed — the same property invariant 11
gives the oracle, and the reason adding identity to an existing chain is not
a fork.

\subsubsection{Three prefixes, and why they are separate}

\begin{verbatim}
| Prefix | Holds |
|---|---|
| `i:did:<address>` | the subject's [`DidDocument`] |
| `i:att:<issuer><schema>` | an issuer's [`CryptographicAttestation`] |
| `i:rev:<issuer><page>` | one [`RevocationPage`] of that issuer's bitmap |
\end{verbatim}

Separate rather than one record per issuer, because they change at different
rates and for different reasons. A revocation is a single bit flipped by the
issuer, possibly many times a day; an attestation root changes when the
issuer reissues, which is rare. Folding them together would mean rewriting
the root record — and journalling the whole of it — on every revocation.

\subsection{state/identity\_exec.rs}
Executing the identity transitions.

\subsubsection{Every one of these acts on the sender's own DID}

Registration, rotation and revocation take their subject from
\texttt{Transaction::sender}, never from a field. There is no way to express "do
this to somebody else's DID", so there is no authorisation check to get
wrong — the signature that made the transaction valid *is* the
authorisation, and the address it produces is the subject.

That is the same construction the rest of the chain uses and it is worth
being explicit about, because the alternative reads as harmless: a \texttt{subject}
field plus a check that it equals the sender is one refactor away from a
\texttt{subject} field plus a check somebody removed.

\subsubsection{Rotation is the interesting one}

The transaction is signed by the key the document currently names — by
construction, since the address is BLAKE3 over that key pair. So the *old*
key authorises the new one, which is exactly what rotation has to mean. The
epoch must advance by exactly one, which is what stops a replayed rotation
reinstalling a key.

A revoked subject cannot rotate. Otherwise revocation would be a state a
subject could leave, and "revoked" would mean "revoked for now".

\subsubsection{Nothing here writes a claim}

\texttt{AnchorAttestation} writes a root. \texttt{SetRevocationBit} flips a bit. Neither
has a field a personal fact could travel in, and
\texttt{tests/identity\_tests.rs} checks committed state for preimages rather than
trusting that reading.

\subsection{state/invariant\_guard/anomaly.rs}
The anomaly checks: conditions a block may satisfy the rules and still be
wrong about.

[Conservation](super::conservation) refuses a block that creates value.
These two refuse nothing — the block commits — and instead halt the module
that produced the anomaly for [\texttt{BREAKER\_BLOCKS}] blocks. The split is the
whole design: a conservation failure has no innocent reading, and an
anomaly does.

[\texttt{BREAKER\_BLOCKS}]: super::limits::BREAKER\_BLOCKS

\subsection{state/invariant\_guard/breaker.rs}
The circuit breaker: which module is in emergency read-only mode, and
until when.

\subsubsection{Why it is consensus state}

A breaker decides which transactions are valid. If one node halted the DEX
and another did not, the two would disagree about whether a block
containing a swap is valid, and that is a chain split rather than a
difference of local policy. So the breaker is derived from committed state
and the block being executed, it is stored in the state root like every
other consensus record, and nothing node-local ever reaches it — no clock,
no configuration, no flag. That is invariant 28.

\subsubsection{Why it lives under the governance prefix}

\texttt{g:guard:<module>} is under \texttt{GOVERNANCE\_PREFIX}, so it folds into the
\texttt{Governance} state-root layer and the generic \texttt{records} undo journal already
covers it. A new prefix would have needed a new layer, a new entry in
\texttt{RECORD\_LAYERS}, and a new undo section — three places to forget, for a
record that is eighteen bytes. A reorg restores it exactly as it restores a
ballot.

Both of those are \texttt{pub(crate)}, so they are named here rather than linked:
an intra-doc link to a private item renders as plain text on a public build
anyway, and rustdoc warns about it.

\subsubsection{Why an expired record is not deleted}

A breaker clears by its height passing, not by a write. The record stays as
the chain's own account of what tripped and when, at eighteen bytes per
module and at most one record per module, and a later trip overwrites it.
Deleting it would cost a write on the block that clears it and lose the only
on-chain evidence that anything happened.

\subsection{state/invariant\_guard/conservation.rs}
Value conservation: the one equation every block has to satisfy.

\subsubsection{The equation}

For every asset, native and registered, what a block moves must equal what
it issues:

\begin{verbatim}
Σ holdings delta  ==  supply delta
\end{verbatim}

The left side is every place on this chain that can hold value, which is a
closed list:

\begin{verbatim}
| Asset | Held in |
|---|---|
| Native | `acct:` balances, `chan:` capacity, the shielded pool's public balance, `g:lock:` stakes, `g:prop:` deposits, `h:lk:` HTLC escrow |
| Registered | `d:bal:` balances, `d:pool:` reserves, `d:ord:` escrow |
| LP share | `d:bal:` balances |
\end{verbatim}

The right side is zero for the native coin — there is no mint operation and
no block reward, so the genesis allocation is the supply forever. For a
registered asset it is the change in \texttt{d:asset:<id>.total\_supply}, which is
non-zero exactly once, at registration. For a share asset it is the change
in its pool's \texttt{total\_shares}.

\subsubsection{Why a delta and not a sum}

Summing every account each block would make block execution O(state), so
execution would slow down for the life of the chain. The overlay already
holds exactly what the block touched, so every term here is a difference
between an overlay value and its committed one: O(touched), which for a
normal block is a handful of keys.

It is also not weaker. Conservation over deltas, applied to every block from
genesis, is an induction: if the total was right before the block and the
block moved a net zero, the total is right after it. An absolute
\texttt{MAX\_SUPPLY} scan would check the same fact more slowly, and only for the
assets it remembered to look at.

\subsubsection{Why a violation fails the block}

This is not an anomaly threshold, and there is no judgement in it. A block
that creates a unit from nothing is invalid in the same way a block with the
wrong state root is invalid — it is refused, nothing is committed, and no
breaker is written, because there is no state to protect. The
[breaker](super::breaker) exists for the other case: a block that is valid
and suspicious.

\subsection{state/invariant\_guard/limits.rs}
The guard's thresholds, and why each is the number it is.

Compiled in, and reachable by no [\texttt{ParameterKey}]. This is invariant 12's
rule applied to the guard: a governable threshold is one proposal away from
being set to a value that disables the check, and a monitor an attacker can
turn off before the attack is not a monitor.

[\texttt{ParameterKey}]: maya\_governance::ParameterKey

\subsection{state/invariant\_guard/mod.rs}
The invariant guard: what has to be true of a block before it commits, and
what happens when it is not.

\subsubsection{Where it runs}

[\texttt{StateDB::stage\_block}] is the single place a block's overlay is produced.
Every commit path funnels through it — \texttt{apply\_block},
\texttt{apply\_block\_journaled}, \texttt{apply\_block\_checked}, and \texttt{preview\_root} — so one
hook at the end of it covers all four. \texttt{preview\_root} is the one that makes
this more than a check: a miner assembling a candidate through
[\texttt{Chain::candidate\_block}](crate::consensus::chain::Chain::candidate\_block)
runs the same guard, so it refuses to build a block that breaks an invariant
rather than minting one the network then rejects.

\subsubsection{The two kinds of check, and why they differ}

\begin{verbatim}
| Kind | Example | Outcome |
|---|---|---|
| [Conservation](conservation) | a block creates a unit of an asset from nothing | the block is **invalid**; nothing commits |
| [Anomaly](anomaly) | a pool is suddenly worth less per share | the block commits and the module is **halted** for [`BREAKER_BLOCKS`](limits::BREAKER_BLOCKS) blocks |
\end{verbatim}

A conservation failure has no innocent reading, so it is refused exactly as
a wrong state root is refused. An anomaly does have one — a threshold is a
judgement about what is unusual, not about what is legal — so the block
stands and the affected module stops accepting new work while somebody
looks.

\subsubsection{What this guard does not claim to catch}

The brief that produced it names re-entrancy, integer overflow and
flash-loan manipulation. Two of those three attacks have no surface on this
chain and the third is not what the phrase usually means, which
\texttt{tests/exploit\_replays.rs} demonstrates rather than asserts:

- \textbf{Re-entrancy} needs a contract-to-contract call. The VM exposes nine
  host functions and none of them calls a contract or moves value, and
  \texttt{Vm::validate} refuses an unknown import at deploy.
- \textbf{Integer overflow} is already refused: \texttt{ledger-math} is checked
  arithmetic, Kani-verified, and every call site turns \texttt{None} into
  [\texttt{BalanceOverflow}](crate::error::NodeError::BalanceOverflow), which
  fails the whole block.
- \textbf{Flash loans} do not exist. \texttt{SwapRoute} is self-funded — no borrow, no
  callback — and \texttt{l2-flash} is a payment-channel network rather than a
  lending one.

What the guard is actually for is the class underneath all three: value
created from nothing by *any* future bug, across transparent accounts, AMM
reserves, channel escrow, governance locks and the shielded pool. Each of
those five is checked today only by its own local rules, and a bug in one of
them is invisible to the other four.

\subsubsection{Invariant 28}

The guard reads only committed state and the block. No clock, no
configuration, no node-local value, ever. A breaker that trips on one node
and not another changes which transactions are valid on each, which is a
chain split — the same reason invariant 24's state-root check reads only
what execution produced.

[\texttt{StateDB::stage\_block}]: crate::state::db::StateDB

\subsection{state/iot.rs}
IoT anchor records: one per device, under \texttt{v:}, one layer of the state root
(invariant 25).

- \texttt{v:dev:<device id>} — [\texttt{DeviceRecord::encode}].

Each batch overwrites its device's record, so state grows with the number
of devices, never with the number of readings.

\subsection{state/iot\_exec.rs}
Executing the IoT anchor transactions.

\subsubsection{Who bears a failure}

- \textbf{Owner transactions} — enrollment, revocation — are the sender's own, so
  a bad proof, a second enrollment, or a revocation by someone else is an
  error, as an HTLC lock that cannot be made is.
- \textbf{Device messages} are relayed by any gateway. A bad signature is the
  relay's error, but everything that merely *loses* — an unknown device, a
  duplicate relay, a batch overtaken by a newer one, a device already
  terminal — is a no-op, because two gateways racing is ordinary and an
  error would void the block the winner is in (invariant 7).
- \textbf{Clone evidence} conflicting with a recorded batch is detected inside
  \texttt{SubmitTelemetry} itself, so a copied key is caught even if nobody files a
  \texttt{ProveEquivocation}.

Tamper reports, equivocation proofs and revocations belong to no breaker
module: halting them while a key is being abused would protect the abuser.

\subsection{state/merkle.rs}
Binary Merkle tree over the account set.

\subsubsection{Construction}

Leaves are the account records ordered by address, so the root is a pure
function of state content and not of insertion order.

Leaf and internal hashes carry distinct domain tags. Without them an
attacker could present an internal node as if it were a leaf — the classic
Merkle second-preimage attack.

An odd node at any level is \textbf{promoted} to the next level rather than
duplicated. Duplicating it is the flaw behind Bitcoin's CVE-2012-2459, where
two distinct trees hash to the same root.

\subsection{state/mod.rs}
Persistent account state: RocksDB-backed storage, the account state
transition rules, and the Merkle state root.

\subsubsection{Account model}

State transitions are account-based: a transaction's \texttt{public\_key} is the
sender, its \texttt{outputs} are credits, and its \texttt{nonce} is that sender's sequence
number. [\texttt{Transaction::inputs}] belongs to a UTXO model and is carried
through the structure but not consumed here; spending authority comes from
the signature plus the nonce, not from referenced outputs.

[\texttt{Transaction::inputs}]: crate::core::Transaction::inputs

\subsection{state/oracle\_exec.rs}
Oracle operations against the block-execution overlay.

\subsubsection{Where each thing happens}

Feed and registry transactions execute in place: they are not races, nobody
is competing to land first, and there is nothing to net them against. A
submission that fails is a submission somebody built wrong, so it is an
\texttt{Err} — the same reasoning that makes a bad signature an error rather than a
no-op.

The beacon is different. It is *staged* during the block and folded once at
the end, by \texttt{StateDB::settle\_oracle}, because the accumulator must advance
exactly once per block whether or not a proof arrived. Folding it where the
transaction sits would make the beacon a function of how many beacon
transactions a miner chose to include.

\subsubsection{What is checked, and in what order}

For a feed submission, in this order and for a reason:

1. The feed exists. Creating one implicitly would let the first quorum to
   sign a name own it, and a typo'd name would become a second feed that
   looks like the first.
2. The round is newer than the stored one. This is the replay guard: without
   it, a quorum's signatures over an old price stay valid forever.
3. Every signature verifies, is from a registered authority, and no
   authority appears twice. A quorum is a count of *distinct* authorities;
   one key repeated is otherwise a quorum on its own.
4. The count clears the threshold.
5. Only then, the median.

Verification is the expensive step and it is fourth on purpose — the three
cheap rejections come first, so a malformed submission costs a lookup rather
than a quorum's worth of post-quantum signature checks.

\subsection{state/proof.rs}
State proofs: what a light client needs to believe a balance.

\subsubsection{Why an account path is not enough}

The state root is not one tree. It is an accounts tree, folded through a
chain of one-sided layers:

\begin{verbatim}
root = accounts_root
root = H(flash      ‖ root ‖ channels_root)    if any channel exists
root = H(dex        ‖ root ‖ dex_root)         if any `d:` record exists
root = H(oracle     ‖ root ‖ oracle_root)      if any `o:` record exists
root = H(governance ‖ root ‖ governance_root)  if any `g:` record exists
root = H(sealed     ‖ root ‖ sealed_root)      if any `m:` record exists
root = H(shielded   ‖ root ‖ pool_state)       if the pool is not empty
root = H(contracts  ‖ root ‖ contracts_root)   if any code or storage exists
root = H(nullifiers ‖ root ‖ nullifiers_root)  if any nullifier exists
\end{verbatim}

Each layer folds in \textbf{only when it holds something}, which is what keeps a
chain that has never traded, never governed, and never shielded at exactly
the root it would have had before those subsystems existed. The cost is that
the shape of the fold is a fact about the chain's *contents*, not a constant.

So a proof carries the layer digests alongside the path. Without them a
verifier reaches \texttt{accounts\_root} and has nothing to compare it to; with them
it reaches the header's \texttt{state\_root} and has checked everything in between.

\subsubsection{What this does and does not prove}

It proves \textbf{inclusion}: this account, with this balance and nonce, was in
the state committed to by that root.

It does \textbf{not} prove the root is canonical. That is the header chain's job,
and it is why [\texttt{crate::state::proof}] is only half of a light client — the
other half follows headers and picks the chain with the most work. A proof
against a root nobody mined is a proof of nothing.

It also does not prove \textbf{absence}. Showing that an account is *not* in the
tree would need the two neighbouring leaves and an ordering argument, and
nothing in this node needs that yet.

\subsection{state/rwa.rs}
RWA records in the state, and their place in the root.

Invariant 25: one prefix in \texttt{RECORD\_LAYERS}, one
[\texttt{StateLayer}](crate::state::proof::StateLayer), and the
undo journal covers everything written through \texttt{put\_record}. The layer folds
only when non-empty, so a chain with no tokenised assets produces the root
it would have produced before the subsystem existed.

\begin{verbatim}
| Prefix | Holds |
|---|---|
| `r:tok:<asset>` | the [`RwaToken`] |
| `r:cap:<asset><page>` | one [`CapTablePage`] |
| `r:leg:<asset><doc>` | a [`LegalAttestation`] |
| `r:dist:<asset><round>` | a [`RevenueDistribution`], so a round cannot be replayed |
| `r:elig:<asset><holder>` | a cached [`Eligibility`] |
\end{verbatim}

Separate records rather than one per asset, because they change at different
rates: a cap table page moves on every transfer, a token record never after
issuance. Folding them together would journal the token on every trade.

\subsection{state/rwa\_exec.rs}
Executing the RWA transitions: issuance, DvP, gated transfers, revenue.

\subsubsection{A DvP that cannot settle is a no-op, never an error}

Invariant 7, and it is not an analogy — it is the same rule:

> A trade that merely *loses* is a no-op, never an \texttt{Err}. A failing
> transaction fails its whole block here, so making any of these an error
> hands every trader a way to void a block.

Delivery-versus-payment means both legs or neither. The obvious
implementation returns an error when a leg cannot settle — and on this chain
that voids the block, so anyone could kill any block by submitting a DvP
they know will fail. The counterparty does not even have to be involved.

So \texttt{settle\_dvp} — \texttt{pub(crate)}, so named here rather than linked —
applies both legs to the overlay or applies
nothing. The nonce advances, the fee is spent, and the state is exactly what
it was. Atomicity comes from the overlay mutation being all-or-nothing, not
from an error unwinding it.

\subsubsection{Eligibility is read, not proved, on the transfer path}

A disclosure proof is a STARK (ADR-008); the joinsplit, the nearest measured
one, is \textasciitilde{}383 KB. On a block with ten thousand transfers, one per transfer
would be gigabytes every node downloads and verifies, forever, for a check
whose answer changes rarely. So the proof is
verified once in its own transaction and cached in \texttt{r:elg:}; the transfer
reads a record and a height.

The cache expires, in \textbf{block height} rather than time (invariant 9): a
miner may write any \texttt{header.timestamp}, so a freshness rule measured in
seconds would read as safety and provide none.

\subsubsection{Distribution places exactly what it debits}

\texttt{ledger\_math::distribute} is Kani-checked to place exactly the total. That
is not a fairness nicety — the invariant guard refuses any block whose value
deltas do not balance, so a rounding rule that lost a base unit would be a
distribution nobody can mine.

\subsection{state/sealed\_exec.rs}
Executing the sealed mempool: taking envelopes in, taking shares in, and
opening what is due.

\subsubsection{The one inequality everything rests on}

An envelope is included at some height \texttt{h} and opened at \texttt{reveal\_height > h}.
The miner who chose the block at \texttt{h} — who chose *whether* to include it and
in what company — could not read it, because the shares that open it are not
valid until \texttt{reveal\_height}. That is the property. Every other rule in this
file exists to keep that inequality true or to keep its consequences
bounded.

\subsubsection{Where the revealed batch executes, and why it matters}

\texttt{StateDB::settle\_sealed} — \texttt{pub(crate)}, so it is named here rather than
linked — runs after every plaintext transaction in the block has been staged
and *before* [\texttt{crate::state::db::StateDB}]'s trading pass. Both halves of
that sentence are load-bearing.

After the plaintext transactions: a miner who assembles the reveal block
learns the plaintexts while building it, and could place a transaction of
their own in the same block. Running the batch last means any such
transaction executes *before* the revealed one, so the miner can back-run
but not front-run. Back-running is ordinary arbitrage; front-running is the
attack.

Before the trading pass: a revealed swap therefore stages into the same
uniform-price batch as every plaintext swap in the block. A miner who saw a
revealed swap coming and raced it settles at the identical clearing price —
which is what turns the residual advantage above from a profit into a
nuisance. See \texttt{docs/dex.md}.

\subsubsection{A revealed action that fails is a no-op, never an error}

This is [\texttt{crate::state}]'s trading rule applied to a strictly worse case. A
failing transaction fails its whole block here, so if a revealed action
could return \texttt{Err}, anyone could seal a payload that is guaranteed to fail
and thereby void the block that opens it — a censorship weapon costing one
transaction fee, aimed at a block chosen days in advance.

So a revealed action that cannot execute simply does not: the envelope is
consumed, the sender's nonce was spent when they submitted it, and nothing
else moves. The sender loses what a failed transaction always loses.

\subsection{state/settlement.rs}
On-chain settlement of Maya Flash channels.

Implements the six channel operations against the block-execution overlay,
so channel writes are as atomic as account writes: a block that fails at any
transaction leaves neither balances nor channel records touched.

\subsubsection{The rule every operation shares}

A channel may never pay out more or less than it escrowed. Every closure is
checked against the recorded capacity before a single unit moves. Without
that check a channel would be a mint: two parties could sign any balances
they liked and the chain would honour them.

\subsubsection{Why a unilateral close does not pay out}

[\texttt{TxKind::DisputeClose}] only *records* the claimed state and starts a
timer. Paying out immediately would make fraud free — publish a stale state
showing yourself richer and walk away before anyone can object. The window
exists so the counterparty can present the revocation secret and take
everything instead.

\subsection{state/shielded.rs}
The shielded pool: commitment tree, anchor window, and proof verification.

\subsubsection{What the pool stores}

Two things, and neither reveals anything about who owns what. The
commitment tree accumulates note commitments — opaque hashes — and never
removes one. The nullifier set records which notes have been spent, without
saying which commitment each nullifier retires.

Because notes are never removed, a Merkle path stays valid as the tree
grows only if the verifier accepts the root the path was built against. That
is what the anchor window is for: a spend proves membership under some root
from the recent past, not necessarily the current one. Without it a wallet
would have to build its proof and get it mined inside a single block, and
essentially every shielded transaction would fail.

\subsubsection{Cost, and why it is capped}

A joinsplit proof is a Plonky3 STARK (ADR-008): about 0.4 MB on the wire
and a verification that hashes its way through FRI queries. The cap on
joinsplits per block is set by size first — sixteen proofs are \textasciitilde{}6 MB of an
8 MiB gossip frame — and bounds verification work for the same reason
[\texttt{MAX\_BATCH\_CLOSURES}](crate::core::payload::MAX\_BATCH\_CLOSURES) does.

\subsubsection{Post-quantum, with no setup}

Until 2026-09-21 this pool verified Groth16 over BLS12-381: pairing-based,
so a quantum adversary could forge a proof and mint hidden supply, and its
parameters came from a setup nobody had run for real. The STARK's
soundness rests only on hash collision resistance (Poseidon2 inside,
Keccak outside) and it has no setup at all — \texttt{maya\_zk\_stark::pool}.

\subsection{state/stateless.rs}
\begin{framed}\noindent\textbf{Status.} This module is a research branch. It is present in the tree and no consensus path reaches it. Nothing described here is part of the running chain.\end{framed}

Stateless verification of transfer blocks. Research branch.

\subsubsection{The switch}

From [\texttt{STATELESS\_ACTIVATION\_HEIGHT}] the block being staged writes one
record, \texttt{sl:accounts-sparse}, and from then on the accounts root under the
state root is \texttt{maya-stateless-core}'s sparse tree instead of
\texttt{state::merkle}'s address-ordered one. The marker sits under its own layer,
so the switch is itself committed, journalled and reverted like any record,
and \texttt{state\_root()} can tell which tree to build without being told a
height. The activation block's *pre*-state is still dense, so the first
block a stateless node can check is the one after it.

\subsubsection{What a witness can decide}

A block whose transactions are all \texttt{TxKind::Transfer}, on a state with no
\texttt{m:}, \texttt{o:} or \texttt{g:} records. That is not caution but the list of everything
\texttt{StateDB::stage\_block} does besides staging transactions:

\begin{verbatim}
| Pass | Reads | Inert when |
|---|---|---|
| `settle_sealed` | envelopes due at this height (`m:`) | the sealed layer is absent |
| `settle_trading` | pairs this block's DEX transactions touched | no transaction is a DEX kind |
| `settle_oracle` | the authority registry (`o:`) | the oracle layer is absent |
| `settle_governance` | every proposal (`g:`) | the governance layer is absent |
| `check_anomalies` | pool records and the shielded pool the block wrote | no transaction is a DEX or shielded kind |
\end{verbatim}

Conservation holds by construction for transfers — a debit of exactly the
output total — and a transfer maps to no breaker module. Anything else is
[\texttt{StatelessError::Unverifiable}]: a stateless node cannot say, which is not
the same as saying no.

\subsubsection{Invalid is a verdict; unverifiable is not}

A witness is not under the block id or the proof of work. A relay can
corrupt it, and a body that fails its header's \texttt{tx\_root} may be a relay's
substitution under an honest header — invariant 24's censorship case. So
both are \texttt{Unverifiable}, and the caller fetches again. \texttt{Invalid} is reserved
for what the block's own contents decide once the witness has matched the
parent root: a bad signature, a broken transfer rule, or a declared state
root the transfers do not produce.

\subsection{state/threat.rs}
Threat-intel records: one indicator per author, one marker per piece of
evidence.

Every key is under \texttt{t:}, one layer of the state root (invariant 25):

- \texttt{t:in:<author>} — [\texttt{ThreatIndicator::encode}];
- \texttt{t:ev:<evidence id>} — the height that recorded it, little-endian.

The marker is what makes a second copy of the same evidence a no-op. Its id
covers the offence kind and the signed bytes but not the signature, so an
author who re-signs one message has produced the same evidence, not more.

\subsection{state/threat\_exec.rs}
Executing an attack attestation.

\subsubsection{What is checked, in order}

1. The author's gossipsub signature over the rebuilt signed bytes, with
   \texttt{verify\_strict}. Strict, because consensus needs one answer: a signature
   libp2p's looser check accepted and this one refuses is evidence nobody can
   attest — evasion by the author, never a way to frame someone else.
2. That the frame decodes and then fails the check the kind names. A frame
   that verifies is a lie about an honest author and is refused.

Both are stateless, so a miner learns whether an attestation is valid
before including it.

\subsubsection{Invalid evidence is an error; repeated evidence is a no-op}

The attestation is its submitter's own transaction with nobody else's
outcome riding on it, so evidence that does not verify fails the block —
as an HTLC lock that cannot be made does (\texttt{crate::state::htlc\_exec}). A
*second copy* of evidence already on chain is different: two honest
observers of the same offence submit it independently as the ordinary case,
and if the later one were an error it would void the block the earlier one
sits in. So it returns, the nonce advances, and nothing is written —
invariant 7's rule applied to a race between reporters.

\subsection{state/undo.rs}
Undo journal for reverting a committed block.

A reorg has to move state backwards, but a [\texttt{rocksdb::WriteBatch}] only
moves it forwards. So each journaled block records the *prior* value of
every account it touched, which is enough to reconstruct the pre-block state
exactly.

An account that did not exist before the block is recorded with
\texttt{existed = false} and deleted on revert, rather than being restored as a
zero-balance record. The distinction matters: a spurious zero account would
change the Merkle state root and make the reverted state disagree with the
chain it is supposed to match.

\subsection{state/vm\_exec.rs}
Contract deployment and invocation as state transitions.

\subsubsection{Reverting}

A contract call runs entirely against an owned [\texttt{ChainHost}]. Its writes
reach the block overlay \textbf{only} when execution succeeds. A call that traps
or exhausts gas therefore leaves nothing behind: there is no partial write to
undo because no write ever left the host.

\subsubsection{Determinism}

The VM is configured deterministically and its wasmtime version is pinned;
see \texttt{maya\_vm::config}. What this layer must not do is introduce
non-determinism of its own, which is why the host is populated from committed
state through ordered maps and nothing else.

\subsection{state/zkml.rs}
The node's answer to \texttt{host\_verify\_zkml\_proof}: there is no verifier.

The halo2/KZG-over-BN254 verifier this module used to call was removed on
2026-09-21 (ADR-008): a pairing-based proof is forgeable by a quantum
adversary, and its SRS was derived from a public seed, so anyone could
forge one today (invariant 22). zkML becomes PLANNED until it is
re-expressed as a transparent STARK.

The host function stays in the VM's ABI — removing an import a deployed
module links against is a harder break than answering it — and it was
dark anyway: \texttt{ZKML\_ACTIVATION\_HEIGHT} is \texttt{u64::MAX}, so
\texttt{ContractHost::verify\_zkml} returns before reaching [\texttt{verdict}]. If it is
ever reached, the answer is \texttt{Malformed}, which stops the call rather than
guessing at a boolean.

\subsection{state\_pruner/archive.rs}
The uploader: archive a batch, read every copy back, and only then prune.

\subsection{state\_pruner/cold.rs}
Fetching a pruned block back, from an archive nobody has to trust.

Three checks, each against something this node already holds:

1. every archive section hashes to its CID ([\texttt{maya\_archive::open\_archive}]);
2. the archive's root is the one the receipt recorded when this node wrote
   it and read it back;
3. the decoded block's id is the id of the \textbf{header this node kept}, and
   its transactions match that header's \texttt{tx\_root}.

The third is what makes an untrusted archive safe. The kept header's proof
of work was checked when it arrived, and the header commits to the
transactions. Before the header carried \texttt{tx\_root}, a body fetched from
anywhere could not have been tied to anything.

\subsection{state\_pruner/mod.rs}
Historical pruning: dropping old block bodies and undo journals, after a
verified copy of them exists somewhere else.

\subsubsection{What "final" means here}

This chain is proof of work. Nothing is ever final, so this module does not
pretend it is. A \textbf{pruned node} instead adopts a local policy: a block at
least [\texttt{PRUNE\_DEPTH}] below the tip on the active chain will not be
reorganised away, and such a node refuses any reorg reaching below its
horizon ([\texttt{crate::error::NodeError::BelowPruneHorizon}]). An archive node
prunes nothing, has no horizon, and is unaffected. Pruning is opt-in and
off by default.

[\texttt{PRUNE\_DEPTH}] is one DAG epoch, 30,000 blocks, about 5.2 days at 15-second
blocks. That is far deeper than any honest reorg on a working network.

\subsubsection{What is kept}

Every header, forever: 144 bytes each, about 4.3 MB per 30,000 blocks.
A header carries \texttt{tx\_root} and \texttt{state\_root}, which is all a pruned body is
ever verified against. Those are the "light Merkle roots" left where the
blocks were. Bodies and undo journals below the horizon are deleted.

\subsubsection{Ordering: archive, verify, then delete}

[\texttt{archive::Archiver::archive}] writes the batch to every configured store
and \textbf{reads each copy back and verifies it} before returning a receipt.
[\texttt{crate::consensus::Chain::prune}] is only ever called with that receipt.
So a body is never deleted before a checked copy of it exists. Pruning
without any archive ([\texttt{ArchivePolicy::None}]) is possible, as it is for a
Bitcoin pruned node, but it has to be chosen explicitly.

\subsection{state\_pruner/snapshot.rs}
State snapshots, and bootstrapping a pruned node from one.

\subsubsection{What a new pruned node downloads}

1. \textbf{Every header}, genesis to tip. 144 bytes each, fully validated: the
   exact retarget rule, proof of work, cumulative work.
2. **The state at a height \texttt{H} at least K below the tip**, as a snapshot:
   every key/value under a committed prefix, in chunks.
3. **The bodies above \texttt{H}**, applied with full validation.

Nothing below \texttt{H} but headers. That is the whole saving.

\subsubsection{Why the snapshot can be believed}

Each chunk is checked against the manifest's chunk hashes. That only proves
the server was consistent. The snapshot is then imported, the state root
recomputed, and the result required to equal \texttt{header(H).state\_root}. That
header's proof of work was checked in step 1, and its state root is checked
by every full node since the tx\_root fix. A server that lies about one
balance, drops one nullifier or changes one contract slot produces a
different root, and the import is refused and wiped. The root covers every
committed prefix (see \texttt{state::commitments}), and a key under any other
prefix is refused on sight.

What this cannot rule out is the SPV limitation: a source that shows only
a lower-work fork of headers. That is the same eclipse problem a light
client has, and the answer is the same: ask more than one peer.

\subsubsection{Snapshots at rest}

[\texttt{Snapshots::take}] makes a RocksDB checkpoint (hard links, cheap) at the
current tip. [\texttt{Snapshots::serveable}] serves the newest one at least K deep,
so what a new node adopts is below the depth a pruned node treats as final.

\section{Ledger arithmetic}
\noindent\texttt{crates/ledger-math\textbackslash{}src}

\subsection{lib.rs}
Ledger value arithmetic for \texttt{custom-l1-node}, isolated so it can be proved.

\subsubsection{Why this is a crate and not a module}

Every operation here is three lines of \texttt{checked\_add} or \texttt{checked\_sub}. That
is not what justifies the crate boundary — verifiability is.

The [Kani Rust Verifier] compiles a crate *and its whole dependency graph*
into a CBMC goto-binary. The node depends on \texttt{rocksdb} (C++ through
\texttt{librocksdb-sys}), \texttt{libp2p}, \texttt{fips204}, and the \texttt{ark-*} SNARK stack. Kani
has no semantics for the C++ and no appetite for the rest, so
\texttt{cargo kani -p custom-l1-node} is not a slow proposition, it is not a
possible one.

A dependency-free leaf crate is. The node calls into this crate rather than
inlining the arithmetic, so what the proofs in \texttt{proofs} (compiled only under
Kani) establish is a property of the code that actually executes when a
block is applied — not of a restatement of it that could drift.

\subsubsection{What is proved, and what is not}

Every scalar function here has an *unbounded* proof: the harness quantifies
over all \texttt{u64} inputs, and CBMC discharges it without a loop bound.

[\texttt{total\_outputs}] is the exception. It folds over a sequence, so its harness
carries a \texttt{\#[kani::unwind]} bound and an accompanying \texttt{kani::assume} on the
length. That proof is bounded — it holds for sequences up to the bound, not
for all sequences. It is stated that way in \texttt{proofs} (compiled only under
Kani) rather than left for a reader to infer.

[\texttt{distribute}] is only partly proved: Kani covers its placement of the
spare, and its flooring is checked by exhaustive enumeration over a small
domain. The split, and why, is at the top of \texttt{proofs}.

\subsubsection{What these functions deliberately do not do}

They return [\texttt{Option}], not the node's \texttt{NodeError}. Keeping \texttt{custom-l1-node}
out of this crate's dependencies is the whole point, and an error type is
not arithmetic: the caller owns the mapping from "this would have wrapped"
to whichever \texttt{NodeError} variant its context calls for.

[Kani Rust Verifier]: https://model-checking.github.io/kani/

\subsection{proofs.rs}
Kani harnesses for every operation in this crate.

Run them with:

\begin{verbatim}
cargo kani -p maya-ledger-math
\end{verbatim}

\subsubsection{How these are written}

Most harnesses check the operation against a \textbf{wider-type oracle}: the same
arithmetic performed in \texttt{u128}, where no \texttt{u64} ledger quantity can overflow.
That is a stronger statement than "the result did not wrap", because it pins
down what the answer must be rather than only what it must not be. A
\texttt{checked\_add} replaced with a \texttt{saturating\_add} still never wraps; it fails
the oracle immediately.

\subsubsection{Bounded versus unbounded}

Stated plainly, because the difference is the difference between a proof and
a very good test:

- [\texttt{credit}], [\texttt{debit}], [\texttt{advance\_nonce}], [\texttt{combined\_balance}] and
  [\texttt{settle\_pool}] are proved \textbf{unbounded}. Their harnesses quantify over
  every \texttt{u64} input, with no loop and no unwind bound. There is no input
  these functions have not been checked on.
- [\texttt{total\_outputs}] folds over a sequence, so its harness carries a
  \texttt{\#[kani::unwind]} bound and is proved for sequences up to that length
  only. It is a \textbf{bounded} proof. The bound is set at
  [\texttt{PROVED\_OUTPUT\_COUNT}] below, which is smaller than the 65,536 outputs the
  wire codec will accept, and that gap is real rather than an oversight —
  see the constant's documentation.
- [\texttt{distribute}]'s placement of the spare is proved here, \textbf{bounded} at
  three holders but at every magnitude. Its flooring and sort are not: a
  harness through them has to decide a symbolic 128-bit division per holder
  and did not finish in fifteen minutes. Those are checked exhaustively over
  a small domain by \texttt{every\_small\_distribution\_is\_exact\_and\_fair} in the
  crate's unit tests instead, which is a finite enumeration, not a proof.

[\texttt{credit}]: super::credit
[\texttt{debit}]: super::debit
[\texttt{advance\_nonce}]: super::advance\_nonce
[\texttt{combined\_balance}]: super::combined\_balance
[\texttt{settle\_pool}]: super::settle\_pool
[\texttt{total\_outputs}]: super::total\_outputs
[\texttt{distribute}]: super::distribute

\section{Post-quantum primitives}
\noindent\texttt{crates/crypto-pq\textbackslash{}src}

\subsection{agility/audit.rs}
The security-level audit.

Run against a policy, it reports every suite that is below
[\texttt{MIN\_PQ\_SECURITY\_BITS}](super::MIN\_PQ\_SECURITY\_BITS) and still able to
sign — on testnet and mainnet [\texttt{SuitePolicy::permitted}] already refuses
those, so in practice this fires on devnet, which is where Ed25519 lives —
and every deprecated suite that accounts still need to leave. Findings are
data for an operator or a governance proposal; the audit changes nothing.

\subsection{agility/migration.rs}
Moving accounts off a deprecated suite.

\subsubsection{The schedule}

1. \textbf{Deprecation} (\texttt{SuitePolicy::with\_deprecation}). The suite still
   signs. An account can [\texttt{rotate}](MigrationState::rotate) to a key in an
   active suite, authorized by its current key, or
   [\texttt{commit\_recovery}](MigrationState::commit\_recovery) — register the hash
   of an upgraded public key without switching yet.
2. \textbf{Sunset.} The suite stops signing. [\texttt{sweep\_step}](MigrationState::sweep\_step)
   walks the accounts a bounded number per block and moves every balance
   still held under a sunset key into the vault.
3. \textbf{Vault.} A vault entry opens only to the key whose commitment the
   account registered while its old key was still trusted, and only with a
   signature by that new key in a suite that is active. The old key cannot
   open it — after sunset, it is assumed broken.

An account that registered nothing stays locked. Proving ownership after
the fact needs a proof of knowledge of the wallet seed, which is a STARK
over the key derivation: PLANNED, not built, and said so here rather than
papered over with a path that trusts the broken key.

\subsubsection{No downtime}

The sweep is incremental: each call examines at most \texttt{budget} accounts,
resuming from a cursor, so a block's migration work is bounded and every
other account keeps transacting while it runs. \texttt{tests/agility\_migration\_tests.rs}
runs that with 1,000,000 accounts.

\subsection{agility/mod.rs}
The crypto-agility engine — ADR-007.

Three parts, each a pure function of its inputs so the node can run it
inside consensus without a clock or a random number:

- [\texttt{SuitePolicy}] — the governance-chosen default suite and every suite's
  deprecation schedule, by block height. Changing it returns a new policy;
  nothing is mutated in place, so a rejected proposal leaves no trace.
- [\texttt{audit}] — flags every live suite whose claimed post-quantum security
  is below [\texttt{MIN\_PQ\_SECURITY\_BITS}], and every deprecated suite accounts
  still need to leave.
- [\texttt{migration}] — moves accounts off a deprecated suite: a window in which
  the old key can register an upgraded one, then a sweep of whatever is
  left into a vault only an upgraded key can open.

Heights are \texttt{u64}; nothing here reads wall-clock time (Standing Order 4).

\subsection{envelope.rs}
The suite-tagged signature envelope — ADR-007.

\begin{verbatim}
version:u8 ‖ suite:u8 ‖ pk_len:u32le ‖ pk ‖ sig_len:u32le ‖ sig
\end{verbatim}

Variable-length on the wire, exact-length per suite: both lengths are
checked against the registry row *before* anything is allocated for them,
so a hostile length prefix costs the decoder nothing.

A pre-envelope hybrid signature (the node's v1 transaction encoding) is a
valid suite \texttt{0x30} envelope once the two-byte header and the lengths are
prepended — [\texttt{SignedEnvelope::from\_legacy\_hybrid}] — so history keeps
verifying after the switch.

\subsection{hqc.rs}
HQC-192 (draft FIPS 207) key encapsulation — the second KEM.

A concrete, non-generic wrapper over \texttt{hqc-kem}, shaped deliberately like
[\texttt{crate::kem}] so the two can be combined without either side of the
handshake special-casing which is which.

\subsubsection{Draft, not standard}

NIST \textbf{selected} HQC in March 2025 as the backup key-encapsulation
mechanism to ML-KEM. The specification is \textbf{draft FIPS 207}; it is not a
published standard, and the parameters and encodings can still move.

Every name in this module says draft, and the wire protocol carries a
version byte. That is the same discipline \texttt{crate::kem} applied when it
refused to ship round-3 Kyber under a FIPS 203 label: a handshake whose
advertised protocol name misstates its primitive is a trap for whoever
audits it next, and "draft standard shipped as standard" is that trap with
a longer fuse.

\subsubsection{Why a second KEM at all}

ML-KEM rests on Module-LWE — a structured lattice assumption. HQC rests on
the hardness of decoding random quasi-cyclic codes. The two are unrelated,
so an advance against one is not an advance against the other. That is the
same reasoning that put ML-DSA and SLH-DSA on every transaction (see
\texttt{docs/hybrid-signatures.md}), applied to the transport.

\subsubsection{Why HQC-192 and not HQC-128}

The transport's other KEM is ML-KEM-768, which is NIST level 3. Pairing it
with a level 1 code-based KEM would mean the combined channel is only as
strong as level 1 against a break in ML-KEM — which is the exact scenario
the second KEM exists for. Matching the level is the only choice that makes
the pairing mean what it says.

The cost is bytes, and it is not small. See [\texttt{HANDSHAKE\_OVERHEAD}].

\subsubsection{Sizes and cost}

\begin{verbatim}
| | ML-KEM-768 | HQC-192 |
|---|---|---|
| encapsulation key | 1,184 B | **4,514 B** |
| ciphertext | 1,088 B | **8,978 B** |
| shared secret | 32 B | 32 B |
\end{verbatim}

HQC-192 adds 13,492 bytes to a handshake that was 2,272 — roughly seven
times the bytes. \texttt{network::pq::measure} reports the segment counts and the
wall clock; the headline is that the largest single message is four segments
against a ten-segment initial congestion window, so the extra bytes cost
bandwidth but \textbf{not} a round trip. See \texttt{docs/pq-transport.md}.

\subsubsection{What this does not provide}

Confidentiality only, exactly as [\texttt{crate::kem}]. Peer authentication stays
classical ed25519 in the Noise layer below, and nothing here changes that.

\subsection{kem.rs}
ML-KEM-768 (FIPS 203) key encapsulation.

A concrete, non-generic wrapper over \texttt{ml-kem}. Used by the node's P2P
transport to layer post-quantum confidentiality over its Noise session —
see \texttt{custom\_l1\_node::network::pq}.

\subsubsection{What a KEM is doing in a transport}

libp2p authenticates and encrypts connections with Noise XX over X25519.
That is sound today and gives perfect forward secrecy per connection. It is
not sound against an adversary who records traffic now and decrypts it once
a cryptographically relevant quantum computer exists — "harvest now, decrypt
later". A ledger is an unusually attractive target for that: gossip traffic
reveals which peer originated which transaction, and that metadata does not
expire.

ML-KEM-768 closes it. A fresh keypair per connection, encapsulated inside
the already-authenticated Noise channel, yields a shared secret an adversary
must break Module-LWE to recover — *in addition to* breaking X25519 to get
at the Noise layer underneath. Recording the traffic buys nothing.

\subsubsection{Why `ml-kem` and not `pqcrypto-kyber`}

\texttt{pqcrypto-kyber} implements \textbf{round-3 Kyber}, not FIPS 203. Its last
release predates the standard, and the two are not interchangeable: FIPS 203
changed the shared-secret derivation (the ciphertext hash was dropped from
the KDF), changed encodings, and added the seed-based key format. Shipping
it under a FIPS 203 label would be false.

\texttt{ml-kem} is also pure Rust, which the infra/docker/Dockerfile's static-musl build depends
on — the same reason \texttt{fips204} and \texttt{slh-dsa} were chosen over their PQClean
equivalents.

\subsubsection{Sizes and cost}

\begin{verbatim}
| | ML-KEM-768 |
|---|---|
| encapsulation key | 1184 B |
| ciphertext | 1088 B |
| shared secret | 32 B |
| keygen | 52 µs |
| encapsulate | 42 µs |
| decapsulate | 51 µs |
\end{verbatim}

A full connection setup is \textasciitilde{}144 µs and \textasciitilde{}2.3 KB on the wire. Against a 15 s
target block time that is 0.001 \% of one block interval, and one core
sustains roughly 6,900 handshakes per second — orders of magnitude above
what a gossip mesh of tens of peers asks for.

\subsubsection{What this does not provide}

Confidentiality only. The peer on the other end is authenticated by the
Noise layer below, using classical ed25519, and nothing here changes that.
An adversary with a quantum computer *present at connection time* could
still impersonate a peer. That is an active attack requiring the machine to
exist and be on the wire, not the recording attack this closes.

\subsection{kem\_suite/dual.rs}
The dual KEM: ML-KEM and HQC together — ADR-009.

\begin{verbatim}
ss = SHA3-256( "maya2c.dualkem.v1" ‖ id
             ‖ ss_mlkem ‖ ss_hqc
             ‖ ct_mlkem ‖ ct_hqc
             ‖ ek_mlkem ‖ ek_hqc )
\end{verbatim}

Lattices and codes fail independently, so an adversary who breaks one
still faces the other: the output is a PRF of both secrets. Both
ciphertexts are bound (the brief), and so are both encapsulation keys — the
X-Wing argument, which stops a ciphertext for one key being replayed
against another when a component KEM is not ciphertext-binding.

Wire format: \texttt{ek = ek\_mlkem ‖ ek\_hqc}, \texttt{ct = ct\_mlkem ‖ ct\_hqc}.

\subsection{kem\_suite/hqc.rs}
HQC-128 and HQC-256 — \textbf{draft standard}.

NIST selected HQC in March 2025 as the code-based backup to ML-KEM; FIPS
207 is not final. The backend is pinned to a release candidate that
matches the reference implementation's KATs (\texttt{tests/kem\_kat\_tests.rs}),
and every name here says *draft* so nobody reads it as a finished standard.

\textbf{Not constant-time.} \texttt{benches/dudect.rs} measured decapsulation of a
valid versus a random ciphertext at |t| = 26.9 (null control 2.5): a
remote attacker who can time decapsulation of chosen ciphertexts has a
side channel. ADR-009's addendum records it; the \texttt{hqc} feature that gates
this module is off by default and nothing on a transport path enables it.

Randomness is drawn from the OS through \texttt{getrandom} and fed to the
backend's deterministic entry points, so an entropy failure is an error
the caller sees rather than a panic inside the backend's RNG adapter.

\subsection{kem\_suite/ml\_kem.rs}
ML-KEM-768 and ML-KEM-1024 (FIPS 203).

Randomness comes from \texttt{getrandom} into FIPS 203's seeded entry points
(\texttt{KeyGen\_internal(d, z)}, \texttt{Encaps\_internal(m)}), so an entropy failure is
an error rather than a panic in the backend's RNG adapter — and the same
entry points are the ones the ACVP vectors exercise.

\subsection{kem\_suite/mod.rs}
KEM suites — ADR-009.

ML-KEM-768/1024 (FIPS 203, final), HQC-128/256 (\textbf{draft}: NIST's 2025
selection, FIPS 207 not yet published), X-Wing (ML-KEM-768 + X25519, per
draft-connolly-cfrg-xwing-kem), and two dual-KEM combiners that pair an
ML-KEM with an HQC so that breaking one family is not enough.

This is the primitive layer; wiring any of it into the transport is Master
Prompt 7. The pre-existing [\texttt{crate::kem}] (ML-KEM-768) and [\texttt{crate::hqc}]
(HQC-192) modules keep serving the current handshake unchanged.

Every suite takes and returns byte slices with exact per-suite lengths,
checked on the way in: the bytes come from a peer.

\subsection{kem\_suite/xwing.rs}
X-Wing: ML-KEM-768 + X25519 (draft-connolly-cfrg-xwing-kem).

The reason to carry a classical half at all: ML-KEM is new code. A bug in
a young implementation is a likelier failure today than a break of
X25519, and X-Wing's combiner is secure if *either* half is. The
construction, key expansion and combiner are the draft's, via the \texttt{RustCrypto}
crate, and \texttt{tests/kem\_kat\_tests.rs} runs the draft's own vectors.

\subsection{lib.rs}
Post-quantum primitives, isolated behind concrete APIs.

Two modules, one for each thing the chain needs:

- [\texttt{sig}] — SLH-DSA-SHA2-128s (FIPS 205), the hash-based half of every
  transaction signature. Its lattice partner, ML-DSA-65, lives in
  \texttt{custom\_l1\_node::crypto::keys}.
- [\texttt{kem}] — ML-KEM-768 (FIPS 203), the key encapsulation the P2P transport
  layers over its Noise session.
- [\texttt{suite}] — the signature-suite registry (ADR-007): Ed25519, ML-DSA-65/87,
  SLH-DSA-SHA2-128s / SHAKE-256f, and the ML-DSA + SLH-DSA hybrid, each
  behind one trait and a one-byte id.
- [\texttt{envelope}] — the suite-tagged signature encoding.
- [\texttt{kem\_suite}] — ML-KEM-768/1024, HQC-128/256 (draft), X-Wing, and the
  ML-KEM + HQC dual-KEM combiners (ADR-009).
- [\texttt{agility}] — default-suite policy, the security audit, and migration
  off a deprecated suite.

\subsubsection{Why this crate exists at all}

Both backends are generic over their parameter set, so their hot code
monomorphizes into whichever crate *instantiates* it. A
\texttt{[profile.dev.package.slh-dsa]} or \texttt{[profile.dev.package.ml-kem]} override —
the trick that keeps \texttt{fips204} usable in dev builds — therefore does
nothing: it optimizes a crate that contains none of the work. Both were
measured:

\begin{verbatim}
| | dependency override only | instantiating crate optimized |
|---|---|---|
| SLH-DSA-SHA2-128s sign | 4468 ms | 140 ms |
| ML-KEM-768 handshake | 5805 µs | 785 µs |
\end{verbatim}

So every instantiation is confined here, behind APIs that name concrete byte
arrays rather than type parameters, and the root \texttt{Cargo.toml} optimizes
*this* crate in dev. That keeps the node's own code unoptimized and
debuggable, which is the property the root profile overrides exist to
preserve — raising \texttt{[profile.dev]} workspace-wide would sacrifice it.

Keeping the generics in one place has a second benefit: a caller cannot
reach for the wrong parameter set by accident, because no other parameter
set is nameable from outside.

\subsection{sig.rs}
SLH-DSA-SHA2-128s (FIPS 205) key handling.

A concrete, non-generic wrapper over \texttt{slh-dsa}. The node uses this beside
ML-DSA-65 (FIPS 204), never instead of it — see
\texttt{custom\_l1\_node::crypto::hybrid} for the construction that binds the two.

\subsubsection{Why a second signature scheme at all}

ML-DSA and SLH-DSA are both post-quantum, but they are not post-quantum for
the same reason, and that is the entire point of carrying both.

ML-DSA-65's security rests on Module-LWE, a structured lattice assumption
roughly fifteen years old. It is believed hard. It is not *proven* hard, and
the structure that makes it fast — the ring, the NTT — is also the structure
a future attack would exploit. A cryptanalytic break of Module-LWE would
forge every transaction this chain has ever accepted.

SLH-DSA's security rests on nothing but the preimage and collision
resistance of SHA-2. There is no algebraic structure to attack. If SHA-2
falls, far more than this chain is already lost. The two schemes therefore
fail independently, which is what makes signing under both worth its cost:
a forgery needs to break lattices *and* hashes, not either one.

\subsubsection{Why the `s` parameter set and not `f`}

FIPS 205 offers "small" and "fast" variants at each security level. Both
were measured on the reference hardware before this crate was written, in
one release-profile run so the two columns are comparable:

\begin{verbatim}
| | SHA2-128s | SHA2-128f |
|---|---|---|
| signature | 7856 B | 17088 B |
| sign | ~196 ms | ~9 ms |
| verify | ~0.22 ms | ~0.56 ms |
\end{verbatim}

(\texttt{cargo bench --bench hybrid\_signing} reports a faster \texttt{128s} sign than this
— around 105 ms — because criterion warms the cache across samples. The
ratio between the two parameter sets is what this table is for, and that
holds either way.)

\texttt{f} signs twenty times faster, which looks decisive until you ask who pays.
Signing happens once, in a wallet, per transaction. Verification happens on
every node, for every transaction, in every block, forever — and \texttt{s}
verifies 2.6x faster while producing a signature less than half the size.
For a replicated ledger the asymmetry runs the other way from most systems:
the one-time cost is the cheap one.

\subsubsection{Every 32-byte string is a valid verifying key}

Unlike ML-DSA, where a public key is a structured encoding that can be
malformed, an SLH-DSA verifying key is just \texttt{PK.seed ‖ PK.root} — two hash
outputs. There is no validity condition to check and
[\texttt{VerifyingKey::from\_bytes}] is infallible. A garbage key is a key that
verifies nothing, not a key that fails to decode.

\subsubsection{Deterministic signing}

Signatures use FIPS 205's deterministic variant (\texttt{opt\_rand = PK.seed}),
matching the ML-DSA half. This is forced by consensus: the transaction id
hashes the signature, so a hedged signer would give one authorized payload
two different ids.

\subsection{suite/ed25519.rs}
Suite \texttt{0x01}: Ed25519, for devnet and legacy data only.

Zero post-quantum security — Shor's algorithm recovers the key from the
public key — so the registry marks it \texttt{mainnet\_allowed: false} and the
audit flags it. It exists because devnet tooling and some legacy records
use it, and reading legacy data is the point of crypto-agility.

**\texttt{verify\_strict} only.** Plain \texttt{verify} accepts small-order public keys
and the mixed-order encodings that make Ed25519 signatures malleable
between implementations; in a consensus rule, "two correct nodes disagree
about validity" is a fork. \texttt{verify\_strict} rejects both.

\subsection{suite/hybrid.rs}
Suite \texttt{0x30}: ML-DSA-65 \textbf{and} SLH-DSA-SHA2-128s.

Valid only if both halves verify. A forger has to break Module-LWE *and*
the hash function, so the suite survives a cryptanalytic surprise in either
family alone — the reason every transaction on the chain has carried this
pair.

**Byte-compatible with \texttt{custom\_l1\_node::crypto::hybrid}.** The public key
is \texttt{ml\_dsa\_pk ‖ slh\_pk}, the signature \texttt{ml\_dsa\_sig ‖ slh\_sig}, and keys
derive from a 32-byte chain key under the node's own domain strings, so an
existing account's key and every signature it ever made are a valid \texttt{0x30}
key and signature unchanged. \texttt{crates/node/tests/suite\_parity\_tests.rs} pins
that against the node's \texttt{fips204} implementation.

\textbf{Both halves are always checked.} No early return after the first
failure, so the time taken does not tell a forger which half they got
wrong.

\subsection{suite/ml\_dsa.rs}
Suites \texttt{0x10} and \texttt{0x11}: ML-DSA-65 and ML-DSA-87 (FIPS 204).

\textbf{Deterministic signing, empty context.} FIPS 204 permits both hedged and
deterministic signing; the node has always signed deterministically
(\texttt{custom\_l1\_node::crypto::keys}), and a suite that signed differently would
make the \texttt{0x30} hybrid's lattice half disagree with the node byte for byte.
Domain separation lives inside the signed message, where the caller puts
it, so the FIPS context string is empty — a second place to put a domain is
a second encoding of the same intent that also verifies.

Keys are held in seed form (\texttt{ξ}, 32 bytes). The expanded key is recomputed
from it, which costs a few microseconds and means only 32 secret bytes ever
need wiping.

\subsection{suite/mod.rs}
Signature suites: a closed registry, one trait, and a dispatcher.

Every signature the chain accepts is named by a one-byte [\texttt{SuiteId}], and
the id travels with the signature in an [\texttt{crate::envelope}]. An algorithm
can then be retired or introduced without reinterpreting a single byte that
was written under the old one — the property ADR-007 calls crypto-agility.

\subsubsection{Why the registry is closed}

A suite that exists on some nodes and not others is a fork, so the set of
ids is a \texttt{\#[repr(u8)]} enum and an unknown byte is a decode error, never a
fallback. Governance chooses *among* compiled suites ([\texttt{crate::agility}]);
adding one is a code change with its own KAT vectors.

\subsubsection{Sizes}

Every length below is asserted at compile time against the backing crate's
own constants, and the ACVP tests in \texttt{tests/acvp\_tests.rs} check the
encodings against NIST's. They are final FIPS 204/205 sizes: ML-DSA-65 is
3,309 bytes, not the 3,293 of the 2023 draft.

\subsection{suite/slh\_dsa.rs}
Suites \texttt{0x20} and \texttt{0x21}: SLH-DSA-SHA2-128s and SLH-DSA-SHAKE-256f
(FIPS 205).

Stateless hash-based signatures: their security rests on the hash
function alone, which is why they are the conservative half of the hybrid
and the cold-vault suite. \texttt{128s} is small and slow to sign; \texttt{256f} is
category 5 and fast to sign, at 49,856 bytes a signature — acceptable for a
vault that signs rarely, unacceptable for a payment.

Signing is deterministic (\texttt{opt\_rand} absent) with an empty context, for
the same reasons as ML-DSA: see \texttt{ml\_dsa.rs}.

This is the crate that instantiates \texttt{slh-dsa} (invariant 2), and these are
concrete types, so the monomorphized code lands here, where the dev
profile optimizes it.

\section{Verifiable random function}
\noindent\texttt{crates/vrf\textbackslash{}src}

\subsection{ecvrf.rs}
RFC 9381 \texttt{ECVRF-EDWARDS25519-SHA512-TAI}.

\subsubsection{What the domain separators are doing}

Every hash in this file begins with the suite octet and a one-byte purpose
tag, and ends with a trailing zero. That looks like ceremony and is not.
Three different hashes here take curve points as input — hash-to-curve, the
challenge, and the output — and without the tags a value computed for one
could be presented as a value computed for another. The trailing byte closes
the other end: without it, a variable-length input could be re-cut so that
two different inputs produce one hash.

The tags are consensus. Changing one changes every proof the chain has ever
accepted.

\subsubsection{Why the proof is deterministic}

The nonce comes from hashing the secret key's nonce prefix together with the
hashed input, never from an RNG. Two consequences, both load-bearing:

- A prover cannot grind. Given a key and an input there is one proof and one
  output, so a beacon operator who dislikes today's randomness has nothing
  to try again with. That is the entire reason a VRF is usable as a beacon.
- A repeated proof does not leak the key. With a random nonce, two proofs of
  the same input under a reused nonce would expose the secret scalar by
  subtraction — the failure that has broken ECDSA deployments repeatedly.

\subsubsection{Timing}

Proving is the only side with a secret, and every scalar operation it
performs is a constant-time \texttt{dalek} primitive. Verification uses variable-
time multiplication deliberately: everything it touches is public, and
constant time there would cost every node speed to protect nothing.

\subsection{error.rs}
Why a VRF operation failed.

No \texttt{thiserror} here, unlike the node. This crate is meant to be linkable
into a standalone beacon operator with as little behind it as possible, and
a derive macro is a proc-macro dependency in a crate whose whole argument
for existing is that it links light.

\subsection{keys.rs}
VRF key material.

\subsubsection{The secret key is a seed, not a scalar}

RFC 9381's edwards25519 suite reuses RFC 8032's key format: the stored 32
bytes are a seed, and the scalar used for arithmetic is derived from it by
hashing and clamping. The other half of that hash becomes the nonce prefix,
which is what makes proving deterministic — a VRF whose proof depended on
fresh randomness would leak the secret key across two proofs of the same
input, and a VRF that used a *predictable* nonce would leak it outright.

Callers therefore never see a scalar. They see a seed, and every derivation
from it happens here.

\subsection{lib.rs}
The chain's verifiable random function: RFC 9381
\texttt{ECVRF-EDWARDS25519-SHA512-TAI}.

A VRF is a keyed hash whose output nobody can predict without the secret key
and nobody can dispute once the proof is published. That combination is what
makes it usable as a randomness beacon: the holder cannot choose the output,
because there is exactly one valid output per (key, input), and everyone
else can check it.

\subsubsection{This is a classical primitive on a post-quantum chain}

Stated first, because it is the most important thing about this crate and
the least visible from its API.

\texttt{custom-l1-node} removed ed25519 from transaction authorization on purpose
and replaced it with ML-DSA-65 alongside SLH-DSA, so that "an adversary who
broke either scheme in isolation gets nothing". This crate reintroduces
curve25519 — not for authorization, but for randomness.

There is no standardized, practical post-quantum VRF today. So this is a
considered trade-off rather than an oversight, and two things bound it:

1. \textbf{A break costs future unpredictability, not past outputs.} An adversary
   who recovers a VRF secret key can predict randomness from that moment on.
   Randomness already consumed by an already-committed block stays exactly
   as sound as it was.
2. **Every key carries a [\texttt{VrfScheme}] tag from the first block.** A
   post-quantum VRF lands as a second scheme beside this one, exactly as
   SLH-DSA landed beside ML-DSA. Without the tag on day one, that migration
   would be a hard fork of every stored key — which is the mistake this
   exists to avoid, and it costs one byte.

\subsubsection{Why this is a crate and not a module}

Not for the reason \texttt{ledger-math} and \texttt{dex} are crates: this one has a
dependency graph, so Kani cannot compile it and no proof harness is coming.

It is a crate because the wallet, the CLI, and any external beacon operator
need to *produce* proofs without linking RocksDB, \texttt{libp2p}, and the SNARK
stack. A beacon that could only run inside a full node would be a beacon
nobody could operate.

\subsubsection{What pins the implementation}

RFC 9381's own test vectors, in \texttt{tests/rfc9381\_vectors.rs}. This is
consensus code and there is no second implementation on this chain to
disagree with it, so the check has to come from outside the repository or it
is a check against itself.

\subsection{scheme.rs}
Which VRF a key and a proof belong to.

\subsubsection{Why a one-byte tag exists before there is anything to distinguish}

There is exactly one scheme today, so the tag discriminates nothing. It
exists because of what happens when there are two.

The chain's signature layer went through this. ML-DSA arrived, then SLH-DSA
arrived beside it, and the wire format had to grow a version byte to tell
the frames apart — which invalidated every signature that predated it,
because the old frames had no way to say what they were. The transaction
decoder still carries the rejection path for versions 1 through 4 as a
result.

A VRF key stored without a scheme tag has the same problem waiting in it. A
post-quantum VRF is not standardized today and will be; on the day it lands,
either every stored key already says what it is, or the upgrade is a hard
fork of the key registry. One byte now, or a migration later.

\section{Transparent proofs and the shielded pool}
\noindent\texttt{crates/zk-stark\textbackslash{}src}

\subsection{config.rs}
The one STARK configuration — ADR-008.

- \textbf{PCS:} FRI over BabyBear with the *hiding* variant, so proofs reveal
  nothing about the witness beyond the public values.
- \textbf{Commitments and transcript:} Keccak-f[1600]. The outer layer of the
  proof rests on the hash with the longest cryptanalysis record; Poseidon2
  is used only *inside* the AIRs, where it is cheap to arithmetise.
- \textbf{FRI:} blowup 4 (\texttt{log\_blowup = 2}), 100 queries, 16 bits of query
  grinding — Plonky3's \texttt{new\_benchmark\_zk} preset. The security level this
  gives is *computed*, not asserted: [\texttt{crate::security\_bits}] runs
  Plonky3's own estimator, and \texttt{reports/02-crypto.md} records both the
  conjectured and the proven figure.
- \textbf{Blinding randomness:} fresh from the OS for every proving
  configuration. A fixed seed would make proofs deterministic in the
  witness, which is exactly what zero-knowledge must not be.

\subsection{credential.rs}
Selective disclosure — ADR-008.

> I know a claim value \texttt{v}, a subject, a schema and a blinding factor such
> that \texttt{leaf = compress(compress(v, schema, 0.. ‖ subject), blind)} is in
> the issuer's tree under \texttt{issuer\_root}; \texttt{v} satisfies the public
> predicate; and the position of that leaf in the revocation tree under
> \texttt{revocation\_root} holds the *unrevoked* leaf.

Public: the two roots and the predicate. The value, subject, schema,
blinding and index stay private. The issuer's signature is not in here —
the chain verifies it natively when the root is anchored, which is why
only membership under an anchored root has to be proved.

Values are below \texttt{2\textasciicircum{}29} so that \texttt{v = bound + d} in the field is the
integer equation (\texttt{bound + d < 2\textasciicircum{}30 < p}). The predicate is coarse on
purpose — \texttt{AtLeast(18)}, not a birth date — because a fine predicate
leaks through the public input however private the value is.

Unlike the Groth16 presentation this replaces, it is post-quantum and has
no setup.

\subsection{dual.rs}
Two Merkle paths side by side, one Poseidon2 permutation for each per row.

Shared by the credential and sanctions statements, which both need two
paths whose *positions* are related: the credential's leaf and its
revocation bit sit at the same index of two trees, and a sanctions gap is
two leaves at adjacent indices of one tree.

Layout, \texttt{DEPTH = 16}, 32 rows:

\begin{verbatim}
| row | path A | path B |
|---|---|---|
| 0 | leaf pre-image, first hash | leaf pre-image, first hash |
| 1 | `compress(row-0 digest, second)` → leaf | the same → leaf |
| 2 … 17 | Merkle levels 0 … 15 | Merkle levels 0 … 15 |
| 18 … 31 | padding | padding |
\end{verbatim}

Each path carries its own bit and sibling columns and an index accumulator
\texttt{acc}, which starts at 0 and adds \texttt{bit · 2\textasciicircum{}level} on each level row, so on
the last level row it holds the leaf's index — small enough (\texttt{< 2\textasciicircum{}16}) to
be compared in the field.

\subsection{gadgets/key.rs}
Knowledge of a hash-based signing key's secret, bound to a message.

The key pair is \texttt{pk = compress(sk, KEY)}. The proof shows knowledge of an
\texttt{sk} hashing to the public \texttt{pk}, and the message is a public value: the
verifier's Fiat–Shamir transcript absorbs every public value before it
draws a challenge, so a proof made for one message fails for another. That
makes the proof a signature of knowledge — post-quantum, because it rests
only on Poseidon2 preimage resistance.

This is the key-ownership piece of the shielded pool: a note's owner
proves they hold the \texttt{sk} behind the \texttt{pk} committed in the note.

\subsection{gadgets/merkle.rs}
Merkle path: a leaf is a member of the tree with a given root.

One row per level, \texttt{depth} rows, \texttt{depth} a power of two. Row \texttt{i} holds the
node entering level \texttt{i} (\texttt{cur}), its sibling, and the path bit: with the
bit 0 the permutation input is \texttt{cur ‖ sib}, with the bit 1 \texttt{sib ‖ cur}, and
the next row's \texttt{cur} is this row's digest. The first row's \texttt{cur} is the
public leaf, the last row's digest the public root.

Set membership is this statement with the set committed as the tree.

\subsection{gadgets/mod.rs}
Reusable AIR gadgets.

Each gadget is a small, standalone statement with its own AIR, native
witness check, prover and verifier, and each is also the building block
the shielded pool's AIRs repeat:

\begin{verbatim}
| Gadget | Proves | Public |
|---|---|---|
| [`range`] | a committed value is below `2^bits` (≤ 60) | the commitment |
| [`merkle`] | a leaf is in the tree with a given root (set membership) | leaf, root |
| [`key`] | knowledge of the secret behind a hash-based public key, bound to a message | public key, message |
\end{verbatim}

\subsection{gadgets/poseidon.rs}
One Poseidon2 permutation per trace row, embedded at a column offset.

Every gadget and pool AIR in this crate is a row-per-hash design: columns
\texttt{[0, POSEIDON\_COLS)} of each row hold the upstream \texttt{p3-poseidon2-air}
permutation columns, evaluated through Plonky3's \texttt{SubAirBuilder}, and the
gadget's own columns follow. The gadget then constrains what goes *into*
each permutation (\texttt{inputs}) and what comes *out* (the last round's
\texttt{post}), which is all it needs to know about the hash.

\subsection{gadgets/range.rs}
Range proof: a committed value is below \texttt{2\textasciicircum{}bits}.

The value is split into two 30-bit limbs so every limb, and every sum the
constraints form, stays below the BabyBear modulus (\texttt{p ≈ 2\textasciicircum{}31}) — there is
no wraparound for a prover to hide an out-of-range value in. The
commitment is \texttt{C = P(lo, hi, 0, 0, 0, 0, 0, 0 ‖ r)[..8]} with an 8-element
blinding \texttt{r}, so the proof reveals \texttt{C} and nothing else about the value.

Row 0 carries the statement; row 1 is an honest permutation of zero that
only exists because the hiding PCS needs more than one row.

\subsection{hash.rs}
Poseidon2 over BabyBear, width 16 — the one hash inside every AIR here.

The round constants are the standard ones \texttt{p3-baby-bear} ships
(\texttt{BABYBEAR\_POSEIDON2\_RC\_16\_*}, the Horizen Labs instance), not random
ones: a proof system whose hash constants came from a seed somebody chose
is a proof system with a trapdoor somebody might have chosen. The native
permutation and the AIR's [\texttt{round\_constants}] are built from the same
arrays, and \texttt{tests/gadget\_tests.rs} checks the AIR's output equals the
native permutation's.

Two-to-one compression is the truncated permutation \texttt{P(l ‖ r)[..8]} — the
construction Plonky3's own Merkle trees use. A digest is 8 field elements,
about 248 bits, so collision resistance is \textasciitilde{}124 bits classically and a
generic quantum collision search (BHT) gets no further than \textasciitilde{}83; the
commitment and Merkle hashing in this crate rest on that.

\subsection{lib.rs}
Transparent, post-quantum zero-knowledge proofs — ADR-008.

One proof system: Plonky3 uni-STARKs over BabyBear with hiding FRI and
Keccak commitments. No trusted setup, no pairing, nothing a quantum
computer's discrete-log algorithm breaks; soundness rests on hash
collision resistance. \texttt{tests/pqc\_zk\_tests.rs} checks the dependency graph
has no Groth16, BN254, BLS12-381 or halo2 left in it.

- [\texttt{hash}] — Poseidon2 (width 16), the in-circuit hash, with the standard
  constants.
- [\texttt{gadgets}] — reusable AIRs: range proof, Merkle path (set membership),
  and knowledge of a hash-based signing key's secret.
- [\texttt{pool}] — the shielded pool: notes, commitments, nullifiers, the
  commitment tree, and the 2-in-2-out joinsplit that shields, transfers
  and unshields.
- [\texttt{credential}] — selective disclosure: membership, a predicate on the
  value, and an unrevoked bit at the same index.
- [\texttt{sanctions}] — non-membership in a sorted list, by an adjacent-leaf gap.

Every constraint has a negative test that tampers with one witness
consistently — everything downstream recomputed — so only the constraint
under test can refuse it (invariant 23's lesson, carried over from the
zkML circuit this crate replaces).

\subsection{pool/joinsplit.rs}
The joinsplit AIR: two notes in, two notes out, with transparent edges.

128 rows, one Poseidon2 permutation each; roles fixed by periodic selector
columns the verifier supplies, so no witness can switch a check off.

\begin{verbatim}
| rows | role |
|---|---|
| `36i + 0` | input `i` key: `sk ‖ KEY` → address |
| `36i + 1` | input `i` value layer: `lo, hi, 0.. ‖ address` |
| `36i + 2` | input `i` commitment: `layer ‖ rho ‖ rand` |
| `36i + 3 … 36i + 34` | input `i` Merkle path, 32 levels, up to the anchor |
| `36i + 35` | input `i` nullifier: `sk ‖ NULLIFIER, rho, 0..` |
| `72 + 2j`, `73 + 2j` | output `j` value layer and commitment |
| 76–127 | padding: honest permutations, nothing else |
\end{verbatim}

\textbf{Dummy inputs.} An unused input slot carries a zero-value note with
\texttt{is\_dummy = 1}: its value is forced to zero and its Merkle root is not
compared with the anchor (there is no leaf to find), but its key,
commitment and nullifier are computed like any other, so the published
nullifier still binds to *something* and nothing on chain tells a dummy
from a real spend. Marking a real note dummy gains nothing: its value
becomes zero.

\textbf{Conservation} over 27-bit limbs, with \texttt{A = in₀ + in₁ + public\_in} and
\texttt{B = out₀ + out₁ + public\_out + fee}:

\begin{verbatim}
A_lo + 4·2^27 = B_lo + c·2^27        A_hi + c = B_hi + 4        0 ≤ c ≤ 7
\end{verbatim}

which says \texttt{A = B} as integers: \texttt{A\_lo ≤ 3(2\textasciicircum{}27−1)}, \texttt{B\_lo ≤ 4(2\textasciicircum{}27−1)}, so
every side is below \texttt{11·2\textasciicircum{}27 < p ≈ 15·2\textasciicircum{}27} and nothing wraps.

\textbf{Public values} (57): anchor, both nullifiers, both output commitments,
the limbs of \texttt{public\_in}, \texttt{public\_out} and \texttt{fee}, and the recipient as
eleven 24-bit words — bound by the transcript, so a miner who rewrites the
recipient invalidates the proof.

\subsection{pool/mod.rs}
The shielded pool's statement — ADR-008.

One AIR, [\texttt{joinsplit::JoinSplitAir}]: two notes in, two notes out, value
entering and leaving through public amounts. Shielding (transparent → note),
transfer (notes → notes) and unshielding (notes → transparent) are the same
statement with different edges, built by [\texttt{wallet}]. The state machine —
anchor window, nullifier set, pool balance — belongs to the chain
(\texttt{crates/node/src/state/shielded.rs}); this crate proves and verifies.

\subsection{pool/note.rs}
Keys, notes, commitments and nullifiers.

\begin{verbatim}
address = compress(sk, KEY)
cm      = compress(compress(lo, hi, 0.. ‖ address), rho ‖ rand)
nf      = compress(sk, NULLIFIER, rho, 0, 0, 0)
\end{verbatim}

Values are two 27-bit limbs, so a note holds less than \texttt{2\textasciicircum{}54} (about
1.8·10\textasciicircum{}16 base units). 27 rather than 30 bits because the joinsplit adds up
to four limbs in one equation and every such sum must stay below the
BabyBear modulus, \texttt{p ≈ 15·2\textasciicircum{}27} — see \texttt{joinsplit.rs}.

\texttt{rho} (4 elements, \textasciitilde{}124 bits) makes every note unique, so every nullifier
is; \texttt{rand} hides the commitment. The nullifier needs \texttt{sk}, so only the
owner can compute it, and it is a function of the note, so spending a note
twice publishes the same nullifier twice.

\subsection{pool/tree.rs}
The note-commitment tree: incremental, depth 32, stored as a frontier.

A node keeps only the *frontier* — for each level, the last left child
still waiting for a right sibling — plus the leaf count. That is 32
digests however many notes exist, and it is enough to append and to
compute the root. Wallets, which do need paths, rebuild them from the
leaves they watched ([\texttt{merkle\_path}]).

Empty positions are the zero digest and an empty subtree is the hash of
two empty ones, so the root of a partly filled tree is well defined.

\subsection{pool/wallet.rs}
Building joinsplits from a wallet's point of view.

Every joinsplit has exactly two inputs and two outputs; the builder pads
with zero-value dummies, so an observer cannot tell a one-note spend from
a two-note one.

\subsection{pool/witness.rs}
Joinsplit witness → trace and public values, with the native checks.

\subsection{proof.rs}
Proving, verifying, and measuring what a proof is worth.

\subsection{sanctions.rs}
Proving an identifier is \textbf{not} on a sanctions list, without revealing it.

> I know a 32-byte identifier \texttt{x} and two leaves \texttt{lo < x < hi} at adjacent
> positions \texttt{i}, \texttt{i + 1} of the sorted list committed to by \texttt{root}.

Adjacency is what makes a gap a proof: two *neighbours* of a sorted list
leave no room for \texttt{x} between them. The two sentinels (all-zero and
all-\texttt{0xff}) make every identifier bracketable. Identifiers compare as
big-endian bytes — \texttt{<[u8; 32]>::cmp} natively — held in the AIR as sixteen
16-bit limbs, and \texttt{x − lo − 1 ≥ 0} and \texttt{hi − x − 1 ≥ 0} are proved by
limb-wise borrow chains whose differences are range-checked, so a negative
result cannot hide in the field.

What this does not prove: that \texttt{x} is the party in the payment (bind it
by committing to \texttt{x} in the payment), or that \texttt{root} is the right list
(the verifier must know which root it expects).

\section{Contract execution}
\noindent\texttt{crates/vm\textbackslash{}src}

\subsection{cache.rs}
Compiling a contract once instead of on every call.

\subsubsection{What this is, and what it is not}

It is \textbf{not} a JIT. The VM has been a JIT since it was written: \texttt{Config} at
\texttt{crate::config} sets \texttt{Strategy::Cranelift}, so wasmtime compiles every module
to native x86-64 or aarch64 before it runs, and there is no interpreter in
this tree to be slower than.

What was missing is that the compile happened \textbf{again on every call}.
\texttt{Vm::call} handed the same bytes to \texttt{Module::new} each time and Cranelift
optimised them from scratch. This holds the result.

\subsubsection{Why a cache cannot change gas}

Gas is wasmtime fuel (invariant 21), and fuel is instrumentation the compiler
injects per Wasm operator. The same bytes under the same configuration
produce the same instrumentation, so a cached module charges exactly what a
freshly compiled one charges. \texttt{a\_cached\_call\_burns\_identical\_fuel} asserts
that rather than assuming it, because "gas is unchanged" is the entire safety
argument for caching anything on a consensus path.

\subsubsection{The configuration is part of the key}

Not just the bytecode. A cache keyed on bytes alone would, after a change to
\texttt{deterministic\_config} — NaN canonicalisation, a proposal flag, the stack cap
— keep serving modules compiled under the *old* settings until the process
restarted. Two nodes would then be running the same contract under two
compilers, which is the one way a cache can fork a chain.

So the key is \texttt{BLAKE3(config digest || bytecode)}, and a configuration change
invalidates every entry by construction rather than by anybody remembering.

\subsubsection{Bounded, and evicting the least recently used}

An unbounded cache is a node that dies from the number of distinct contracts
it has seen — which is a number an attacker chooses, by deploying junk. So
there is a cap, and the entry not touched for longest goes first: a contract
called once by an attacker should not evict one called every block.

\subsubsection{Not shared between engines}

A \texttt{Module} is bound to the \texttt{Engine} that compiled it, and instantiating it
against another is an error rather than a subtle wrong answer. The cache
lives inside \texttt{Vm} beside its engine for that reason, and the config digest in
the key is a second guard on the same thing.

\subsection{config.rs}
Deterministic engine configuration.

Every validator must reach an identical result — same output, same gas, same
success or failure — for the same contract and input. WebAssembly is
*mostly* deterministic, but several optional features and defaults are not,
and each has to be switched off explicitly.

\subsubsection{The version pin}

wasmtime makes no guarantee that fuel accounting is stable across releases.
The per-instruction schedule can change with any version. A validator on one
version and a validator on another can therefore disagree about whether a
call ran out of gas — which is a state divergence and a chain split.

The dependency is pinned to an exact version for that reason, and
[\texttt{WASMTIME\_VERSION}] records it so a test can assert the pin has not drifted.
Bumping wasmtime is a hard fork, not a dependency update.

\subsubsection{What each setting prevents}

\begin{verbatim}
| Setting | Without it |
|---|---|
| `consume_fuel` | contracts can loop forever |
| `cranelift_nan_canonicalization` | NaN bit patterns differ between CPUs |
| `wasm_simd` off | SIMD lowering varies by host feature detection |
| `wasm_threads` off | shared memory makes execution order observable |
| `wasm_reference_types` off | reachability, and thus GC timing, is observable |
| fixed memory ceiling | a host-dependent limit forks on memory-heavy calls |
\end{verbatim}

\subsection{error.rs}
VM errors.

\subsection{host.rs}
The host interface exposed to contracts.

\subsubsection{Why this is a trait}

The VM crate deliberately does \textbf{not} depend on the node crate. The node
depends on the VM in order to execute contracts, so a dependency the other
way would be a cycle. Everything the VM needs from chain state arrives
through [\texttt{HostState}], which also makes the VM testable against an in-memory
stub instead of a database.

\subsubsection{Determinism obligations on an implementor}

Every method must be a pure function of committed chain state. No clocks, no
randomness, no network, no filesystem, no iteration order that depends on
hashing. A single non-deterministic host function undoes every precaution the
engine configuration takes.

\subsection{lib.rs}
\subsubsection{Maya VM}

A sandboxed, deterministic WebAssembly execution engine for the Maya2C chain,
built on wasmtime.

\subsubsection{Determinism is the whole job}

Every validator must agree on a contract call's output, its gas, and whether
it succeeded. Disagreement on any of those is a state divergence and a chain
split. Three things enforce it:

1. [\texttt{config::deterministic\_config}] disables every non-deterministic
   WebAssembly proposal and pins compilation to Cranelift.
2. The wasmtime dependency is pinned to an \textbf{exact} version, because fuel
   accounting is not stable across releases. See
   [\texttt{config::WASMTIME\_VERSION}].
3. [\texttt{host::HostState}] implementors must be pure functions of committed
   state. One host function that reads a clock undoes the other two.

\subsubsection{Layering}

This crate does not depend on the node crate. The node depends on *it* to
execute contracts, so the reverse would be a cycle. Chain state reaches the
VM through the [\texttt{host::HostState}] trait, which also makes the VM testable
against an in-memory stub.

\subsubsection{Gas}

Gas maps one-to-one onto wasmtime fuel and is a **halting bound, not a
price**. Exhausting it traps the call and discards its writes; nothing is
billed. A fee market is a separate design.

\subsection{runtime.rs}
Contract execution.

\subsubsection{ABI}

A contract exports:

\begin{verbatim}
memory
input_ptr()    -> i32   address of its static input buffer
input_cap()    -> i32   capacity of that buffer
invoke(len:i32)-> i64   (out_ptr << 32) | out_len ; negative means failure
\end{verbatim}

The host writes the call input into the contract's own buffer and passes the
length. There is no allocator in the ABI: a \texttt{\#![no\_std]} contract has no heap
by default, and requiring one would exclude the smallest contracts for no
benefit.

\subsubsection{Gas}

Gas maps one-to-one onto wasmtime fuel. Gas is a \textbf{halting bound}, not a
price — running out traps, the call fails, and every staged write is
discarded. Nothing is billed to the caller.

\subsubsection{Reverting}

Host writes go into the [\texttt{HostState}] the caller supplies. Because execution
either returns a result or an error, and the node applies state only on
success, a trap partway through leaves nothing behind. The
\texttt{gas\_exhaustion\_reverts\_all\_writes} test pins that.

\subsection{zkml.rs}
The zkML verification host function's contract with the rest of the VM:
its bounds, its price, and its answers.

The verifier itself is not here. \texttt{maya-vm} links no pairing library — the
node supplies verification through [\texttt{crate::HostState::verify\_zkml}], the
same seam that keeps chain types out of the VM. That keeps the VM's
dependency set to wasmtime, and it means a host that has not opted in has
no verifier at all rather than a disabled one.

\subsubsection{The price is measured, and it is large on purpose}

Gas is wasmtime fuel, about one unit per wasm instruction, and fuel cannot
see native work. A host function that verifies a pairing-based proof does
\textasciitilde{}6 ms of native work (release build, \texttt{crates/zkml-prover/benches/verify.rs}); a tight
guest loop on the same machine burns \textasciitilde{}16.6 million fuel per millisecond
(\texttt{crates/vm/tests/fuel\_calibration\_tests.rs}). So the honest charge is on the order
of a hundred million fuel, and [\texttt{ZKML\_VERIFY\_BASE\_FUEL}] is set at roughly
9 ms' worth for headroom.

That is fifteen times the VM's default gas ceiling, which means a contract
calling this must ask for a larger \texttt{gas\_limit}. That is correct: a call that
verifies a proof *is* fifteen times as much work as the default call budget,
and a price that pretended otherwise would let a contract loop over
verifications at a fraction of what they cost the validators running it.

The per-byte charge covers copying the buffers out of guest memory and
hashing the key. The verification itself does not grow with the byte
counts: the verifier accepts one circuit shape, reads a fixed number of
points, and refuses trailing bytes afterwards.

\subsubsection{Charged before any work}

The fuel is taken first. A call that cannot pay traps without a byte having
been read and without a pairing having been computed — otherwise
"out of gas" would arrive after the work it was meant to bound.

\section{Trading engine}
\noindent\texttt{crates/dex\textbackslash{}src}

\subsection{amm.rs}
Constant-product automated market maker.

\subsubsection{The invariant, stated precisely}

A swap must never decrease \texttt{reserve\_base * reserve\_quote}. That product is
the pool's solvency: it is what guarantees the reserves can honour every
outstanding share, and a swap that lowered it would have paid a trader out
of the providers' capital rather than out of the curve.

Fees are what make it *increase*, and that is the whole mechanism by which
providers are paid — see [\texttt{crate::fees}].

Every function here recomputes the product and refuses to return a pool that
violates it. That check is redundant with the arithmetic — the rounding
rules make a violation unreachable — and it stays because "unreachable"
rests on an argument about integer division that a future edit can quietly
break.

\subsubsection{A pool cannot be emptied}

Draining the last unit of a reserve would require infinite input, so
[\texttt{Pool::swap\_exact\_out}] rejects an output equal to the whole reserve rather
than computing a price for it. This is a property of the curve, not a
configured limit.

\subsubsection{The locked minimum}

[\texttt{MINIMUM\_LIQUIDITY}] shares are minted to nobody on the first deposit. The
attack it closes is specific: without it, the first depositor can mint one
share, donate a large amount directly into the reserves, and thereby make
one share worth so much that every later depositor's contribution rounds
down to zero shares and is absorbed. Locking a floor of shares makes the
share price bounded from the start, so the rounding never has that much
room.

\subsection{batch.rs}
Uniform-price batch clearing: the pool's defence against ordering.

\subsubsection{What this replaces}

Executed one at a time in the order a miner chooses, swaps against a curve
are trivially extractable. Put a buy in front of a victim's buy and a sell
behind it, and the victim's own price impact pays for the round trip. The
attack needs nothing but the ability to decide where in the block a
transaction goes, which is precisely the thing a miner has.

Slippage bounds do not fix this. A bound caps the loss; it does not remove
the incentive, and a bound tight enough to remove it is a bound that fails
constantly on ordinary volatility.

\subsubsection{What happens instead}

Every swap on a pair in one block is collected and settled at \textbf{one price}.
Position within the block stops existing as a variable. An attacker who
front-runs is in the same batch as the victim and receives the same price
they engineered; the sandwich costs them the fee and returns nothing.

The price is found by netting first. Buyers and sellers in the same batch
cross against each other, and only the imbalance reaches the curve. The
curve's execution price on that imbalance becomes everyone's price.

\subsubsection{What it does not solve, stated plainly}

- \textbf{A miner spanning consecutive blocks.} Batching is per block. A miner
  who mines two in a row can still put a trade in the first and unwind it in
  the third. Defeating that needs commit–reveal, which costs a block of
  latency for every trade; it is not implemented here.
- \textbf{Censorship.} A miner who drops a transaction rather than reordering it
  is unaffected by anything in this module.
- \textbf{Cross-pair routing.} Two pools are two batches.

\subsubsection{Coincidence of wants is free}

Only the imbalance touches the curve, so only the imbalance pays a fee.
Volume that crosses internally pays nothing, to anybody. That is a real
consequence and not an oversight: a batch that is perfectly balanced settles
at the spot price with no fee revenue at all. The providers are compensated
for the capital the batch actually used, which on a balanced batch is none.

\subsubsection{Rounding surplus}

Every payout is floored, so the batch almost always leaves a unit or two
behind. Those go to the reserves via [\texttt{crate::amm::Pool::donate}], raising
the share price. They are not distributed, not tracked, and not claimable —
they came out of trades against the providers' capital, and raising the
share price is how that capital is paid.

\subsection{book.rs}
The limit order book: an in-memory index built to be crossed quickly.

\subsubsection{Why the ordering key is the whole design}

Price-time priority is a total order on resting orders, and the fastest way
to maintain a total order is to make it the sort key rather than to compute
it. Every order here is filed under \texttt{(price\_key, sequence)}:

- \texttt{sequence} is a monotonic arrival counter, so ties in price break in
  arrival order without a timestamp — and without a timestamp there is
  nothing for a miner to backdate.
- \texttt{price\_key} is the price for an ask and \texttt{u64::MAX - price} for a bid, so
  that *both* sides are best-first in ascending order. Without the
  inversion, one side would have to be iterated backwards, and every
  function that walks a book would need to know which side it was holding.

Matching is then two forward cursors over sorted maps, and the "best order"
is whatever the cursor is already pointing at.

\subsubsection{The same key orders the database}

[\texttt{sort\_key}] emits the pair big-endian. RocksDB iterates keys
lexicographically, so a prefix scan over \texttt{<pair><side>} yields resting
orders in exactly the priority order this module matches them in — the book
is rebuilt in priority order for free, with no sort step and no opportunity
for two nodes to disagree about it.

Big-endian, specifically. The little-endian encoding used for every amount
on the wire would order \texttt{256} before \texttt{1}, and the priority queue would be
nonsense.

\subsubsection{This structure is a working index, not state}

A [\texttt{Book}] is built from committed records, crossed, and dropped. The
authoritative form of an order is its stored record; this is the shape that
form is put into to be worked on. That is why [\texttt{Book}] has mutating methods
while [\texttt{crate::amm::Pool}] has none — mutating a scratch index is local,
mutating a pool is a ledger write.

\subsection{error.rs}
Failure modes of the trading engine.

These are *arithmetic and rule* failures, not chain errors. The crate
returns its own type for the same reason [\texttt{maya\_ledger\_math}] returns
\texttt{Option}: keeping \texttt{custom-l1-node} out of the dependency graph is the point
of the crate boundary, and the mapping from "this trade cannot be
represented" to whichever \texttt{NodeError} variant a context calls for belongs to
the caller.

[\texttt{maya\_ledger\_math}]: https://docs.rs/maya-ledger-math

\subsection{fees.rs}
The fee split: what the liquidity providers keep and what the chain takes.

\subsubsection{Where an LP fee actually lives}

It is never paid out. The fee stays in the reserves, so \texttt{k} grows and every
outstanding share redeems for slightly more than it did before. There is no
accrual counter, no per-provider accumulator, and no claim transaction.

The alternative — a per-share fee accumulator, incremented on every swap and
settled when a provider withdraws — is what a protocol reaches for when fees
are paid in an asset the pool does not hold. Here they are paid in the asset
the pool does hold, so the accumulator would be a second representation of a
number the reserves already carry, kept in step by hand, and drifting by a
unit every time a division was inexact.

The consequence to be aware of: a provider's share count never changes, so
"how much have I earned" is \texttt{shares * reserves / supply} now minus what it
was at deposit, and nothing on chain records the second term. That is a
wallet's job, not the ledger's.

\subsubsection{Why the protocol cut comes off the input}

The protocol's share is taken from \texttt{amount\_in} *before* the curve sees it,
not skimmed from the pool afterwards. Skimming afterwards would mean
withdrawing from the reserves, which lowers \texttt{k} — and then "did this swap
decrease \texttt{k}" stops being a usable invariant, because the honest answer
becomes "yes, by exactly the amount we intended". An invariant with an
exception is not an invariant.

\subsection{lib.rs}
The trading engine for \texttt{custom-l1-node}: a constant-product AMM and a
price-time limit order book, with no I/O and no chain types.

\subsubsection{Why this is a crate and not a module}

The same argument that produced [\texttt{maya\_ledger\_math}]. Every function here
decides how much value moves, and a rounding error in one of them is a mint
or a burn rather than a cosmetic defect. That is the class of code worth
proving, and the [Kani Rust Verifier] compiles a crate *together with its
whole dependency graph* — which rules out anything that reaches RocksDB's
C++ or the \texttt{ark-*} stack.

So this crate has no dependencies, and the node calls into it rather than
inlining the arithmetic. What \texttt{proofs} (compiled only under Kani) establishes is a property of the
code that runs when a block is applied, not of a restatement of it.

The visible cost of that boundary is that nothing here hashes. Pair
identifiers, order identifiers, and LP asset identifiers arrive as opaque
32-byte values; deriving them is the node's job.

\subsubsection{The two venues}

[\texttt{amm}] is a constant-product pool: \texttt{x * y >= k}, fees retained in the
reserves. [\texttt{book}] and [\texttt{matching}] are a limit order book cleared by
price-time priority at the resting order's price.

They are two independent venues on the same pair, not one routed venue. A
swap reaches the pool; a limit order rests in the book; the block-end
matching pass crosses the book against itself. Routing a single order across
both — taking book liquidity up to the pool's marginal price and the
remainder from the curve — is a real feature and is deliberately *not*
implemented here, because a half-built router that silently picks the worse
venue is worse than no router. See \texttt{docs/dex.md}.

\subsubsection{Determinism}

Every function is a pure function of its inputs. There is no floating point,
no hashing, no iteration order that depends on an address's bit pattern, and
no time source. Two nodes handed the same inputs produce the same trades, or
the chain forks.

\subsubsection{Rounding}

One rule, applied without exception: **when a division is inexact, the
remainder goes to the pool.** A trader receiving output gets a floor; a
trader supplying input pays a ceiling. The alternative is not "fairer", it
is a slow leak that a bot can drain in a loop.

[\texttt{maya\_ledger\_math}]: https://docs.rs/maya-ledger-math
[Kani Rust Verifier]: https://model-checking.github.io/kani/

\subsection{matching.rs}
Crossing the order book.

\subsubsection{The rule}

While the best bid is at or above the best ask, the two trade. The price is
the \textbf{maker's} — the order that arrived first, identified by the lower
sequence number. The party who moved second is the one who chose to accept a
price that was already on the book, so they get that price and not their
own.

That single rule is what makes queue position worth having. Without it,
posting early would buy nothing, and there would be no reason for anyone to
rest an order rather than wait.

\subsubsection{Bounded work}

Matching is the only place in the state transition where one transaction can
cause an amount of work that is a function of *other people's* orders. A
deep book crossed by one large order is thousands of fills, every one of
which every node must redo.

So the caller supplies a ceiling, and reaching it stops the pass with
[\texttt{MatchOutcome::limit\_reached}] set rather than by returning an error. A
crossed book left behind is not a failure — it is a book with more to do
next block, which is exactly what a ceiling means. Turning it into an error
would make one large order able to invalidate a block.

\subsubsection{Fees on a book trade}

Only the protocol's cut applies. The LP rate in a [\texttt{FeeSchedule}] pays
liquidity providers for the use of pooled capital, and a book trade uses
none: the counterparty is another trader. Charging it here would be taking a
fee on behalf of a party that did not participate.

The cut comes off the \textbf{taker's} output. The maker posted a price and had
it accepted; the taker consumed liquidity that was resting there for them.

\subsection{math.rs}
Wide-intermediate integer arithmetic.

Every product formed here is \texttt{u64 * u64}, which overflows \texttt{u64} and does not
overflow \texttt{u128}. Computing in \texttt{u128} and reducing back is what makes the
reserve arithmetic exact instead of merely usually-correct: the naive
\texttt{reserve\_out * amount\_in / (reserve\_in + amount\_in)} in \texttt{u64} wraps for any
pool holding more than about four billion units on each side, which is not a
large pool.

The two rounding directions are separate functions rather than a flag,
because which one a call site wants is a *safety* property. A payout that
rounds up and a charge that rounds down are the same bug — value created
from nothing — and naming them apart makes that visible at the call site.

\subsection{proofs.rs}
Kani harnesses for the trading arithmetic.

Run them with:

\begin{verbatim}
cargo kani -p maya-dex
\end{verbatim}

\subsubsection{What is worth proving here}

Not "the swap returns roughly the right number" — a test does that better
and faster. What a model checker is for is the *edges*: the reserve sizes
where a \texttt{u64} product wraps, the fee rates where an intermediate rounds to
zero, the one input in \texttt{2\textasciicircum{}64} where a floor and a ceiling disagree by enough
to matter. Those are exactly the inputs nobody writes a test for, and
exactly the inputs an attacker looks for.

Every property below is a *conservation* or *monotonicity* statement. None
of them says what a swap should return; they say what it must never be
possible for one to do.

\subsubsection{Bounded versus unbounded}

Stated plainly, because the difference is the difference between a proof and
a very good test:

- [\texttt{mul\_div\_floor\_and\_ceil\_bracket\_the\_true\_quotient}],
  [\texttt{mul\_div\_rounding\_differs\_by\_at\_most\_one}] and
  [\texttt{isqrt\_is\_the\_floor\_of\_the\_real\_root}] are \textbf{unbounded} over every input
  their arguments can take.
- Every pool harness constrains the reserves and the fee rate with
  \texttt{kani::assume}. Those are not simplifications: each one mirrors a check
  the function itself performs and returns an error for, so what is proved
  is the behaviour on inputs the function accepts. The one genuine
  restriction is the reserve *ceiling*, noted at [\texttt{PROVED\_RESERVE\_CEILING}].
- Nothing in [\texttt{crate::matching}] or [\texttt{crate::batch}] is proved. Both loop
  over attacker-supplied collections, so a harness for either would be
  bounded at a length far below the ceiling the code actually accepts, and a
  bounded proof of a loop is worth less than the property tests in
  \texttt{crates/dex/tests/} that run it at full width. Said here rather than left to be
  inferred from the absence.

\subsection{types.rs}
Identifiers, units, and the scale every price is quoted in.

\section{Governance}
\noindent\texttt{crates/governance\textbackslash{}src}

\subsection{error.rs}
Why a governance operation was refused.

\subsection{lib.rs}
On-chain governance: the proposal lifecycle, the vote tally, and the bounds
a proposal may never escape.

\subsubsection{Why this is a crate and not a module}

The same argument that produced [\texttt{maya\_ledger\_math}] and \texttt{maya-dex}. What
this crate decides is which changes to the chain's own rules are
permissible, which makes it the last place a mistake should be hard to
check — and the [Kani Rust Verifier] compiles a crate together with its
whole dependency graph, so anything reaching RocksDB's C++ cannot be model
checked at all.

So it has no dependencies, and the node calls into it rather than inlining
the rules. The visible cost is that nothing here hashes or stores: proposal
identifiers, authors, and the changes themselves are chain records, and this
crate sees only what it needs to decide what may happen next.

\subsubsection{The one idea worth reading before the code}

\textbf{Governance must not be able to make governance unsafe.}

A system that can change every rule can change the rules that protect the
process of changing rules, and the first move of a hostile majority is not
to steal — it is to remove the checks that would have let anyone react. So
the quorum floor, the approval floor, the minimum voting period, and the
minimum timelock live in [\texttt{limits}], are compiled into the binary, appear in
no [\texttt{params::ParameterKey}], and are reachable by no transaction. Changing
them is a release, which is to say it is a decision every node operator
makes individually by choosing what to run.

Everything that *is* governable carries a hard range in [\texttt{params}], checked
when a proposal is made and again when it executes.

\subsubsection{What self-amendment means here, precisely}

A governed rule is a \texttt{u64} in consensus state with an activation height.
That covers scalar limits, rates, resource bounds — and, using the same
table, flipping between two rules the binary already ships, which is how
\texttt{crypto::dag::registry} already decides whether a block's digest is
ArgonBlake or the DAG.

It does \textbf{not} cover fetching native code from chain state and running it.
See [\texttt{params}] for why that is not a limitation to be worked around.

[\texttt{maya\_ledger\_math}]: https://docs.rs/maya-ledger-math
[Kani Rust Verifier]: https://model-checking.github.io/kani/

\subsection{limits.rs}
The bounds governance may never escape.

\subsubsection{Why a governed chain needs ungoverned floors}

Governance that can change anything can break everything, and the first move
of a hostile majority is not to steal — it is to remove the checks that would
have let anyone react. Set the timelock to zero and the exit window
disappears; set the quorum to one and the majority stops needing to be one.

So the parameters that protect the process are \textbf{not proposable at all}.
They are compiled into the binary, they appear in no
[\texttt{crate::params::ParameterKey}], and there is no transaction that changes
them. Changing them is a release and a hard fork, which is to say it is a
decision every node operator makes individually by choosing what to run.

That is a real limitation and it is the point: a chain whose safety margins
can be voted away has margins only until somebody wants them gone.

\subsubsection{Everything else is range-checked}

Parameters that *are* proposable carry a hard \texttt{[min, max]} in
[\texttt{crate::params}], checked when the proposal is made and \textbf{again} when it
executes. Twice, because a range could be tightened by a release between
those two moments, and a value that was legal when proposed must not become
law after it stopped being legal.

\subsection{params.rs}
What a proposal is allowed to change, and to what.

\subsubsection{The shape of self-amendment here}

Every governable rule is a \texttt{u64} under a [\texttt{ParameterKey}], stored in
consensus state, and read at execution time by whichever subsystem owns it.
A change carries an activation height, so the value that applies at height
*h* is a function of *h* alone — which is what makes two nodes replaying the
same chain reach the same answer.

That covers two of the three shapes a rule change can take:

\begin{verbatim}
| Shape | Example |
|---|---|
| A scalar limit or rate | the pool protocol fee, a per-block ceiling |
| A resource bound the VM enforces | contract memory pages, module size |
\end{verbatim}

The third — \textbf{flipping between two rules} — uses the same table. The binary
ships both implementations and a parameter decides which applies above a
height. \texttt{crypto::dag::registry} already works this way, and
\texttt{core::block::pow\_hash\_at} reads it: below the activation height the digest
is ArgonBlake, at or above it is the DAG. Governance moving such a height is
the same operation as governance moving a fee.

\subsubsection{What is not here, and will not be}

\textbf{Native code.} There is no key whose value is a program. A node that
fetched executable code from chain state and ran it would be handing
arbitrary code execution on every machine in the network to whoever wins a
vote, and no timelock makes that acceptable.

Substrate-style runtime upgrades work because a Substrate runtime *is*
sandboxed WebAssembly. This node's consensus logic is native Rust compiled
into the binary. Shipping the new logic and letting governance choose when it
activates is the honest version of the same idea, and it is what the table
above does.

\textbf{The wasmtime engine.} \texttt{maya\_vm::config} pins an exact version and says
why: the fuel schedule is not stable across releases, so two nodes on two
versions can disagree about whether a call ran out of gas. Governance can
move the VM's *limits*; moving its engine is a release.

\subsubsection{Every value is checked twice}

Once when the proposal is made, and again when it executes. A release
between those two moments could tighten a range, and a value that was legal
when proposed must not become law after it stopped being legal.

\subsection{proofs.rs}
Kani harnesses for the governance rules.

Run them with:

\begin{verbatim}
cargo kani -p maya-governance
\end{verbatim}

\subsubsection{What is worth proving here}

Not that a vote is counted — a test does that better. What a model checker is
for is the claim that \textbf{no input reaches a state the rules forbid}, and that
is exactly the shape of every property below. A governance system is only as
good as its worst reachable configuration, and the worst one is by definition
the one nobody wrote a test for.

\subsubsection{Bounded versus unbounded}

Stated plainly, because the difference is the difference between a proof and
a very good test:

- Every harness here is \textbf{unbounded} over the inputs it quantifies. There
  are no loops in the code they cover, so there is no unwind bound and no
  sequence length to assume.
- What they do *not* cover is a sequence of operations. Each proves a
  property of one transition; the composition of many transitions is covered
  by the tests in \texttt{proposal.rs} and \texttt{tests/governance\_tests.rs}. Said here
  rather than left to be inferred from the absence.

\subsection{proposal.rs}
The proposal lifecycle.

\begin{verbatim}
           ┌──────────► Cancelled          (proposer withdraws, before close)
           │
 Voting ───┼──────────► Rejected           (quorum missed, or approval missed)
           │
           └──► Queued ─┬──► Executed      (timelock elapsed, grace not)
                        └──► Expired       (grace elapsed, nobody executed)
\end{verbatim}

\subsubsection{Every transition is a function of height}

Nothing here happens because time passed; it happens because a block at some
height asked. Two nodes replaying the same chain therefore make the same
transitions in the same blocks, which is the only property that matters —
a proposal that executed on one node and not another is a fork.

\subsubsection{Why `Queued` is a state and not a delay}

A passed proposal does not execute. It becomes executable, later. The gap is
the exit window, and its purpose is to give everyone who disagrees time to
act before the change binds them: sell, withdraw, fork, or simply decline to
run the release. \texttt{MIN\_TIMELOCK\_BLOCKS} is a floor rather than a default for
exactly that reason — a governance system whose timelock can reach zero has
no dissent mechanism at all.

\subsubsection{Why `Expired` exists}

A proposal that passed and was never executed must not stay executable
forever. Otherwise a change nobody remembers wanting lands years later into
a chain it was never argued about.

\subsection{tally.rs}
Counting votes.

\subsubsection{Two thresholds, measured against different things}

\textbf{Quorum} is turnout against eligible stake: did enough of the chain show
up for the result to mean anything. \textbf{Approval} is \texttt{for} against
\texttt{for + against}: of those who expressed a preference, did a majority want it.

Abstentions count toward quorum and not toward approval, which is what makes
abstaining a distinct act rather than a slower way of voting against. A
holder who thinks a question should be settled but has no view can say so.

\subsubsection{Why quorum is measured against locked stake alone}

Voting weight is locked stake plus recent work credit, but the quorum
denominator is locked stake only. Work credit decays per address at a
different moment for each address, so a chain-wide total of it would drift
from the sum of its parts — and a quorum denominator that is quietly wrong
is worse than one that is conservatively simple.

The consequence is that turnout can exceed 100\% of the denominator when
miners vote. That only makes quorum easier to reach, never harder, and the
binding threshold is approval, which is unaffected.

\subsubsection{Overflow}

Every weight is capped at [\texttt{MAX\_WEIGHT}] and every accumulation is checked.
A tally that wrapped would not merely be wrong, it would be wrong in the
direction of whoever arranged the wrap.

\section{Threshold-encrypted mempool}
\noindent\texttt{crates/mev\textbackslash{}src}

\subsection{cipher.rs}
Sealing a transaction to a committee, and opening it once the committee agrees.

\subsubsection{Why hybrid, and not "encrypt the transaction as a curve point"}

Textbook ElGamal encrypts a group element. A transaction is a few hundred
bytes of structure, so encrypting it directly would mean an encoding from
byte strings to curve points, which is either lossy, variable-time, or
both — and would cap the plaintext at 32 bytes anyway.

So this is a KEM/DEM construction. The ElGamal half transports a shared
point; a hash of that point keys ChaCha20-Poly1305, which carries the actual
bytes. The threshold property lives entirely in the KEM half, and the DEM
half is an off-the-shelf AEAD already in this tree.

\subsubsection{Why the nonce is derived, not stored}

The AEAD nonce comes out of the same KDF as the key, seeded by the ephemeral
scalar. A fresh scalar per ciphertext means a fresh nonce per ciphertext,
structurally — there is no field an encoder could fail to randomize and no
counter to get out of step. Nonce reuse under ChaCha20-Poly1305 is
catastrophic and total; the cheapest defence is to make it unrepresentable.

\subsubsection{What the associated data is for}

The AEAD's associated data binds a ciphertext to the height it was sealed
for. Without it, a ciphertext observed in one block could be replayed into
another and the committee would dutifully decrypt it there. The transaction
inside is still nonce-protected, so a replay could not double-spend — but it
could resurrect a trade at a price its sender never agreed to, which is the
same class of harm this whole crate exists to prevent.

\subsection{committee.rs}
Who can decrypt, and how many of them it takes.

\subsection{error.rs}
Why a sealed transaction could not be built, opened, or believed.

\subsection{lib.rs}
A mempool a miner can order but cannot read.

\subsubsection{The problem}

A miner choosing the order of a block sees every transaction's contents
first. That is not a bug in any one implementation; it is what "the miner
assembles the block" means. It is also the whole supply of maximal
extractable value: front-running a large swap, sandwiching it on both sides,
reordering a liquidation to win the collateral. \texttt{dex}'s uniform-price batch
clearing removes the *profit* from reordering swaps within a pair, which is
a large share of the problem and not all of it — a miner who reads a
transaction can still delay it, censor it, or race it with one of their own.

\subsubsection{The approach}

Encrypt each transaction to a committee key with a \texttt{t}-of-\texttt{n} threshold. A
miner includes ciphertext, so the order is fixed while the contents are
unknown. After the block commits, \texttt{t} committee members publish decryption
shares, anyone recombines them, and the transactions execute in the order
that was already settled.

The property this buys, stated exactly: **ordering is committed before
contents are known.** It is not confidentiality — everything is public a
block later — and it is not censorship resistance, since a miner can still
drop a ciphertext it cannot read. It removes content-dependent ordering,
which is the specific power MEV extraction runs on.

\subsubsection{What it costs, in the order the costs bite}

\textbf{Liveness depends on the committee.} If fewer than \texttt{t} members publish
shares, sealed transactions cannot be opened. This crate's answer is that
they then \textbf{expire unopened} — the nonce does not advance, nothing moves,
the sender resubmits.

The tempting alternative is a plaintext fallback: after some deadline, let
the sender reveal in the clear and execute anyway. That is worse, and the
reason is worth being precise about. A miner who wants to read a transaction
would simply not include the shares, wait out the deadline, and read it —
turning the fallback into a switch the attacker controls. A scheme whose
safety property can be waited out does not have that property. Expiry fails
toward *nothing happening*, which is recoverable; the fallback fails toward
*the attack succeeding*, which is not.

Sealing is opt-in besides. A sender who needs guaranteed inclusion more than
they need protection sends an ordinary transaction and always could.

\textbf{The setup is trusted.} [\texttt{Committee::generate}] holds the whole key for
the length of one function call. Distributed key generation removes that and
is not built — see the method's own documentation for what a compromised
setup does and does not cost.

\textbf{It is a classical primitive.} ristretto255 falls to a quantum adversary,
on a chain whose signatures and transport are post-quantum precisely because
that adversary is assumed to arrive. This is a deliberate, bounded
exception: the confidentiality here has a lifetime of roughly one block, and
an adversary able to run Shor's algorithm against a ciphertext within one
block interval is an adversary who has already broken far more urgent things
elsewhere. Post-quantum threshold encryption exists in the literature and
not in a standardized, audited Rust implementation; when one lands, this is
the crate to replace. The rest of the chain does not depend on this
assumption — a break here costs MEV protection and nothing else.

\subsubsection{Layout}

- [\texttt{shamir}] — splitting a scalar so that any \texttt{t} of \texttt{n} pieces rebuild it.
- [\texttt{committee}] — the public key, the member keys, and the setup.
- [\texttt{cipher}] — sealing, decryption shares and their proofs, and combining.

\subsection{shamir.rs}
Splitting one scalar into \texttt{n} pieces of which any \texttt{t} suffice.

A degree-\texttt{t-1} polynomial over the scalar field is fixed by any \texttt{t} of its
points and completely undetermined by any \texttt{t-1}. Put the secret at \texttt{f(0)},
hand member \texttt{i} the point \texttt{(i, f(i))}, and both halves of that sentence
become the security property: \texttt{t} members can reconstruct, \texttt{t-1} learn
nothing — not "learn less", nothing, in the information-theoretic sense.

\subsubsection{Index zero is not a member}

\texttt{f(0)} *is* the secret. A member issued index \texttt{0} would be issued the whole
key, so [\texttt{split}] numbers members from \texttt{1} and [\texttt{lagrange\_at\_zero}] refuses
an index of \texttt{0}. This is the kind of off-by-one that does not produce a
wrong answer — it produces a correct answer to the wrong question.

\subsubsection{Reconstruction happens in the exponent}

Nothing in this crate ever calls a function that returns the reassembled
secret scalar, and none exists. Members interpolate their *decryption
shares* — points on the curve — so the committee key is never materialized
anywhere after the setup that created it. See [\texttt{crate::cipher::combine}].

\section{Payment channels}
\noindent\texttt{crates/l2-flash\textbackslash{}src}

\subsection{channel.rs}
Bidirectionally signed payment channels.

\subsubsection{What makes a state binding}

A channel state is only enforceable once \textbf{both} parties have signed it.
One signature is a proposal; two are an agreement. Every update therefore
carries two signatures over identical bytes, and [\texttt{SignedState::verify}]
checks both against the registered participants.

\subsubsection{Revocation}

Advancing from state *N* to *N+1* requires revealing the revocation secret
for *N*. Holding a counterparty's revoked secret is what makes fraud
punishable: if they later publish state *N* on chain, the secret proves it
was superseded and the whole channel balance is forfeit to the victim.

Only the *commitment* ever goes on chain — a hash per state. Publishing the
secrets themselves would grow on-chain storage without bound.

\subsubsection{Two signing forms, deliberately}

- \textbf{Update form} commits to \texttt{(seq, balances, htlc\_root)} and is what parties
  exchange off-chain while HTLCs are in flight.
- \textbf{Settlement form} commits to \texttt{(seq, balances)} only, and is the format
  the L1 understands.

They are different domains, so an update signature can never be replayed as a
settlement. A state is only settleable once its HTLC set is empty; otherwise
settling would silently discard in-flight value, which is why
[\texttt{Channel::settlement\_closure}] refuses.

\subsection{error.rs}
Errors produced by the Flash layer.

\subsection{htlc.rs}
Hash Time-Locked Contracts.

An HTLC conditionally reserves value inside a channel: the receiver claims it
by revealing a preimage before an expiry height, otherwise the sender
reclaims it afterwards. Chaining one across several channels with the *same*
hash lock is what makes a multi-hop payment atomic — revealing the preimage
to claim the last hop necessarily reveals it to every upstream hop.

\subsubsection{Expiry must decrease along the route}

Each hop's expiry is strictly earlier than the hop before it. If they were
equal, an intermediary could be left having paid downstream while its own
incoming HTLC expired in the same block — losing the amount outright. The
decreasing ladder guarantees every node has a window to claim upstream after
being claimed downstream.

\subsection{lib.rs}
\subsubsection{Maya Flash}

A state-channel network for micro-payments over the Maya2C L1.

Two parties lock capacity into a channel on chain, then exchange
bidirectionally signed states off chain at no per-payment cost. Only the
opening and the final settlement touch the base chain, so a channel carrying
ten thousand payments still costs two on-chain transactions.

\subsubsection{Modules}

- [\texttt{channel}] — bidirectionally signed state, revocation, settlement
- [\texttt{htlc}] — hash time-locked contracts, the primitive behind atomic
  multi-hop payments
- [\texttt{routing}] — the directed capacity graph and its pathfinder
- [\texttt{payment}] — assembling an HTLC chain across a route

\subsubsection{Trust model}

A channel participant never has to trust its counterparty, only the chain.
Every state is signed by both sides; advancing revokes the previous one. If a
counterparty publishes a revoked state, the revocation secret proves the
fraud and forfeits their entire channel balance. The security of the whole
construction rests on a defrauded party being able to get that proof on chain
inside the dispute window — which is why the window is a channel parameter
rather than a constant.

\subsubsection{What this layer does not do}

There is no gossip protocol here: [\texttt{routing::ChannelGraph}] is a local view a
caller populates. Real networks discover topology over the wire, and a router
is only as good as the freshness of the capacities it was told about.

\subsection{payment.rs}
Assembling an atomic multi-hop payment from a route.

Every hop carries an HTLC locked to the \textbf{same} hash. That single shared
lock is what makes the payment atomic: the only way for the final recipient
to claim is to reveal the preimage, and once revealed, every upstream hop can
claim with it too. There is no state in which the recipient is paid but an
intermediary is not.

\subsubsection{The expiry ladder}

Expiries decrease along the route, by [\texttt{HOP\_EXPIRY\_DELTA}] per hop. An
intermediary's incoming HTLC must outlive its outgoing one, or it could be
claimed downstream after its own claim window upstream had already shut —
paying out with no way to be paid. The ladder is what removes that exposure.

\subsection{routing.rs}
Multi-hop routing across the channel network.

\subsubsection{On the graph's shape}

A payment-channel network is \textbf{not} a DAG. It is an undirected multigraph,
densely cyclic — any two nodes may hold several channels, and cycles are the
norm rather than the exception.

What *is* acyclic is any individual route. This module models each channel as
\textbf{two directed edges with independent capacity} — what A can push to B is
not what B can push to A — and searches with no node revisited, so every
route it returns is a simple, acyclic path. That is the accurate description;
calling the search itself a "DAG pathfinder" would misdescribe both the input
and the algorithm.

\subsubsection{Algorithm}

Dijkstra over the directed capacity graph, minimising total fee. Edges whose
directional capacity cannot carry the amount are skipped rather than
penalised: a channel that cannot forward the payment is not a worse route, it
is not a route at all.

Fees accumulate *backwards* — an intermediary forwards less than it receives,
keeping the difference — so the search runs from the destination outward and
grows the amount as it goes. Searching forward would require knowing the
final amount before the fees that determine it.

\section{SPV light client}
\noindent\texttt{crates/light-client\textbackslash{}src}

\subsection{chain.rs}
Following headers, and believing proofs against them.

\subsubsection{Fork choice is on work, never on length}

A thousand headers mined at trivial difficulty are cheaper than ten mined at
real difficulty, so a client that preferred the longer chain would prefer the
cheaper lie. [\texttt{HeaderChain}] accumulates \texttt{work\_from\_target} — the node's own
function, not a reimplementation — and the tip is whichever branch has the
most of it.

Ties keep the incumbent, matching \texttt{consensus::chain}: switching on equal work
would let anyone with a header at the current height flip the tip for free.

\subsubsection{Why the difficulty floor is checked here}

Without it, none of the above matters. An attacker mines at
\texttt{target = 0xFF..FF}, where every digest qualifies, and produces an arbitrarily
long chain in an instant. The floor is what makes each header cost something,
and it is the single check that separates SPV from credulity.

It is deliberately *not* a full difficulty-retarget validation. Recomputing
the retarget would need every header in the window, which a client syncing
from a checkpoint may not have. The floor is the weaker, cheaper property,
and the gap between them is real: a chain could sit at the floor forever and
this client would follow it. Anyone who needs more should run a full node.

\subsection{error.rs}
Why a light client refused a header or a proof.

\subsection{lib.rs}
An SPV light client: headers, fork choice, and state proofs.

\subsubsection{What "light" buys and what it costs}

A full node downloads every transaction, executes it, and computes the state
root itself. It needs no trust at all — it checks everything.

A light client downloads only \textbf{headers}, which are 144 bytes each, and
verifies two things:

1. Each header's proof of work is valid, and the chain it is on has more
   cumulative work than any competitor. This is the same fork-choice rule
   \texttt{consensus::chain} uses, and it is not weakened here.
2. A [\texttt{custom\_l1\_node::state::AccountProof}] reproduces the \texttt{state\_root} a
   header commits to.

What it gives up is \textbf{validity}. A light client cannot tell whether the
transactions that produced a state root were themselves legal. If a majority
of hash power produces a chain in which somebody minted coins from nothing,
every header is valid, every proof verifies, and the light client believes
it. That is the SPV trade-off in one sentence, and no amount of engineering
removes it — it is why this crate is a convenience and not a substitute.

\subsubsection{What it does not give up}

Three things a naive SPV implementation gets wrong, and this one does not:

- \textbf{Difficulty is checked, not assumed.} A header claiming an easy target
  is cheap to produce. [\texttt{HeaderChain}] refuses any target easier than the
  network's floor, so a fake chain has to do real work.
- \textbf{Work, not length, decides.} A thousand easy headers must not outrank ten
  hard ones. Fork choice is on cumulative work, computed with the node's own
  \texttt{work\_from\_target}.
- \textbf{A proof is bound to a root, and a root to a header.} Verifying a proof
  against a root nobody mined proves nothing, so [\texttt{LightClient::verify\_account}]
  takes a *height*, looks the header up, and uses its committed root.

\subsubsection{What this is not}

It is not a bridge to another chain. "Bridge" in the milestone list means the
client can be embedded in something that needs to check Maya2C state without
running a node — a wallet, an exchange's deposit watcher, a rollup's
settlement checker. The cross-chain half of a bridge is a message-passing
protocol on top of this, and it is not built here.

\subsection{stateless.rs}
\begin{framed}\noindent\textbf{Status.} This module is a research branch. It is present in the tree and no consensus path reaches it. Nothing described here is part of the running chain.\end{framed}

Checking blocks, not only headers. Research branch.

[\texttt{LightClient}] follows headers and believes that whatever produced a state
root was legal — the SPV trade-off in \texttt{lib.rs}. For transfer blocks a
witness removes that belief: [\texttt{StatelessValidator}] re-runs each block's
transfers against its witness and checks the declared root, holding 32 bytes
of state between blocks and one block's witness while it works.

What it still does not do:

- \textbf{Proof of work.} That is [\texttt{crate::HeaderChain}]'s job, and its DAG cache
  is tens of megabytes. [\texttt{LightClient::verify\_block}] joins the two: a
  work-checked header, and a block checked against it.
- \textbf{Anything but transfers,} or any chain whose oracle, governance or
  sealed mempool holds records. Those answer
  [\texttt{LightClientError::BlockUnverifiable}], which is not a verdict — see
  \texttt{custom\_l1\_node::state::stateless}.
- \textbf{Fetch witnesses.} No gossip topic carries them yet.

\section{Pool protocol}
\noindent\texttt{crates/stratum-v2\textbackslash{}src}

\subsection{codec.rs}
Bounds-checked reader and writer for Stratum V2's primitive types.

Deliberately the same shape as \texttt{src/core/codec.rs}: every read is range
checked, every collection length is capped before anything is reserved, and
[\texttt{Reader::finish}] refuses trailing bytes so one value cannot have two
encodings. The node's reader could not simply be reused because it returns
\texttt{NodeError}, and this crate does not depend on the chain — but the
discipline is copied on purpose, not reinvented.

\subsubsection{Endianness}

Stratum V2 is little-endian throughout, which happens to match the node's
wire format (\texttt{src/core/block.rs:33-41}). Integers here are therefore LE.

\textbf{Except targets and hashes.} See [\texttt{Reader::read\_bytes32}].

\subsection{error.rs}
Errors raised while framing or parsing Stratum V2 messages.

Every variant here is reachable from bytes a peer chose, so each one carries
enough detail to tell an operator *which* peer sent *what* nonsense without
having to reproduce it. None of them is a panic: a pool that aborts on a
malformed frame is a pool one connection can take down.

\subsection{frame.rs}
The Stratum V2 frame header, unchanged from the specification.

\begin{verbatim}
┌────────────────┬──────────┬────────────┬─────────────────┐
│ extension_type │ msg_type │ msg_length │ payload         │
│ U16            │ U8       │ U24        │ msg_length bytes│
└────────────────┴──────────┴────────────┴─────────────────┘
  bit 15 of extension_type is the channel_msg flag
\end{verbatim}

Six bytes, little-endian, exactly as SV2 specifies. This layer is the one
part of the protocol Maya2C adopts without alteration — nothing in it
mentions a Bitcoin header, so there is nothing to adapt.

\subsubsection{The extension namespace}

Extension type \texttt{0x0000} is the reserved SV2 mining protocol, and this is not
that: Maya2C's mining messages carry a \texttt{state\_root}, a \texttt{tx\_root} where SV2
puts its merkle root, and a 256-bit target where SV2 has \texttt{nbits}. Using \texttt{0x0000} would announce
wire compatibility that does not exist, and the failure mode — a stock SV2
client parsing our \texttt{NewMiningJob} as its own — is a client mining garbage
rather than a client reporting an error.

So Maya2C's messages live under [\texttt{MAYA\_EXTENSION\_TYPE}], \texttt{0x4D41}, which is
\texttt{"MA"} in ASCII and legible in a hex dump.

\subsection{lib.rs}
Stratum V2 framing and mining messages for Maya2C.

\subsubsection{What this is, and what it deliberately is not}

This is Stratum V2's *architecture* — binary framing, a channel abstraction,
per-channel targets, batched share acknowledgement, and the separation of
template choice from pool operation — carrying a mining sub-protocol shaped
around Maya2C's header.

It is \textbf{not} wire-compatible with stock Stratum V2 clients, and it does not
pretend to be. SV2's mining messages assume a Bitcoin header: they carry a
\texttt{merkle\_root}, a compact \texttt{nbits}, a rollable \texttt{version}, and a coinbase to
hold an extranonce. Maya2C's 144-byte header (\texttt{src/core/block.rs}) keeps
only the root, as \texttt{tx\_root}. No SRI or Braiins client could mine this chain even against a
byte-perfect implementation of the specification — it would need an
ArgonBlake hasher and a Maya2C header builder, at which point it is a
different client.

Given that, announcing compatibility would have been the harmful choice. So
Maya2C's messages live under their own extension id
([\texttt{frame::MAYA\_EXTENSION\_TYPE}]) and a stock client's frames are refused with
[\texttt{error::Sv2Error::UnknownExtension}] rather than misparsed. Message
*numbering* still tracks the specification slot for slot, so the
correspondence stays legible; [\texttt{messages::mining}] tabulates every field that
had to change and why.

\subsubsection{No chain dependency}

Nothing here depends on \texttt{custom-l1-node}. Every byte this crate reads was
chosen by a stranger, so it is built to be fuzzed and reviewed in isolation,
the way \texttt{ledger-math} is built to be model-checked in isolation. Protocol
values meet chain types one layer up, in the pool daemon.

\subsubsection{Decoding discipline}

Copied from \texttt{src/core/codec.rs} rather than reinvented:

- every read is range-checked, so a truncated frame is an error and never a
  panic — a pool that aborts on a malformed frame is a pool that one
  connection can take down;
- a declared payload length above [\texttt{frame::MAX\_PAYLOAD\_LEN}] is refused
  before anything is reserved, because with 50,000 connections a length
  field is an allocation primitive;
- [\texttt{codec::Reader::finish}] rejects trailing bytes, so one message cannot
  have two encodings. For a share submission that matters concretely: two
  encodings would be two share identities for one piece of work, and the
  duplicate check keys on decoded fields.

\subsection{messages/common.rs}
Connection setup, shared by every sub-protocol.

These three messages are the part of Stratum V2 that needed no adaptation:
they negotiate versions and describe the peer, and say nothing about what a
block header looks like. Field order and message numbering follow the
specification exactly.

\subsection{messages/mining.rs}
The mining sub-protocol, adapted to Maya2C's header.

\subsubsection{What changed from the specification, and why}

Stratum V2's mining messages are shaped around a Bitcoin header. Maya2C's is
144 bytes of `prev\_hash ‖ state\_root ‖ timestamp ‖ nonce ‖ difficulty\_target
‖ tx\_root\texttt{ (}src/core/block.rs\texttt{) — no version field, no compact }nbits`, and
no coinbase to hide an extranonce in. Four adaptations follow.
Message numbering still matches SV2 slot for slot, so the correspondence
stays readable.

\begin{verbatim}
| SV2 field | Here | Why |
|---|---|---|
| `merkle_root` in `NewMiningJob` | `tx_root` **and** `state_root` | The header commits to the transaction tree, as in SV2, and also to the post-execution state, which the chain checks |
| `nbits` (U32 compact) in `SetNewPrevHash` | `target` (32 bytes) | Maya2C targets are full 256-bit values compared bytewise (`src/crypto/pow.rs:16-18`); there is no compact form to pack into |
| `version` in jobs and shares | *absent* | The header has no version field to roll |
| `SetExtranoncePrefix` | [`SetNonceRange`] | No coinbase means no extranonce; see below |
| `ntime` (U32) | `timestamp` (U64) | The header's timestamp is a `u64` |
\end{verbatim}

\subsubsection{Nonce ranges instead of extranonces}

In Bitcoin, each miner gets a distinct search space by varying the coinbase
extranonce. Maya2C has no coinbase, so the only field a miner may vary is
the 8-byte \texttt{nonce} — and with tens of thousands of connections on one job,
every one of them would otherwise start at zero and walk the same path.

Each channel is therefore assigned a disjoint half-open nonce range at open
time, the way \texttt{src/consensus/miner.rs:6-10} already partitions across
threads. Two channels cannot collide by construction rather than by luck,
duplicate detection becomes exact, and a submitted nonce outside a channel's
range is rejected without hashing anything — which, at 25 ms per ArgonBlake
verification, is the cheapest rejection the pool has.

\subsection{messages/mod.rs}
Every message this build speaks, and the mapping between message types and
frames.

Message numbers follow the Stratum V2 specification slot for slot, even
where a field set had to change (see [\texttt{mining}]). Keeping the numbering
means anyone holding the spec can read a hex dump of this protocol; renaming
the slots would have bought nothing and cost that.

Job Declaration messages (\texttt{0x50}–\texttt{0x5b}) are deliberately absent. They are
Phase 6 work, and an empty variant that decodes into nothing is worse than
an honest [\texttt{Sv2Error::UnknownMessageType}].

\section{Pool daemon}
\noindent\texttt{bins/pool-service\textbackslash{}src}

\subsection{api.rs}
The JSON API.

Read-only, and deliberately so. Everything that changes pool state arrives
over the authenticated Stratum port or is decided by the daemon's own loops.
An HTTP endpoint that could retarget a channel or trigger a payout would be
a second, weaker way into the parts of the pool that handle money.

\subsubsection{Shapes are the ones the dashboard renders}

The same types [\texttt{crate::ui}] uses, serialised. Two representations of a rig
would be two things to keep in step, and the one that drifts is always the
one without a page rendering it.

\subsubsection{Reported figures keep their prefix in JSON too}

\texttt{reported\_hashrate}, \texttt{reported\_power\_milliwatts}, and the rest are named
that way in the wire format, not only on the page. A consumer writing an
alert against this API needs to know which numbers the pool measured and
which a rig asserted.

\subsection{bin/pool.rs}
\texttt{pool} — the Maya2C mining pool daemon.

\begin{verbatim}
pool --node-rpc http://127.0.0.1:8545 \
     --treasury treasury.key \
     --reward-per-block 1000000000 \
     --fee-address <HEX>
\end{verbatim}

\subsubsection{Three listeners, three audiences}

The Stratum port is for miners and speaks the post-quantum transport. The
dashboard is for miners and operators and speaks HTTP. The exporter is for
Prometheus and is off unless asked for, because what it serves is
operational data that belongs inside the pod network — the same reasoning
\texttt{bins/maya2c-node/src/main.rs} applies to the node's own exporter.

\subsubsection{The password is not a flag}

\texttt{MAYA\_POOL\_TREASURY\_PASSWORD}, matching the wallet. A password on a command
line lands in shell history and is visible in the process list; an
environment variable is neither by default. See [\texttt{treasury::PASSWORD\_ENV}].

\subsubsection{It refuses to start rather than pay wrongly}

\texttt{--reward-per-block} has no default. This chain mints no block reward
(\texttt{docs/stratum-v2.md} §5), so what a found block distributes is operator
policy funded from a treasury account, and a plausible-looking default would
be a policy invented by the binary rather than chosen by the operator.

\subsection{channel.rs}
Mining channels: nonce ranges, per-channel state, and the registry.

\subsubsection{Disjoint nonce ranges are the whole design}

Bitcoin gives each miner its own search space through the coinbase
extranonce. Maya2C has no coinbase, so the only field a miner may vary is
the header's 8-byte nonce (\texttt{src/core/block.rs:23}). With tens of thousands
of channels on one job, every one of them would otherwise start at zero and
walk the same path — the pool would pay many miners for the same hash.

Each channel is therefore assigned a disjoint half-open range at open time,
the way \texttt{src/consensus/miner.rs} already partitions across threads. Three
things follow, and all three are load-bearing:

1. Two channels cannot collide by construction rather than by luck.
2. Duplicate detection is exact, and its memory is bounded by the range.
3. A nonce outside a channel's range is refused \textbf{without hashing} — at
   25 ms a verification, the cheapest rejection the pool has.

\subsubsection{Ranges are recycled, and that is safe}

[\texttt{NonceAllocator}] returns a closed channel's range to a free list. Without
recycling, a pool with churn exhausts its ranges after a few million channel
opens and starts refusing connections for no reason a miner can see. It is
safe because duplicate detection is keyed by \texttt{(job, nonce)} *within* a
channel, and a recycled range belongs to a new channel with its own empty
set — there is no cross-channel duplicate to miss.

\subsection{config.rs}
Pool configuration, and the checks that run before anything binds.

\subsubsection{Why validation is aggressive here}

Most of these values are policy about other people's money. A pool that
starts with \texttt{fee\_rate = 1.5} does not fail — it quietly takes 150\% of every
reward and pays miners nothing, and the first report of the problem comes
from a miner rather than from the process. So every field with a meaningful
range is range-checked at load, and [\texttt{PoolConfig::validate}] refuses rather
than clamps.

\subsubsection{The treasury is the part with no default}

\texttt{docs/stratum-v2.md} §5 is explicit that this chain has no block reward:
\texttt{apply\_block\_checked} credits no subsidy and fees burn to \texttt{FEE\_SINK}. PPLNS
therefore accrues *share weight*, and coin only enters at settlement, out of
an account the operator funds. There is no defensible default for whose
account that is or what a block is worth, so both are required and the
daemon refuses to start without them.

\subsection{daemon.rs}
The daemon: the Stratum listener, and the loops around it.

\subsubsection{Transport}

The pool reuses \texttt{src/network/pq} — Noise-style XX over X25519 with an
ML-KEM-768 exchange layered inside, ChaCha20Poly1305, 64 KiB frames.
\texttt{docs/stratum-v2.md} §6 gives the reasoning: SV2 specifies
\texttt{Noise\_NX\_secp256k1\_…}, and adopting it would introduce secp256k1 — the one
classical primitive this codebase deliberately avoids — to buy interop that
is unreachable anyway.

\texttt{pq::handshake::respond} is generic over *futures* \texttt{AsyncRead}/\texttt{AsyncWrite}
and a Tokio \texttt{TcpStream} implements Tokio's, which is the entire reason
\texttt{tokio-util}'s compat layer is a dependency.

\subsubsection{One task per connection, and no `select!` over the reader}

Job pushes arrive on their own task and leave through an outbound queue, so
the read loop never sits in a \texttt{select!} alongside a partially consumed
frame. A cancelled frame read would discard bytes already taken off the
socket and desynchronise the stream — a failure that appears as "this one
miner sends garbage" long after the cause.

\subsubsection{Verification happens inline, per connection}

The read loop awaits the validator rather than spawning. That serialises one
connection's shares, which is the correct backpressure: vardiff is holding
each channel near one share a minute, so a rig waiting on its own share is
waiting on work it just asked for. Spawning per share would let one
connection queue thousands of verifications and starve every other.

\subsection{error.rs}
Pool errors.

Split the way \texttt{apps/explorer/src/error.rs} splits its own: by the operational
response, not by the module that raised it. An on-call engineer reading
\texttt{ledger:} reaches for the disk, and one reading \texttt{node rpc:} reaches for the
chain. A single \texttt{Internal(String)} would have told them neither.

[\texttt{PoolError::Treasury}] is the variant worth pausing on. It covers every way
the payout path can refuse to spend — an unfunded account, a locked
keystore, a spend cap reached — and all of them mean *stop paying*, never
*retry harder*. Folding them into a generic failure would let a retry loop
treat "the treasury is empty" as a transient condition.

\subsection{job.rs}
Work templates, and the ids miners submit against.

\subsubsection{Jobs go stale fast here}

\texttt{TARGET\_BLOCK\_TIME} is fifteen seconds (\texttt{src/consensus/difficulty.rs:32}).
A ten-minute chain can afford to measure its stale-grace window in seconds;
this one cannot. Job push has to be sub-second and the grace window is one
job deep — see [\texttt{crate::channel::ChannelState::current\_job}] for why one and
not zero, and not three.

\subsubsection{The pool builds empty blocks}

\texttt{get\_mining\_candidate} returns a header over the current state root and no
transaction list (\texttt{src/rpc/server.rs:147}), and \texttt{submit\_block} takes a whole
block. The pool therefore submits \texttt{Block::new(header, vec![])}, which is what
the node's own miner does. This is not a limitation being worked around: on
a chain that burns fees to \texttt{FEE\_SINK} there is no fee revenue to collect by
including transactions, so template construction is the node's business and
not the pool's.

\subsection{ledger/memory.rs}
In-memory ledger, for tests.

Proves the logic — window walks, maturity, reversal, batch lifecycle — and
proves nothing about durability. Anything that matters about surviving a
crash is only demonstrated against [\texttt{super::rocks::RocksLedger}].

One \texttt{Mutex} around one struct, deliberately. Finer-grained locking would buy
throughput this type does not need and would let a test observe a state no
ordering of the real store can produce.

\subsection{ledger/mod.rs}
The share ledger, behind a trait.

Two implementations: [\texttt{rocks::RocksLedger}] for deployment and
[\texttt{memory::MemoryLedger}] for tests. The trait exists so the PPLNS window
walk, the maturity rules, and the reorg reversal path can be tested without
a database — the same reasoning \texttt{apps/explorer/src/store/mod.rs} gives for its
own split, and the same honest limit applies: \texttt{MemoryLedger} proves the
logic and cannot prove the durability.

\subsubsection{Why RocksDB and not the explorer's PostgreSQL}

The two stores hold different kinds of data. Everything the explorer indexes
can be rebuilt by re-reading the chain; a share credit exists nowhere but
here, and a lost one is a miner who worked for nothing. That argues for the
storage engine the operator already runs and already knows how to back up,
which on this project is RocksDB.

The cost is stated rather than hidden: one node, no replication, and a
restore is a file restore.

\subsubsection{The trait is synchronous}

RocksDB is a blocking API, and wrapping it in \texttt{async\_trait} would produce
futures that block their executor thread anyway — the appearance of async
without the property. \texttt{src/rpc/server.rs} already calls \texttt{StateDB} straight
from request handlers for the same reason. Every read a request handler
makes here is a point lookup or a short bounded scan; the one unbounded walk
([\texttt{ShareLedger::window}]) runs on the payout task, never on a request.

\subsubsection{Append-only, and what that buys}

Nothing rewrites a [\texttt{ShareRecord}]. Credits are derived from the records by
reading a window backwards from the tip, and a record that could be edited
is a credit that could be moved after the fact. Records are eventually
*pruned* — dropped wholesale from the old end — which is a different
operation from editing and cannot change what a retained window says.

\subsection{ledger/rocks.rs}
The durable share ledger.

\subsubsection{Key layout}

One column family, prefixed keys, matching \texttt{src/state/db.rs} rather than
inventing a second convention in the same repository:

\begin{verbatim}
s<sequence be64>      ShareRecord          the append-only share log
b<block id>           FoundBlock           blocks the pool found
c<block id>           Vec<PayoutEntry>     what that block credited, verbatim
a<address>            MinerBalance         unpaid / immature / paid
p<batch id be64>      PayoutBatch          payout batches
m:sequence            u64                  next share sequence
m:batch               u64                  next batch id
\end{verbatim}

Sequences and batch ids are stored \textbf{big-endian} so RocksDB's byte order is
numeric order. With little-endian keys, share 256 would sort before share 2
and every window walk would read the wrong shares — silently, and in a way
that only shows up as miners being paid the wrong amounts.

\subsubsection{What is atomic, and why those things}

Three operations write more than one key and each uses a single
[\texttt{WriteBatch}]:

- \textbf{recording a block} writes the block, its credit list, and every affected
  balance. A block stored without its credits pays nobody; credits stored
  without their block can never be matured or reversed.
- \textbf{maturing} and \textbf{orphaning} move balances and flip the block's state
  together, so a crash cannot leave a block marked mature whose credits are
  still immature — or worse, the reverse.
- \textbf{debiting a batch} moves every recipient at once.

The credit list is stored rather than recomputed. Recomputing it at maturity
would re-walk a ledger that has moved on, and the second answer would not
match the first.

\subsection{lib.rs}
\subsubsection{Maya2C pool service}

A mining pool: a Stratum V2 listener that validates shares against the
chain's own proof-of-work rule, an append-only ledger that records what each
share was worth, a PPLNS engine that turns those records into credits, and a
payout engine that settles them as ordinary signed transfers.

\subsubsection{Layout}

- [\texttt{config}] — policy, and the checks that run before anything binds
- [\texttt{model}] — the pool's own view of miners, shares, and payouts
- [\texttt{channel}] — open channels, disjoint nonce ranges, duplicate memory
- [\texttt{vardiff}] — per-channel difficulty, the pool's load governor
- [\texttt{job}] — work templates and the ids miners submit against
- [\texttt{share}] — verification, through the chain's rule rather than a copy
- [\texttt{ledger}] — the append-only share ledger, behind a trait
- [\texttt{pplns}] — window accounting and the split
- [\texttt{treasury}] — the key, the nonce, and the spend caps
- [\texttt{payout}] — the batch state machine that settles credits on chain
- [\texttt{telemetry}] — measured hash rate, efficiency, and reported rig health
- [\texttt{metrics}] — Prometheus instrumentation
- [\texttt{node}] — the JSON-RPC client the pool drives the chain through
- [\texttt{server}] — JSON API, dashboard routes, and the WebSocket
- [\texttt{ui}] — Leptos server-rendered views
- [\texttt{daemon}] — the Stratum listener and the loops that tie it together

\subsubsection{Three decisions worth knowing about}

\textbf{There is no block reward to split.} \texttt{docs/stratum-v2.md} §5 establishes
it: \texttt{apply\_block\_checked} credits no subsidy and fees burn to \texttt{FEE\_SINK}.
Share credits therefore accrue in *weight*, and coin enters only at
settlement, from a treasury the operator funds. [\texttt{config::PoolConfig}]
refuses to start without that decision having been made explicitly.

\textbf{Share validation is the scaling wall, not the connection count.} A share
costs 25.4 ms of Argon2id below the DAG fork. [\texttt{vardiff}] is what keeps the
aggregate rate survivable, and backpressure sheds load by raising targets —
never by dropping shares a miner already did the work for.

\textbf{Payouts are idempotent through the treasury nonce.} A batch reserves a
nonce and writes itself down *before* it is signed, and a resumed batch
rebroadcasts the same bytes rather than signing new ones. A crash at any
point produces one transaction or none, never two.

\subsection{metrics.rs}
Prometheus instrumentation for the pool.

Built on the node's registry conventions (\texttt{src/metrics/mod.rs}) and served
by the node's exporter through
[\texttt{custom\_l1\_node::metrics::server::Exposition}] — one accept loop, one
routing table, one set of decisions about what an exporter should refuse to
serve.

\subsubsection{Labels are a closed set, and here that matters more than usual}

\texttt{src/metrics/mod.rs} refuses open-ended label values because a hostile peer
could otherwise create unbounded time series and exhaust the scraper. The
pool has the same exposure through a wider door: a rejection reason derived
from error text, or a per-worker label taken from the \texttt{user\_identity} a
stranger chose, would let one connection mint series at will. Rejection
reasons are therefore a fixed \texttt{\&'static str} set, and **nothing here is
labelled per worker.** Per-rig figures belong to the dashboard and the JSON
API, which page and can be authorised; a metric labelled by rig name is a
cardinality bomb with 50,000 fuses.

\subsubsection{Measured and reported stay separate here too}

\texttt{maya\_pool\_hashrate\_hashes\_per\_second} is derived from accepted share weight.
\texttt{maya\_pool\_reported\_power\_watts} is a sum of numbers rigs sent about
themselves. An operator alerting on efficiency needs to know which is which,
and the metric names are where that gets said.

\subsection{model.rs}
The pool's own view of miners, shares, and payouts.

Everything here is \texttt{serde}, because these are the types the JSON API and the
dashboard render. Nothing here reaches consensus; the chain types stay in
\texttt{custom\_l1\_node} and meet these only in [\texttt{crate::payout}].

\subsubsection{Two identities, deliberately separate}

A \textbf{miner} is an address — the thing that gets paid. A \textbf{worker} is one
rig belonging to that miner, and exists only so an operator can tell which
of their machines stopped. Credits accrue to the miner; statistics accrue to
the worker. Keeping them apart is what makes splitting a farm across more
workers payout-neutral, which [\texttt{crate::pplns}] tests explicitly.

\subsection{node.rs}
The node, as the pool sees it.

A thin wrapper over the five JSON-RPC methods the pool actually uses. It is
its own type rather than \texttt{l1\_wallet::client::NodeClient} because that client
covers the wallet's three methods and reports failures as \texttt{anyhow} — the
pool needs \texttt{get\_mining\_candidate} and \texttt{submit\_block} as well, and needs a
typed [\texttt{PoolError::NodeRpc}] so a retry loop can tell an unreachable chain
from an unusable ledger.

\subsubsection{Every call here can fail, and the pool must keep serving when it does}

An unreachable node means work goes stale and blocks cannot be submitted. It
does \textbf{not} mean accepted shares are lost: those are already in the ledger,
and the credits they carry are safe. The error variant this module returns
is the one that says so.

\subsection{payout.rs}
Settling share credits on chain.

\subsubsection{The state machine}

\begin{verbatim}
                   treasury refuses
  Pending ──────────────────────────────► Failed   (credits returned)
     │
     │ signed against a reserved nonce
     ▼
  Signed ─── broadcast ──► Submitted ─── buried ──► Confirmed
                                ▲                │
                                └── reorg ───────┘
\end{verbatim}

Nothing moves backwards out of \texttt{Submitted} except into \texttt{Submitted} again.
Once the bytes are on the network the pool's only options are to wait or to
rebroadcast *those* bytes; signing replacement bytes for the same nonce
would be a second transaction competing with the first, and whichever won
the pool would have double-counted the outcome.

\subsubsection{How the pool knows a payout landed}

There is no \texttt{get\_transaction} RPC to ask. What there is, is
\texttt{src/state/db.rs:418}: a transaction is valid only when \texttt{tx.nonce} equals
the sender's current nonce exactly, so transactions from one account are
strictly sequential. The pool is the only spender from the treasury, so the
account's nonce advancing past a batch's reserved nonce means *that batch's
transaction* was included — there is nothing else it could have been.

A reorg that unwinds the payout moves the nonce back down, which the watcher
reads as "not landed after all" and rebroadcasts the same bytes into the
nonce slot that is free again.

\subsubsection{One batch at a time}

Sequential nonces make pipelining a trap: several batches in flight means
predicting nonces for transactions not yet accepted, and one rejection
strands every batch behind it. The engine therefore starts a new batch only
when nothing is open.

\subsection{pplns.rs}
Pay-Per-Last-N-Shares accounting.

\subsubsection{The window is measured in weight, not in shares}

A share count is the wrong unit the moment vardiff exists. Two miners on
different targets submit shares worth different amounts of work, and a
window of "the last hundred thousand shares" would pay them the same for
them. \texttt{N} here is a quantity of *work*, expressed as a multiple of network
difficulty, and a share contributes its own weight to filling it.

That is also what makes the scheme resistant to the obvious manipulation.
Splitting one rig across ten workers produces ten times the shares at a
tenth the difficulty each, for the same total weight and the same payout —
\texttt{tests::splitting\_a\_farm\_across\_workers\_changes\_nothing} states it as an
assertion rather than a claim.

\subsubsection{Why the denominator is what accumulated, not `N`}

A young pool has not yet earned \texttt{N} worth of shares. Dividing by \texttt{N} anyway
would distribute less than the reward and silently strand the difference in
the treasury. The denominator is therefore always the weight the window
actually contains — which for a full window is \texttt{N} plus at most one boundary
share, and for a short one is everything there is.

\subsubsection{Rounding, and the invariant that has to hold}

Integer division loses units. Handing the remainder to nobody would leak
value out of every payout; handing it to the largest miner would make the
split depend on rounding rather than on work. This uses the largest-remainder
method — the fractional parts are ranked and the spare units go to the
closest-missed claims, ties broken by address so two runs of the same input
produce the same output.

The invariant [\texttt{split}] enforces is exact: the credits sum to the reward, to
the unit. Not approximately, and not with a tolerance.

\subsection{server.rs}
The dashboard and JSON API server.

An Axum application over [\texttt{crate::state::PoolState}], rendering
[\texttt{crate::ui}] and serving [\texttt{crate::api}]. Structurally the explorer's
server, and for the same reasons — including the one about slow clients:
a WebSocket subscriber that falls behind reports \texttt{Lagged} and is skipped
ahead rather than waited for. Share accounting must never block on a
browser.

\subsubsection{This port is not the mining port}

Miners connect over the post-quantum Stratum transport on its own listener.
Nothing here can open a channel, move a target, or start a payout — every
route is a read. A dashboard that could act on the pool would be a second
and much weaker door into the parts that handle money.

\subsubsection{And it is not the metrics port either}

Three listeners, three audiences: miners, operators, and Prometheus. The
exporter carries per-pool operational data and stays inside the pod network
(\texttt{src/metrics/mod.rs} gives that reasoning); this one is the page a miner
opens.

\subsection{session.rs}
One connection's protocol state machine.

Deliberately free of I/O. [\texttt{Session::handle}] takes a decoded message and
returns what should happen next; the socket, the encryption, and the
validator queue all live elsewhere. That split is what makes the protocol
testable at all — every branch below can be exercised with a \texttt{Vec<Message>}
and no listener, and the 50-worker simulation drives these sessions directly
rather than through 50 sockets.

\subsubsection{What a connection is allowed to do}

- \textbf{Before setup}, only \texttt{SetupConnection}. A channel opened by a peer that
  has not negotiated a version is a channel whose protocol assumptions are
  unknown.
- \textbf{Only its own channels.} A message naming a channel this connection did
  not open is refused with \texttt{unknown-channel-id}, not looked up in the global
  registry. Without that check any connection could retarget, close, or
  submit against any other miner's channel by guessing a small integer.
- \textbf{No unsolicited work.} Jobs are pushed by the daemon, never requested.

\subsubsection{Expensive work leaves by a different door}

A share that survives the integer checks becomes a
[\texttt{SessionAction::Validate}], not a reply. Hashing on a connection task would
put 25 ms of Argon2id in front of every other message that connection has
queued, and at 50,000 connections would put it in front of the runtime.

\subsection{share.rs}
Share verification.

\subsubsection{The pool uses the chain's own rule, never a copy of it}

Everything expensive here goes through \texttt{BlockHeader::pow\_hash\_at} and
\texttt{meets\_target} from \texttt{custom\_l1\_node}. That is the point of this crate
depending on the node at all: a pool that reimplements the proof-of-work
rule is a pool that eventually credits work the chain would reject, and the
divergence shows up as blocks the pool submits and the network refuses —
after the shares behind them have been paid for.

\texttt{pow\_hash\_at} also picks the rule by height, ArgonBlake below the DAG
activation and hashimoto above it (\texttt{src/core/block.rs:142}). Hashing at the
wrong height is the same class of error as hashing with the wrong function.

\subsubsection{Order of operations is a cost decision}

One verification is 25.4 ms of Argon2id below the fork. Everything that can
be decided from integers is decided first, in
[\texttt{crate::channel::ChannelState::precheck}], and only what survives reaches
this module. Within it, the single hash answers both questions — share
target and network target — because the digest is the same one.

\subsubsection{A block is a share that got lucky}

There is no separate "solution" submission. A share whose digest also meets
the network target *is* a block, and it is credited as a share as well as
submitted as a block. Crediting it only as a block would quietly steal the
share weight from the miner who found it.

\subsection{state.rs}
State every part of the daemon shares.

\subsubsection{Lock discipline}

Three locks, and an order: \texttt{jobs} before \texttt{channels} before \texttt{telemetry}.
Nothing takes them in another order, and nothing holds one across an
\texttt{.await}. Both rules exist for the same reason — a connection task that
blocks while holding the channel registry stalls every other connection,
and at 50,000 of them that is the whole pool.

The order is worth stating rather than assuming, because the natural way to
write one particular piece of code violates it: after a vardiff retune the
dashboard wants the channel's new bit count, and reaching for the channel
registry while already holding telemetry is the obvious way to get it.
[\texttt{crate::validator}] reads it out under the channel lock instead and
releases that lock before touching telemetry.

\texttt{std::sync::Mutex} rather than Tokio's, deliberately. Every critical section
here is a map lookup and a few field updates; a Tokio mutex would add a
future and a scheduler hop to work that finishes in nanoseconds, and its one
real advantage — being held across an await — is a thing this module
forbids anyway.

\subsubsection{Events}

A broadcast channel, carrying what the dashboard needs to redraw. A slow
browser falls behind and its receiver reports \texttt{Lagged}; the handler skips
ahead rather than stalling, exactly as \texttt{apps/explorer/src/server.rs} does. Pool
accounting must never block on a client — one wedged tab would otherwise
stop shares being credited.

\subsection{telemetry.rs}
Per-worker statistics: what the pool measured, and what the rig claimed.

\subsubsection{Measured and reported are kept apart everywhere}

Two hash rates appear on every rig. One is derived from accepted share
weight and is the pool's own arithmetic; the other arrived in a message the
rig composed. Averaging them, or showing whichever is larger, would destroy
the only interesting comparison an operator can make — a rig claiming 400
MH/s and delivering 40 is a rig with a problem, and the dashboard's job is
to make that visible rather than to smooth it away.

Power, temperature, and fan speed have no measured counterpart at all. They
are carried under \texttt{reported\_} names and are structurally unable to reach
[\texttt{crate::pplns}]: nothing in this module is an input to a credit.

\subsubsection{Measuring hash rate from share weight}

A share worth \texttt{2\textasciicircum{}b} of weight is evidence of about \texttt{2\textasciicircum{}b} hashes. Dividing
accumulated weight by elapsed time gives hashes per second directly, with no
need to trust anything the rig said.

The estimate is smoothed over a sliding window rather than computed per
share. Share arrival is Poisson: an instantaneous rate from one interval has
a variance as large as its mean, so a per-share figure would swing by
multiples between consecutive shares and read as a broken rig.

\subsection{treasury.rs}
The treasury: the key that signs payouts, and the limits on what it may do.

\subsubsection{Where the money comes from}

Not from a block reward — this chain has none (\texttt{docs/stratum-v2.md} §5).
Miners are paid out of an ordinary account the operator funds, by ordinary
signed transfers that need no consensus change. Everything the pool computes
upstream of here is denominated in share weight; this is where weight
becomes coin.

\subsubsection{The key is hot, and the limits exist because of that}

A daemon that signs payouts is a daemon holding a spending key on a machine
with a public mining port. That is a real exposure and it is not resolved by
being careful. What is here instead is a bound on the damage:

- [\texttt{PoolConfig::max\_batch\_value}] caps what one batch can move;
- the balance floor check refuses to sign against an account that cannot
  cover the batch, so a drained treasury fails loudly rather than emitting
  transactions the mempool will reject;
- the password never appears in an argument, matching the wallet's own
  reasoning (\texttt{bins/l1-wallet/src/main.rs}): a password on a command line is in shell
  history and in the process list.

Moving the signer out of the daemon entirely is the real fix and is out of
scope here. It is named in \texttt{docs/pool-service.md} rather than left implied.

\subsubsection{Nonces are the idempotency mechanism}

\texttt{src/state/db.rs:418} requires \texttt{tx.nonce == sender.nonce} exactly, so
transactions from one account are strictly sequential and a nonce is a slot
that exists once. A batch reserves its slot before it is signed and keeps
it across restarts, which is what makes "crashed after signing, before
broadcasting" recoverable: the resumed daemon rebroadcasts the same bytes
rather than signing a second transaction for the same slot.

Consequently \textbf{one batch is in flight at a time}. Pipelining several would
mean predicting nonces for transactions not yet accepted, and a single
rejection in the middle would strand every batch behind it.

\subsection{ui.rs}
Leptos server-rendered views for the miner dashboard.

\subsubsection{SSR without a client bundle}

The same choice \texttt{apps/explorer/src/ui.rs} makes, for the same reason: a hydrating
build needs \texttt{cargo-leptos} and \texttt{wasm-bindgen} to produce a WASM bundle, and
that is a build-tooling change rather than a rewrite. Signals still drive
the render; liveness comes from the WebSocket at \texttt{/ws}, which the page
subscribes to and uses to refresh in place.

\subsubsection{What a miner is shown, and what they are not}

Every rig figure is labelled by where it came from. \textbf{Measured} hash rate
is the pool's own arithmetic over accepted share weight. \textbf{Reported} power,
temperature, and fan speed are what the rig said about itself and cannot be
checked — [\texttt{crate::telemetry}] explains why they are kept apart, and the
page says so in words rather than only in a class name.

Balances are shown in three columns and never as one total. \texttt{immature} can
still be taken away by a reorg; presenting it beside \texttt{unpaid} as a single
number would make an ordinary reversal look like theft.

\subsection{validator.rs}
The share validation pipeline.

\subsubsection{Why this is a separate stage at all}

One share is 25.4 ms of Argon2id below the DAG fork
(\texttt{src/crypto/argon\_blake.rs}). Running that on the task that owns a socket
would put it in front of every other message that connection has queued, and
at 50,000 connections would put it in front of the Tokio runtime itself.
Shares therefore cross a bounded channel into a pool of blocking workers,
and the connection task waits for a small reply rather than doing the work.

\subsubsection{Backpressure raises targets; it never drops shares}

A full queue makes [\texttt{Validator::submit}] wait, which propagates back to the
miner as a slower acknowledgement — honest, and self-limiting. Beyond that,
sustained depth triggers [\texttt{crate::vardiff::Vardiff::apply\_pressure}] across
every channel, which reduces the *arrival rate* rather than the service
rate.

Dropping submissions would be the easy shed and the wrong one: a dropped
share is work a miner already did and will not be paid for, and a pool that
does that under load is a pool miners leave under load.

\subsubsection{Concurrency is bounded by memory, not by cores}

Argon2id over 32 MiB means each concurrent verification holds 32 MiB. Eight
workers is a quarter-gigabyte of working set, which is the same arithmetic
\texttt{src/consensus/miner.rs} does for mining threads and the same cap it lands
on. \texttt{spawn\_blocking}'s default pool of 512 threads would be sixteen
gigabytes.

\subsection{vardiff.rs}
Per-channel difficulty control.

\subsubsection{This is the load governor, not a courtesy}

\texttt{docs/stratum-v2.md} §4 does the arithmetic that makes vardiff mandatory
here rather than optional: every accepted share costs a 25.4 ms Argon2id
pass over 32 MiB. At 50,000 connections, one share per second per channel is
1,270 cores of pure hashing and one share per minute is 21. Nothing else in
the pool moves that number.

\subsubsection{Steering on an EWMA of the interval, not on a share count}

The obvious controller counts shares in a fixed window and retargets on the
count. It behaves badly at exactly the moment it matters: a rig that has
just connected, or one whose target was just raised, spends the first window
with an unrepresentative count and gets retargeted on it. An exponentially
weighted mean of the *gaps between shares* has no window boundary to sit
inside, and its response to a step change is smooth.

\subsubsection{Why moves are one bit at a time}

A bit is a factor of two. Larger jumps converge faster on paper and overshoot
in practice, because the measurement lags the change: the EWMA is still full
of samples taken at the old target when the next decision is made. Single-bit
moves with a minimum number of samples between them is the standard, dull,
correct answer.

\subsubsection{Backpressure raises targets; it never drops shares}

[\texttt{Vardiff::apply\_pressure}] is what the validator pool calls when its queue
is filling. Shedding load by dropping submissions would be shedding it onto
the miners, who did the work and would not be credited for it. Raising
targets sheds the same load onto arithmetic that has not happened yet.

\section{Lattice proof of work}
\noindent\texttt{crates/lattice-pow\textbackslash{}src}

\subsection{basis.rs}
\begin{framed}\noindent\textbf{Status.} This module is a research branch. It is present in the tree and no consensus path reaches it. Nothing described here is part of the running chain.\end{framed}

The lattice a block is mined against.

A basis is [\texttt{LatticeParams}] plus the \texttt{n - 1} coefficients \texttt{x\_i}. It is
*not* an \texttt{n × n} matrix: the Hermite normal form of a Goldstein–Mayer
\texttt{q}-ary lattice is determined by its first column, so storing the matrix
would be storing \texttt{n² - n} zeros and \texttt{n - 1} ones alongside the \texttt{n - 1}
numbers that actually vary.

\subsubsection{Where the coefficients come from}

Not from here. This crate cannot hash — see the crate documentation for why
— so the expansion of a block header into coefficients lives in
\texttt{custom-l1-node}'s \texttt{crypto::lattice}, and arrives through
[\texttt{Basis::new}] as plain integers.

That the coefficients are derived from the header rather than chosen by the
miner is the property the whole scheme rests on: a miner who picked the
lattice would pick one they had already reduced. This crate cannot check
that property — it never sees a header — and the caller must not assume it
does.

\subsection{error.rs}
\begin{framed}\noindent\textbf{Status.} This module is a research branch. It is present in the tree and no consensus path reaches it. Nothing described here is part of the running chain.\end{framed}

Why a lattice solution was rejected.

These are *rule* failures, not chain errors. The crate returns its own type
for the same reason [\texttt{maya\_dex}] does: keeping \texttt{custom-l1-node} out of the
dependency graph is the point of the crate boundary, and the mapping from
"this vector is not in the lattice" to whichever \texttt{NodeError} variant a
context calls for belongs to the caller.

Every variant names one distinguishable way a solution can be wrong. They
are kept separate rather than collapsed into a single \texttt{Invalid} because a
validator that cannot say *which* rule a block broke cannot be debugged, and
because the rejection tests assert on the specific variant — a test that
only checks "some error" passes when the wrong rule fired.

[\texttt{maya\_dex}]: https://docs.rs/maya-dex

\subsection{lib.rs}
\begin{framed}\noindent\textbf{Status.} This module is a research branch. It is present in the tree and no consensus path reaches it. Nothing described here is part of the running chain.\end{framed}

Lattice proof-of-work verification for \texttt{custom-l1-node}: does this vector
lie in the block's lattice, and is it short enough?

\subsubsection{Status}

**This crate is a research branch. Nothing in the node calls it, and the
consensus rule it would serve is not wired to any activation height.** It
exists so that the verification half of a lattice proof of work can be
written, proved, and reviewed independently of the question of whether such
a proof of work is *sound*, which it is not yet known to be. See
\texttt{docs/lattice-pow.md} for what is unresolved; the short version is in
[Why this is not yet a consensus rule](\#why-this-is-not-yet-a-consensus-rule).

\subsubsection{Why this is a crate and not a module}

The same argument that produced [\texttt{maya\_ledger\_math}] and [\texttt{maya\_dex}]. Every
function here decides whether a block is admissible, and a bug in one of
them either accepts a block nobody paid for or rejects one somebody did.
That is the class of code worth proving, and the [Kani Rust Verifier]
compiles a crate *together with its whole dependency graph* — which rules
out anything reaching RocksDB's C++ or the \texttt{ark-*} stack.

The visible cost of that boundary is that nothing here hashes. The basis
coefficients arrive as plain integers; expanding them from a block header is
the node's job, in \texttt{custom-l1-node}'s \texttt{crypto::lattice}.

\subsubsection{The lattice}

A Goldstein–Mayer random \texttt{q}-ary lattice in Hermite normal form — the family
the Darmstadt SVP challenges are drawn from. For dimension \texttt{n} and modulus
\texttt{q}, the basis rows are

\begin{verbatim}
b_0 = ( q,   0, 0, …, 0 )
b_i = ( x_i, 0, …, 1, …, 0 )      1 <= i < n, the 1 in column i
\end{verbatim}

so the whole lattice is named by \texttt{q} and the \texttt{n - 1} coefficients \texttt{x\_i}, and
\texttt{det(L) = q}. That compactness is the point: a basis that would otherwise be
\texttt{n²} integers on the wire is \texttt{n - 1}.

\subsubsection{Membership is a divisibility test, not a matrix solve}

Writing \texttt{v = Σ c\_i b\_i} gives \texttt{v\_i = c\_i} for \texttt{i >= 1} and
\texttt{v\_0 = c\_0·q + Σ\_\{i>=1\} c\_i·x\_i}. Eliminating \texttt{c} leaves

\begin{verbatim}
v ∈ L   ⟺   v_0 ≡ Σ_{i=1}^{n-1} x_i · v_i   (mod q)
\end{verbatim}

This is why the miner sends \texttt{v} and not its coefficient vector: the
coefficients are recoverable, so carrying them would be redundant bytes a
miner could vary to grind. Verification is one pass of modular arithmetic in
exact integers — no floating point, no reduction, no Gram–Schmidt.

That asymmetry is the whole reason this shape of proof of work is checkable
at all: finding a short \texttt{v} is the miner's exponential problem, and
confirming one is \texttt{O(n)} additions.

\subsubsection{Determinism}

Every function is a pure function of its inputs. There is no floating point,
no hashing, no allocation-order dependence, and no time source.

This matters more here than it does elsewhere in the tree. Lattice
*reduction* — LLL, BKZ — runs on floating-point Gram–Schmidt and is not
reproducible across platforms, compilers, or library versions. If validation
re-ran reduction, the chain would fork on a rounding difference. It does
not: the validator never reduces anything. It checks a congruence and a sum
of squares, both in exact integer arithmetic, and both bounded by
[\texttt{MAX\_COORDINATE}] so that neither can overflow.

\subsubsection{Why this is not yet a consensus rule}

Recorded here rather than in a design document nobody opens, because the
next person to read this crate will otherwise assume the hard part is done.

Lattice reduction is \textbf{not progress-free}. Nakamoto consensus assumes each
mining attempt is an independent trial, so a miner who has worked for nine
minutes is no closer than one starting now. BKZ is monotone optimisation:
elapsed work strictly improves the basis. A large miner therefore finishes
reductions a small miner must abandon, reward becomes superlinear in miner
size, and cumulative work stops measuring expected trials — which is the
quantity \texttt{consensus::difficulty} compares branches by.

Nothing in this crate fixes that, and nothing in this crate pretends to.
What it provides is the half of the problem that *is* well-posed.

[\texttt{maya\_ledger\_math}]: https://docs.rs/maya-ledger-math
[\texttt{maya\_dex}]: https://docs.rs/maya-dex
[Kani Rust Verifier]: https://model-checking.github.io/kani/

\subsection{params.rs}
\begin{framed}\noindent\textbf{Status.} This module is a research branch. It is present in the tree and no consensus path reaches it. Nothing described here is part of the running chain.\end{framed}

The two numbers that name a lattice family: dimension and modulus.

\subsection{proofs.rs}
\begin{framed}\noindent\textbf{Status.} This module is a research branch. It is present in the tree and no consensus path reaches it. Nothing described here is part of the running chain.\end{framed}

Kani harnesses for the verification rules.

Run them with:

\begin{verbatim}
cargo kani -p maya-lattice-pow
\end{verbatim}

\subsubsection{Bounded versus unbounded}

Stated plainly, because the difference is the difference between a proof and
a very good test:

- [\texttt{LatticeParams::new}] is proved \textbf{unbounded}: the harness quantifies over
  every \texttt{u16} dimension and every \texttt{u32} modulus, with no loop and no unwind
  bound.
- Everything touching a vector or a coefficient list folds over a sequence,
  so those harnesses carry \texttt{\#[kani::unwind]} and are proved at
  [\texttt{PROVED\_DIMENSION}] only. They are \textbf{bounded} proofs.

The bound is small on purpose. The folds here are uniform — one
multiply-and-reduce per coordinate in [\texttt{Basis::contains}], one
square-and-add per coordinate in [\texttt{norm\_squared}] — so the induction from
step \texttt{n} to step \texttt{n + 1} carries no new case. A larger bound buys confidence
rather than coverage, at a cost that grows with the bound.

The honest reading: the *step* is proved over its whole input range, and the
*fold* is proved up to [\texttt{PROVED\_DIMENSION}] coordinates.

[\texttt{LatticeParams::new}]: crate::params::LatticeParams::new
[\texttt{Basis::contains}]: crate::basis::Basis::contains
[\texttt{norm\_squared}]: crate::verify::norm\_squared

\subsection{verify.rs}
\begin{framed}\noindent\textbf{Status.} This module is a research branch. It is present in the tree and no consensus path reaches it. Nothing described here is part of the running chain.\end{framed}

The rule: is this vector a point of the block's lattice, and is it short
enough to have cost what the target says it should have cost?

\section{Batch references and shard scheduling}
\noindent\texttt{crates/blockgraph\textbackslash{}src}

\subsection{batch.rs}
\begin{framed}\noindent\textbf{Status.} This module is a research branch. It is present in the tree and no consensus path reaches it. Nothing described here is part of the running chain.\end{framed}

A batch: the unit that is disseminated ahead of the block that commits to
it.

Transactions are held as opaque byte strings. This crate cannot decode a
Maya2C transaction — that would mean depending on \texttt{custom-l1-node} — and it
does not need to. What it enforces are the bounds that make a batch safe to
accept from a stranger: how many, how large, and none of them empty.

\subsection{error.rs}
\begin{framed}\noindent\textbf{Status.} This module is a research branch. It is present in the tree and no consensus path reaches it. Nothing described here is part of the running chain.\end{framed}

Why a batch, a reference set, or a schedule was rejected.

Rule failures, not chain errors. The crate returns its own type for the same
reason [\texttt{maya\_dex}] does: keeping \texttt{custom-l1-node} out of the dependency
graph is the point of the crate boundary.

Each variant names one distinguishable way an input can be wrong. They stay
separate rather than collapsing into \texttt{Invalid} because a validator that
cannot say which rule a block broke cannot be debugged, and because the tests
assert on the specific variant — a test checking only \texttt{is\_err()} passes when
the wrong rule fired.

[\texttt{maya\_dex}]: https://docs.rs/maya-dex

\subsection{lib.rs}
\begin{framed}\noindent\textbf{Status.} This module is a research branch. It is present in the tree and no consensus path reaches it. Nothing described here is part of the running chain.\end{framed}

Batch references and deterministic shard scheduling: which transactions a
block commits to, and what order they execute in.

\subsubsection{Status}

\textbf{Research branch. No consensus path reaches this code.} It is the
dependency-free core of the throughput work described in
\texttt{docs/blockgraph.md}; the network worker and the header change that would
use it are not written, and no activation height names it.

\subsubsection{What this takes from Narwhal, and what it does not}

Narwhal's contribution is *separating data availability from consensus*: a
block orders references to batches that were already disseminated, so block
propagation costs one digest per batch instead of every transaction byte.
That half transfers to a proof-of-work chain unchanged, and it is the half
implemented here.

The other half does not transfer. Narwhal's availability guarantee \textbf{is} the
quorum certificate — a batch is available because \texttt{2f + 1} of a known
committee signed for it — and Tusk's leader election needs a shared coin over
that same committee. This chain has no committee, no \texttt{f}, and no stake
weight; membership is open and churns. So there are no certificates here, and
availability is a *local* question: a validator that cannot fetch a
referenced batch does not build on the block that referenced it. That is
weaker than Narwhal, and it is stated rather than papered over.

\subsubsection{Why the reference set is canonically ordered}

[\texttt{BatchRefs}] requires its identifiers to be strictly ascending. Two blocks
committing to the same set of batches in different orders would be two
encodings of one block — which is a nonce the miner did not have to pay for.
The same reasoning that makes \texttt{maya-lattice-pow} reject an unreduced basis
coefficient rather than reducing it.

Strictly ascending also makes duplicate detection free: a repeat is a
non-increase, so one pass decides both rules.

\subsubsection{The throughput ceiling this implies}

Worth writing down, because the point of the exercise is a number:

\begin{verbatim}
MAX_BATCH_TRANSACTIONS  x  MAX_REFS_PER_BLOCK  =  512 x 256 = 131,072 tx/block
at TARGET_BLOCK_TIME = 15 s                    ->  ~8,738 tx/s
\end{verbatim}

That is the *structural* ceiling, not a measured rate, and it is roughly a
sixth of the 50,000/s the work was originally scoped against. The binding
constraint is not this crate: a Maya2C transaction carries an ML-DSA-65
signature (3,309 B) and an SLH-DSA-SHA2-128s signature (7,856 B), so 50,000
of them per second is 558 MB/s of signature bytes before any payload. See
\texttt{docs/blockgraph.md}.

\subsubsection{Deterministic execution instead of two-phase commit}

[\texttt{schedule}](fn@crate::schedule) partitions an already-ordered transaction list into *waves*.
Within a wave no two transactions touch a common shard, so a wave may be
executed in parallel in any order; between waves, conflicting transactions
keep their original relative order. Executing the waves in sequence produces
exactly the state the serial order would have produced.

No locks, no coordinator, no prepare phase. Two-phase commit exists to agree
an order at commit time; here the order is already agreed, so there is
nothing left for it to do — and it would be unsafe anyway, being blocking on
coordinator failure and not Byzantine-safe.

[Kani Rust Verifier]: https://model-checking.github.io/kani/

\subsection{proofs.rs}
\begin{framed}\noindent\textbf{Status.} This module is a research branch. It is present in the tree and no consensus path reaches it. Nothing described here is part of the running chain.\end{framed}

Kani harnesses for the bounds and the scheduler.

Run them with:

\begin{verbatim}
cargo kani -p maya-blockgraph
\end{verbatim}

\subsubsection{Bounded versus unbounded}

- [\texttt{shard\_of}] and [\texttt{Access::from\_mask}] are proved \textbf{unbounded}: every
  address byte, every \texttt{u64} mask.
- Everything folding over a list carries \texttt{\#[kani::unwind]} and is proved at
  [\texttt{PROVED\_LENGTH}] only. Those are \textbf{bounded} proofs.

The scheduler harnesses are the ones worth reading. They restate the two
properties the crate documentation argues informally — no conflict shares a
wave, and conflicting transactions keep their order — as post-conditions
over symbolic access sets, so a change that broke either fails here rather
than in whichever integration test happened to generate a witness.

The masks are narrowed to [\texttt{PROVED\_SHARDS}] bits. Conflict is a bitwise
\texttt{AND} and the scheduler's loop is over set bits, so a mask wider than the
number of transactions adds unreachable bits rather than new cases — and
\texttt{u64}-wide symbolic masks are what makes the solver expensive.

[\texttt{shard\_of}]: crate::shard::shard\_of
[\texttt{Access::from\_mask}]: crate::schedule::Access::from\_mask

\subsection{refs.rs}
\begin{framed}\noindent\textbf{Status.} This module is a research branch. It is present in the tree and no consensus path reaches it. Nothing described here is part of the running chain.\end{framed}

What a block commits to: a canonically ordered set of batch identifiers.

This is the change that buys the throughput. A block that carried its
transactions costs every peer the full bytes at propagation time, on the
critical path, after the block was found. A block that carries *references*
costs 32 bytes per batch, and the transactions themselves travelled earlier,
off the critical path, while miners were still searching.

\subsubsection{What is not here}

Availability. Narwhal answers "can I fetch this batch?" with a quorum
certificate from a known committee; this chain has no committee. A validator
that cannot fetch a referenced batch simply does not build on the block that
referenced it. That is a weaker guarantee and it is named in the crate
documentation rather than implied away.

\subsection{schedule.rs}
\begin{framed}\noindent\textbf{Status.} This module is a research branch. It is present in the tree and no consensus path reaches it. Nothing described here is part of the running chain.\end{framed}

Turning an agreed transaction order into waves that may run in parallel.

\subsubsection{The rule}

For transaction \texttt{i}, with \texttt{conflicts(i, j)} meaning their shard sets
intersect:

\begin{verbatim}
wave(i) = 0                                    if no j < i conflicts with i
wave(i) = 1 + max{ wave(j) : j < i, conflict } otherwise
\end{verbatim}

Two properties follow, and together they are the whole correctness argument.

\textbf{Within a wave, nothing conflicts.} If \texttt{i < j} are both in wave \texttt{w} and
share a shard, then when \texttt{j} was assigned, that shard had already been
claimed at wave \texttt{w}, so \texttt{wave(j) >= w + 1}. Contradiction.

\textbf{Conflicting transactions keep their relative order.} If \texttt{i < j} conflict
then \texttt{wave(j) > wave(i)}, directly from the rule.

So executing waves in ascending order — and transactions *within* a wave in
any order, on any number of threads — reaches the state the serial order
would have reached. That is the claim \texttt{serial\_equivalence\_tests.rs} checks
against an executable model, and it is why there is no two-phase commit
here: the order was agreed before execution began, so there is nothing left
to negotiate at commit time.

\subsubsection{Cost}

One pass, \texttt{O(n)} with a 64-bit mask operation per transaction. The
\texttt{last\_claimed} table is [\texttt{SHARD\_COUNT}] words held on the stack, so
scheduling allocates only the output.

\subsubsection{What this does not do}

It does not decide *whether* two transactions conflict — it is handed the
shard sets. A caller that computes an access set too narrowly gets a
schedule that runs them concurrently, and a wrong state root. Computing the
set is the node's job and is the part that has to be conservative: a
transaction whose accesses are not known exactly must declare more, never
fewer.

\subsection{shard.rs}
\begin{framed}\noindent\textbf{Status.} This module is a research branch. It is present in the tree and no consensus path reaches it. Nothing described here is part of the running chain.\end{framed}

Which of the 64 partitions an account's state lives in.

\subsubsection{Shards are not threads}

Worth separating, because the two get conflated. A *shard* is a partition of
state; a *thread* is an execution resource. This module defines the first.
[\texttt{crate::schedule}](fn@crate::schedule) decides what may run concurrently, and how many threads
actually serve a wave is a deployment question that changes no result — the
state root is the same on one core and on sixty-four.

\subsection{shard\_manager/handoff.rs}
\begin{framed}\noindent\textbf{Status.} This module is a research branch. It is present in the tree and no consensus path reaches it. Nothing described here is part of the running chain.\end{framed}

Moving per-shard record caches from one map to another.

The brief called this inter-shard teleportation. What moves is each leaf's
in-memory cache of the records under its range — the working set a lane
keeps hot. State itself does not move: every leaf reads one database and
commits to one root.

\subsubsection{Atomic}

[\texttt{ShardStores::apply}] checks the plan was built for the map the stores are
on before it touches a record, and from then on cannot fail: every record
is placed by the new map's own lookup, so it lands in exactly one leaf and in
key order. There is no half-moved state to observe or to roll back.

That covers the stores, not the stores *and* the manager together. The
sequence is \texttt{observe}, \texttt{apply}, \texttt{commit}, and it holds only if nothing
calls \texttt{observe} between the last two: an interleaved tick makes \texttt{commit}
refuse the plan after \texttt{apply} has already moved the caches. Both take
\texttt{\&mut}, so a caller owning the two keeps them in step by calling them in
order; one sharing them across threads must hold one lock across all three.

\subsubsection{Batched}

Records leave an old leaf in key order and new leaves are key ranges, so the
records going from one old leaf to one new leaf are contiguous. Each such run
is one [\texttt{Relocation}], which is what a lane handing its range to another
would hand over in one message.

\subsubsection{No proof}

Both sides of the handoff are this process. A proof of a relabel that the
receiver can redo would cost more than redoing it — see
[\texttt{crate::shard\_manager}].

\subsection{shard\_manager/map.rs}
\begin{framed}\noindent\textbf{Status.} This module is a research branch. It is present in the tree and no consensus path reaches it. Nothing described here is part of the running chain.\end{framed}

A tiling of the address space by binary prefixes.

A leaf is a prefix of the address's first 32 bits: \texttt{depth} leading bits
fixed, the rest free. Leaves are held in key order and must cover the space
exactly once, so a leaf is a contiguous key range — a RocksDB range scan, as
for the fixed partition — and lookup is a binary search.

Merges are \textbf{buddy} merges only: two leaves merge exactly when they are the
two halves of one parent. So every map is reachable from the root by splits
alone, and there is one map for any set of leaves rather than several that
cover the same space differently.

\subsection{shard\_manager/metrics.rs}
\begin{framed}\noindent\textbf{Status.} This module is a research branch. It is present in the tree and no consensus path reaches it. Nothing described here is part of the running chain.\end{framed}

What one scheduling tick looked like, per leaf.

Three signals, each chosen because it decides something:

- \textbf{Density} ([\texttt{TickSample::load}]): transactions touching each leaf. The
  split trigger.
- \textbf{Straddling} ([\texttt{TickSample::straddling}]): transactions touching *both*
  halves of a leaf. Splitting a leaf whose load mostly straddles turns
  intra-shard conflicts into cross-shard ones and buys no parallelism, so it
  vetoes the split.
- \textbf{Cross-shard delay} ([\texttt{TickSample::waves}], [\texttt{TickSample::cross\_shard}]):
  how many sequential waves the tick needed, and how many transactions
  spanned leaves. A wave is the unit of waiting.

Memory pressure is supplied by the caller in permille. It is this node's own
reading, which is only acceptable because the map is not consensus (see
[\texttt{crate::shard\_manager}]).

\subsection{shard\_manager/mod.rs}
\begin{framed}\noindent\textbf{Status.} This module is a research branch. It is present in the tree and no consensus path reaches it. Nothing described here is part of the running chain.\end{framed}

Elastic shards: a per-node partition that bisects under load and merges
when idle.

\subsubsection{A shard map here is an execution choice, not consensus}

The fixed partition in [\texttt{crate::shard}] is 64 prefixes of the address. This
module lets a node use a *different* tiling of the same address space — from
4 leaves up to 64, deeper where it is hot — and change it while it runs.
That is safe to do locally, and only locally, because of what
[\texttt{crate::schedule}](fn@crate::schedule) already guarantees:

- Two transactions naming a common address land in the same leaf under
  *every* tiling, so they conflict under every map and keep their order.
- Two transactions naming disjoint addresses may or may not share a leaf.
  Sharing one serialises them; not sharing lets them run together. Either
  way the waves reach the serial state.

So the state root does not depend on the map, and inputs no other node can
see — how many threads this machine has, how much memory it has left — are
legitimate triggers. They would not be if the map were consensus: a split
on one node and not another would be a chain split, the failure invariant
28 and execution directive 2 exist to prevent.

Two consequences worth stating:

- **The fixed [\texttt{crate::shard\_of}] stays.** \texttt{src/neural\_gas/features.rs}
  reads it to compute a fee feature, and a feature that moved with each
  node's load would be a consensus value that differs between nodes.
- **The guarantee rests on access sets naming every address a transaction
  touches.** That is the existing contract of [\texttt{crate::schedule}](fn@crate::schedule), and
  \texttt{core::batch::access\_for\_transaction} in the node meets it only for plain
  transfers today. With a set that is too narrow a per-node map is worse
  than a fixed one: two nodes on different maps group the same transactions
  into different waves, so the bug diverges between nodes rather than
  hiding on all of them alike. That set has to be fixed before any consensus
  path reaches the scheduler. See \texttt{docs/blockgraph.md}.

\subsubsection{Pieces}

- [\texttt{map}]: the tiling, lookup, split and buddy merge.
- [\texttt{metrics}]: what one scheduling tick looked like.
- [\texttt{policy}]: when to split and when to merge.
- [\texttt{handoff}]: moving per-shard record caches to a new map.

\subsubsection{"Teleportation" and why there is no proof}

A split or a merge relabels key ranges. All shards live in one process over
one state, so no record leaves the machine and nothing needs a proof: a
zero-knowledge proof that a relabel was done correctly would prove, at a cost
orders of magnitude above redoing it, a fact every node recomputes. The only
proof systems in this tree are also not post-quantum. What does move is each
shard's in-memory record cache, and [\texttt{handoff}] moves it atomically: a plan
built for another map is refused before anything is touched.

\subsubsection{Rounds}

Narwhal's rounds come from certificates, and this chain has no committee
(\texttt{docs/blockgraph.md}). The unit here is the \textbf{tick}: one scheduled block.

\subsection{shard\_manager/policy.rs}
\begin{framed}\noindent\textbf{Status.} This module is a research branch. It is present in the tree and no consensus path reaches it. Nothing described here is part of the running chain.\end{framed}

When to split and when to merge.

\subsubsection{The rules}

- \textbf{Split} a leaf whose load has exceeded \texttt{split\_permille} of one lane's
  capacity for \texttt{split\_streak} consecutive ticks. One tick at or below the
  line resets the streak: "100 consecutive rounds" means consecutive.
- \textbf{Merge} two buddy leaves whose *combined* load has stayed under
  \texttt{merge\_permille} of one lane for \texttt{merge\_streak} ticks. Combined, so the
  merged leaf starts below the line; and \texttt{merge\_permille < split\_permille}
  is enforced, so a merge can never be the first half of a re-split.
- A leaf created by a rebalance does neither for \texttt{cooldown\_ticks}.
- A split is vetoed while memory pressure is at or above
  \texttt{memory\_veto\_permille} (every leaf is a lane and a cache), and for a leaf
  whose load straddles its halves above \texttt{straddle\_veto\_permille} (see
  [\texttt{super::metrics}]).
- Splits and merges never share a tick, and the map stays within
  \texttt{min\_shards..=max\_shards}. Where the bound binds, the busiest leaves split
  first and the idlest pairs merge first, ties broken by key order, so a
  decision is a function of the sample and nothing else.

\subsubsection{Plan, then commit}

[\texttt{ShardManager::observe}] proposes a [\texttt{Rebalance}] without adopting it. The
caller moves its caches ([\texttt{super::handoff}]) and then calls
[\texttt{ShardManager::commit}], which refuses a plan built for another map or
another tick. A handoff that fails leaves the manager on the old map. A
caller must not \texttt{observe} between \texttt{apply} and \texttt{commit}; see
[\texttt{super::handoff}] for why.

\section{API gateway}
\noindent\texttt{crates/api-gateway\textbackslash{}src}

\subsection{allowlist.rs}
Which node methods the gateway will call, and — more importantly — which it
will not.

\subsubsection{Why this exists}

The node's JSON-RPC server has \textbf{no authentication}. That is a reasonable
design for an interface that assumes a trusted network position: it is bound
inside a pod network, and the deployment is what keeps it private.

A gateway breaks that assumption by construction. It terminates public HTTP
and forwards to the node, so every method reachable through it is a method
reachable by anyone who can reach the gateway. Inheriting the node's
"everything is allowed" posture would publish the miner interface to the
internet.

\subsubsection{Default deny}

[\texttt{is\_allowed}] returns \texttt{false} for anything not named here. A method added
to the node does not become reachable through the gateway until somebody
adds it to this list and says why — which is the review step that catches
the next \texttt{submit\_block}.

\subsubsection{What is deliberately excluded}

- **\texttt{get\_mining\_candidate}** — hands out the header a miner is searching. A
  public endpoint for it is free work for anyone building a competing
  template, and it reveals the node's mempool selection.
- **\texttt{submit\_block}** — accepts a block for validation and propagation.
  Exposing it publicly means anyone can make every node on the network do
  full block validation on demand, which is an amplification vector with no
  upside: a real miner talks to its own node.

- **\texttt{get\_headers}\textbf{, }\texttt{get\_snapshot\_manifest}** and
  **\texttt{get\_snapshot\_chunk}** — what a pruned node bootstraps from. A snapshot
  chunk is a slice of the whole committed state and a header range is up to
  2,000 headers a call, so both are cheap to ask for and expensive to
  answer. Node-to-node work, on the node's own port, not a public endpoint.
- **\texttt{get\_tip\_height}** — harmless in itself, and excluded only to keep the
  public surface to what the REST layer already covers.

Both stay reachable on the node's own RPC port, to the operator who is
already inside the trust boundary.

\subsection{error.rs}
What the gateway returns when it cannot answer.

\subsubsection{Errors do not carry upstream detail to the client}

[\texttt{GatewayError::Upstream}] holds the node's message for the log and renders
as a flat "upstream node error" to the caller. A gateway that echoed the
backend's error text would leak the node's internal addresses, RocksDB
paths, and version strings to anyone who could provoke a failure.

\subsection{graphql.rs}
GraphQL over the same node methods the REST layer serves.

\subsubsection{Why the limits are not optional}

GraphQL lets a client compose its own query. Over a chain backed by RocksDB
that is an unbounded-work surface by construction: a nested query is a
request that the server, not the client, pays for, and a public endpoint
with no ceiling is a denial-of-service primitive that ships enabled.

So [\texttt{schema}] installs a depth limit and a complexity limit at construction,
and there is no constructor without them. That is the same discipline the
rest of this tree applies to every length it reads off a wire — the ceiling
is checked before the work is done, not after it hurts.

\subsubsection{What is queryable}

Exactly what the REST layer exposes, because both go through
[\texttt{NodeClient}] and the allowlist behind it. GraphQL is a second shape over
one surface, not a second surface.

Notably absent: anything about the "DAG". \texttt{crypto::dag} in the node is the
memory-hard proof-of-work *dataset*, not a block graph, and it is not chain
state. The block-graph work in \texttt{maya-blockgraph} is a research branch that
no consensus path reaches. There is no DAG state to query, and inventing a
field that returned something plausible would be worse than its absence.

\subsection{lib.rs}
REST and GraphQL gateway in front of a Maya2C node.

\subsubsection{What this is for}

The node speaks JSON-RPC on a port that assumes a trusted network position:
no authentication, and every method equally reachable — including the miner
interface. That is a sound design for something bound to a pod network and
an unsound one to expose publicly.

This crate is the layer that makes a public interface possible: a
default-deny allowlist, size ceilings checked before allocation, GraphQL
depth and complexity limits, and errors that do not echo the backend's
internals to the caller.

\subsubsection{It talks to a node over the network}

Not through \texttt{StateDB}. The gateway is a JSON-RPC client, so it can run on a
different host from the node it serves, and RocksDB's C++ stays out of the
process that terminates untrusted HTTP. That is the same shape \texttt{explorer}
uses.

\subsubsection{The three things it refuses to do}

1. \textbf{Proxy the miner methods.} \texttt{get\_mining\_candidate} and \texttt{submit\_block}
   are not on the allowlist and are named in [\texttt{allowlist::DENIED\_METHODS}]
   so the exclusion has a test behind it.
2. \textbf{Decrypt a sealed transaction.} It holds no committee share and cannot.
   A gateway that could read the threshold-encrypted mempool would be the
   observer that mempool exists to remove.
3. \textbf{Expose an unbounded query.} Every GraphQL schema this crate builds
   carries a depth and a complexity limit; there is no constructor without
   them.

\subsection{node.rs}
The gateway's view of a node: a JSON-RPC client behind a trait.

The trait exists so the REST and GraphQL layers can be tested against a
recorded node without a running chain. That is not a convenience — an
end-to-end test that needs RocksDB, a genesis file, and a mined block in
order to check that a 404 is a 404 is a test nobody runs.

\subsection{rest.rs}
REST routes.

One route per allowlisted node method, plus sealed submission and health.
Nothing here reaches the node directly — every call goes through
[\texttt{NodeClient}], which enforces the allowlist at the single point all
outbound traffic passes.

\subsection{sealed.rs}
Submitting a threshold-encrypted transaction.

\subsubsection{What the gateway does and does not do}

It validates *shape* — hex, length, non-empty — and forwards. It does
\textbf{not} decrypt, and it holds no share of the committee key. That is the
point of a threshold-encrypted mempool: the party that orders transactions
cannot read them, and a gateway that could decrypt would reintroduce
exactly the observer the scheme exists to remove.

So the ciphertext is opaque here, and every check below is a check a proxy
can make without understanding the payload.

\subsubsection{Why there are size limits at all}

The gateway terminates public HTTP. Without a ceiling, a submission
endpoint is a way to make the process allocate whatever the sender feels
like. [\texttt{MAX\_SEALED\_PAYLOAD\_BYTES}] is checked \textbf{before} the hex is decoded,
so a hostile length never becomes a hostile allocation — the discipline
\texttt{core/codec.rs} applies on the wire, applied at the edge.

\section{Language bindings}
\noindent\texttt{sdks/sdk-ffi\textbackslash{}src}

\subsection{bin/uniffi-bindgen.rs}
The bindings generator, as uniffi requires it: a binary in the crate that
defines the interface, so the generator sees the same proc-macro metadata
the library exports.

\subsection{lib.rs}
Python, Kotlin, and Swift bindings over Maya2C's post-quantum signing.

\subsubsection{The security rule this crate is built around}

\textbf{A secret key never crosses the FFI boundary.}

\texttt{HybridSigningKey} in the node is \texttt{Zeroizing}: its bytes are wiped when it
drops. That guarantee ends at the boundary. Hand the secret to Python and it
becomes a \texttt{bytes} object — immutable, freely copied by the interpreter,
resident until garbage collection, and impossible to wipe. Kotlin and Swift
are no better. A wallet SDK that exported key bytes would be quietly
weaker than the Rust it wraps, in the one respect that matters most.

So [\texttt{SigningKey}] is an \textbf{opaque handle}. It can sign, and it can hand out
its public key and address. There is deliberately no \texttt{secret\_bytes()}, and
adding one would undo the reason this crate is shaped the way it is.

The one place secret material may enter is [\texttt{SigningKey::from\_seed}], and
that is the caller's own seed arriving — the SDK is not exporting anything
the caller did not already hold.

\subsubsection{What is verified, and what is only generated}

uniffi emits Python, Kotlin, and Swift from one interface. Emitting is not
the same as working:

\begin{verbatim}
| Target | Generated | Compiled and tested |
|---|---|---|
| Python | yes | **yes** — `sdks/sdk-ffi/tests/python/` |
| Kotlin | yes | no — no `kotlinc` on the build host |
| Swift | yes | no — no Swift toolchain on the build host |
\end{verbatim}

That table is in the README too. Generated-and-never-compiled bindings are
not a working SDK, and reporting all three as done would be a green tick
over something nobody ran.

\subsubsection{Cost}

A hybrid signature is 11,165 bytes — ML-DSA-65 at 3,309 and
SLH-DSA-SHA2-128s at 7,856 — and both must verify. Signing is dominated by
the hash-based half. Callers building interactive software should sign off
the UI thread; see \texttt{sdks/sdk-ffi/README.md} for measured figures.

\section{Browser bindings}
\noindent\texttt{sdks/sdk-wasm\textbackslash{}src}

\subsection{lib.rs}
Browser-side signing and verification for Maya2C, via wasm-bindgen.

\subsubsection{Why this crate does not use the node's `hybrid` module}

\texttt{custom\_l1\_node::crypto::hybrid} is the authoritative implementation, and
reaching for it would be the obvious right answer. It is not available:
\texttt{custom-l1-node} links RocksDB, which is C++, and C++ does not target
\texttt{wasm32-unknown-unknown}. There is no feature flag that removes it, because
the node *is* the database.

So the construction is composed here from the same two primitives the node
composes it from — \texttt{fips204} for ML-DSA-65 and \texttt{maya-crypto-pq} for
SLH-DSA-SHA2-128s — both of which are pure Rust and both of which build for
wasm.

\textbf{That is duplication of consensus-critical encoding.} Duplication is
dangerous when it is undetected, so it is detected: \texttt{tests/parity\_tests.rs}
signs with the node's implementation and verifies with this one, and the
reverse, and compares addresses and encodings byte for byte. Drift fails a
test rather than producing a transaction the chain silently rejects.

The structural fix is to extract \texttt{hybrid} into its own crate that both the
node and this SDK depend on. That is a refactor of a consensus path and is
deliberately not bundled into an SDK change; see \texttt{sdks/sdk-js/README.md}.

\subsubsection{Cost in a browser}

A hybrid signature is 11,165 bytes and the hash-based half dominates
signing. In wasm there is no AVX2 and no SHA extensions, so expect
materially worse than the native figures. \texttt{sdk-js} measures it and reports;
nothing here asserts a bound. Sign off the main thread.

\end{document}
